% SOURCE docs/release/2026-10-29/source/docs/writing-repair/2026-10-29/src/47033306ecbad01f.tex CHAR 0
\chapter{Panel and time-series comparisons for a banking rollout}
\label{exec17:s6}

A regional rollout of a deposit-servicing program creates several plausible
comparisons. An early region can be compared with a later region, with a
weighted combination of unaffected regions, or with its own forecast under the
old operating rule. Each comparison uses a different source of counterfactual
information. They also answer different questions when treatment effects
vary over the life of a relationship, the donor regions face different funding
shocks, or the bank changes the rule that responds to those shocks.

The earlier panel and time-series chapters establish the identification and
estimation foundations. The task here is to connect the resulting estimands
to executable comparisons without discarding heterogeneity, extrapolation risk
or the calendar on which uncertainty arises. A regression's lead coefficient
may contain post-treatment effects. Exact donor balance may increase mean
squared error. Positive persistence may reduce the variance of an adjacent
difference while leaving a much larger uncertainty for a comparison spanning
several periods. An impulse response under the existing feedback rule need
not describe the response after that rule changes.

We develop these distinctions using exact stipulated laws. The amounts are
illustrative dollars per declared unit of exposure; none is an estimate from
a bank. The constructions supply reproducible mathematical references for a
future public-data or generated-data prototype. Literature claims retain their
specific original versions, while the finite examples and the extensions of
their algebra are project derivations.~\citep{exec17_s6_math}

The first part develops the comparison weights, sensitivity budgets and
feedback values that a software implementation must preserve. The later
studies examine competing interventions, changing populations and concrete
R, Python, Stata, Java and C++ routes. Their purpose is to locate where a
sound mathematical comparison can change meaning during estimation,
aggregation or reporting, and to distinguish a source-backed proposal
from a qualified native implementation.

\section{An event-study lead can be made from treatment effects}
\label{exec17:s6-event}

Let $G$ denote the first treatment period, with $G=\infty$ for a never-treated
group. Under an absorbing rollout, the observed outcome in cohort $g$ at time
$t$ has mean
\begin{equation}
 y_{gt}=\mu_{gt}^{0}+1\{t\geq g\}\tau_{gt},
 \qquad \tau_{gt}=E\{Y_t(g)-Y_t(\infty)\mid G=g\}.
 \label{exec17:s6-cohort-law}
\end{equation}
The superscript zero denotes the untreated mean. No anticipation says that
the treated and untreated potential outcomes agree before $g$. It places no
restriction requiring the early and late cohorts to have equal effects after
their own adoption. Sun and Abraham show that pooling relative-time
indicators with unit and time effects can mix treatment effects from other
relative times into a reported coefficient, including a lead. The inspected
original is their September 2020 arXiv version, associated with the 2021
publication; its technical text is not represented as the final journal
version.~\citep{sun2021}

Consider an equally weighted, balanced population table with three groups,
$g=2,3,\infty$, and five periods, $t=0,1,2,3,4$. All untreated means are zero.
All treatment effects are zero except $\tau_{2,2}=20$. Thus the early region
has a one-period implementation gain of USD 20, and neither region has any
anticipatory effect. These premises deliberately eliminate a true pretrend.
They isolate a projection problem that would otherwise be confused with a
failure of the rollout's causal assumptions.

For each of the 15 cells, put an intercept, two group indicators and four
time indicators in a nuisance matrix $Z$. Its seven columns are linearly
independent. Put the five event-time indicators
\begin{equation}
 D_{(g,t),k}=1\{g\ne\infty,\ t-g=k\},
 \qquad k\in(-3,-2,0,1,2)
 \label{exec17:s6-design}
\end{equation}
in a matrix $D$. Every event indicator is zero for the never-treated group.
The omitted relative period is $-1$. The extreme lead $-3$ and lag $2$
are included explicitly, rather than silently pooled or dropped. The complete
$15$ by $12$ design $[Z,D]$ has exact rank 12. Consequently the following
coefficients are uniquely defined; no generalized inverse or accidental
collinearity supplies the result.

Residualize against $Z$ by $M_Z=I-Z(Z'Z)^{-1}Z'$. The event coefficients
in the population least-squares projection are
\begin{equation}
 b=(D'M_ZD)^{-1}D'M_Zy.
 \label{exec17:s6-fwl}
\end{equation}
For this example $y$ consists entirely of the five possible post-treatment
cell effects, ordered as
$(\tau_{2,2},\tau_{2,3},\tau_{2,4},\tau_{3,3},\tau_{3,4})'$.
The corresponding columns of the projection operator are
\begin{equation}
 W=
 \begin{pmatrix}
 13/38&4/19&0&-13/38&-4/19\\
 25/171&11/171&0&-25/171&-11/171\\
 11/18&-1/9&0&7/18&1/9\\
 23/114&29/57&0&-23/114&28/57\\
 -8/171&17/171&1&8/171&-17/171
 \end{pmatrix}.
 \label{exec17:s6-projection}
\end{equation}
Rows follow the five event times in \eqref{exec17:s6-design}. These are
projection weights for this declared table, not universal software weights.
For example, the first row subtracts the late cohort's impact and first-lag
effects from positive weights on the early cohort's corresponding effects.
It is not an average of pre-treatment outcomes or pre-treatment causal
effects. The first two row sums are zero, but their individual entries are
nonzero. A zero row sum cannot protect a lead against heterogeneous effects.

Multiplication by $(20,0,0,0,0)'$ yields
\begin{equation}
 (b_{-3},b_{-2},b_0,b_1,b_2)'
 =\left(\frac{130}{19},\frac{500}{171},
          \frac{110}{9},\frac{230}{57},-\frac{160}{171}\right)'.
 \label{exec17:s6-false-leads}
\end{equation}
Both leads are positive although every pre-treatment causal effect is zero.
Even the lag-two coefficient is negative although the only nonzero effect
in the entire table is positive. The example therefore concerns the meaning
of the coefficients before any standard error is computed.

There are two useful checks on this calculation. Direct solution of the full
normal equations gives a residual orthogonal to all 12 columns. Alternatively,
double demean every event indicator and the outcome: subtract its group
mean and time mean and add its overall mean. Regressing the demeaned outcome
on the five demeaned indicators gives exactly
\eqref{exec17:s6-false-leads}. The two calculations use different
representations of the nuisance adjustment. Both retain rational arithmetic,
so a small numerical tolerance cannot hide the false leads.~\citep{exec17_s6_math}

The weights explain when cancellation would occur. If the two cohorts had
the same impact effect and the same first-lag effect, the positive and negative
entries in the first two rows would cancel pairwise. The witness violates
that homogeneity while satisfying no anticipation and an identical untreated
law. It does not show that an observed bank pretrend is always spurious.
An actual lead can reflect genuine anticipation, a different untreated trend,
changing composition, or regression contamination. Distinguishing these
possibilities requires a target-preserving estimator and substantive knowledge
of assignment and information.

For a regional program, an apparent gain two periods before adoption should
therefore trigger examination of the comparison weights and rollout clock.
It should not automatically certify early customer response or disqualify the
program. Adding more customers to the same misinterpreted projection can
make its standard error smaller while preserving the wrong population
quantity. In particular, a precise lead test cannot turn
\eqref{exec17:s6-fwl} into an estimate of a different causal target.

The operational alternative is to retain cohort-by-calendar-time contrasts,
their admissible untreated comparisons and their joint uncertainty, then
aggregate to the declared exposure target. The earlier cohort identification
argument remains necessary: choosing a heterogeneous-effect estimator does
not establish parallel trends, support or absence of spillovers.
The inspected \texttt{eventstudyinteract} manual describes an interacted
regression followed by cohort-share weighting, and warns about its control
cohort and balanced-panel validity conditions. Its returned interacted
coefficients and its weighted estimates are distinct objects.
Those restrictions should govern a proposed wrapper, including the choice
of saved covariance; the presence of a familiar event-study plot cannot
serve as semantic admission.~\citep{v12_sw_eventstudyinteract}

\section{The horizon determines the trend-sensitivity budget}
\label{exec17:s6-curvature}

Correcting the event-study projection does not settle the counterfactual
trend. We now take as input an admissible treated-minus-comparison mean
series for a fixed population. This can be one cohort or a fixed aggregation
with a common, explicitly defined gap process. It cannot be a sequence whose
cohort composition changes and is then interpreted as the trend of one
unchanging population without further justification.

Index the last untreated period by $-1$, and let $d_t$ be the untreated
treated-minus-comparison gap. Define the last observed untreated slope
$v=d_{-1}-d_{-2}$. For future periods $j=0,\ldots,J-1$, impose the
specific curvature restriction
\begin{equation}
 u_j=d_j-2d_{j-1}+d_{j-2},\qquad |u_j|\leq M.
 \label{exec17:s6-curvature-law}
\end{equation}
The constant $M$ is measured in outcome units per squared calendar period.
It is a maintained sensitivity choice, not the largest observed residual
relabelled as truth. Changing months to quarters changes the restriction.
Rambachan and Roth motivate smoothness and relative-magnitude restrictions
as distinct choices and develop inference that accounts for noisy
pre-treatment information. The construction below specializes a curvature
restriction to one explicit horizon target; it is not an implementation of
their optimized inference procedures.~\citep{rambachan2023}

The bias in a baseline-$-1$ contrast at horizon $j$ is
$b_j=d_j-d_{-1}$. Taking first differences in
\eqref{exec17:s6-curvature-law} and summing twice gives
\begin{equation}
 b_j=(j+1)v+\sum_{k=0}^{j}(j-k+1)u_k.
 \label{exec17:s6-triangular}
\end{equation}
Indeed, the next slope is $v+u_0$, the following slope is
$v+u_0+u_1$, and so on. Summing these slopes counts the curvature
$u_k$ exactly $j-k+1$ times. A distant effect is exposed to an early
change of slope repeatedly. This triangular propagation is the reason a
long rollout evaluation can be much more sensitive than its first-period
estimate, even with the same numerical $M$.

Let $\beta_j$ be the observed post-period contrast and $\tau_j$ the
corresponding causal effect, so $\beta_j=\tau_j+b_j$. A fixed target
$\theta=\sum_{j=0}^{J-1}a_j\tau_j$ may be an average effect, a discounted
total or a signed change between horizons. Define
\begin{equation}
 L=\sum_{j=0}^{J-1}a_j(j+1),\qquad
 c_k=\sum_{j=k}^{J-1}a_j(j-k+1).
 \label{exec17:s6-target-kernel}
\end{equation}
Then $\theta=a'\beta-Lv-\sum_k c_ku_k$. Under only the box
restriction in \eqref{exec17:s6-curvature-law}, its sharp identified set is
\begin{equation}
 \Theta(M)=
 \left[a'\beta-Lv-M\sum_k|c_k|,
       \ a'\beta-Lv+M\sum_k|c_k|\right].
 \label{exec17:s6-sharp-curvature}
\end{equation}
To attain the upper endpoint choose
$u_k=-M\operatorname{sign}(c_k)$, with any permitted value when
$c_k=0$. Reverse the signs for the lower endpoint. Starting from the
two fixed pre-period gaps and recursively generating $d_j$ produces a
permitted path for each endpoint. Every intermediate target value is
attainable by convex combinations of the curvature vectors. Sharpness here
uses unrestricted real outcomes and no additional shape constraints; bounded
or otherwise restricted outcomes can narrow the set.

For nonnegative $a_j$, the radius simplifies to
$M\sum_j a_j(j+1)(j+2)/2$. Signed weights require the absolute values
\emph{after} summing the horizon contributions to each curvature. Bounding
each horizon separately and adding absolute target weights is safe but can
be unnecessarily wide. For example, $a=(1,-2,0)'$ has
$c=(-3,-2,0)'$ and radius $5M$, whereas summing the separate absolute
horizon bounds gives $7M$. For $a=(1,-1,1)'$, the kernel is
$(2,1,1)'$ and the radius is $4M$. A package translation must therefore
preserve the target vector, rather than assemble a monetary result from
unrelated marginal sensitivity intervals.

Take pre-period gaps $(-30,-20,-10)$ at times $(-3,-2,-1)$ and
post-period baseline contrasts $(30,35,40)$. These are mean outcome
amounts in dollars per standardized exposure-period. With equal weights
$a=(1,1,1)'/3$, $L=2$, $c=(2,1,1/3)'$, and
\begin{equation}
 a'\beta=35,\qquad Lv=20,\qquad
 \Theta(M)=[15-(10/3)M,\ 15+(10/3)M].
 \label{exec17:s6-money-sensitivity}
\end{equation}
At $M=3$ the effect set is $[5,25]$. The endpoint paths can be
constructed by setting all three future curvatures to $3$ or $-3$.
Exhaustive evaluation of the eight vertices agrees with the triangular
formula. This is a population identification calculation, so its endpoints
do not carry a 95\% coverage label.~\citep{exec17_s6_math}

Suppose the recurring incremental operating cost is USD 10 in the same
exposure-period units. Subtracting this fixed cost gives the net-effect set
$[5-(10/3)M,5+(10/3)M]$. A strictly positive lower net value requires
$M<1.5$; at $M=1.5$ it is exactly zero. At $M=3$ the net range is
$[-5,15]$, which leaves the sign unresolved even with the population means
known. This threshold explains what substantive trend restriction an
investment recommendation would require. It does not establish that the
restriction is credible or that the illustrative cost is an actual price.

An average exposure-period benefit must also be distinguished from a
three-period total. With the same three equal exposure counts, the total
has target weights $(1,1,1)'$, center 45 and curvature radius $10M$.
Three periods of the same recurring cost total 30, leaving net center 15.
The positivity threshold remains $M<1.5$, because both benefit and cost
were scaled consistently. Subtracting a single cost of 10 from that total
would answer a different economic question. Unequal exposure counts,
discounting or an upfront fixed cost alter the target weights or cost
functional and require recalculation of the kernel.

\subsection{Retaining the uncertainty in the extrapolated slope}

The preceding center contains an estimated pre-period slope in a real study.
Normalize the gap at $-1$ to zero and let
$\beta_{-2}=d_{-2}-d_{-1}=-v$. Assemble the estimate vector
$\widehat\eta=(\widehat\beta_{-2},\widehat\beta_0,\ldots,
\widehat\beta_{J-1})'$ with its full covariance matrix $\Sigma$.
For fixed target weights, define $h=(L,a')'$. The trend-adjusted center
and its sampling variance are
\begin{equation}
 \widehat m=h'\widehat\eta,\qquad s_m^2=h'\Sigma h.
 \label{exec17:s6-noisy-slope}
\end{equation}
The plus sign on the pre-period coefficient is necessary because the
subtracted slope is $v=-\beta_{-2}$. Even if the post-period contrasts
were estimated independently of one another, the common estimated baseline
would generally induce covariance with the pre-period contrast. Removing the
normalized zero coefficient does not remove that sampling dependence.

For a finite normal design with
$\widehat\eta-\eta\sim N(0,\Sigma)$ and known $\Sigma$, a conservative
$1-\alpha$ interval for the fixed target under the fixed restriction $M$ is
\begin{equation}
 \left[\widehat m-z_{1-\alpha/2}s_m-M\|c\|_1,
       \widehat m+z_{1-\alpha/2}s_m+M\|c\|_1\right].
 \label{exec17:s6-normal-interval}
\end{equation}
On the event $|\widehat m-m|\leq z_{1-\alpha/2}s_m$, every target
permitted by the curvature restriction lies in this interval. That event
has probability $1-\alpha$ under the stipulated Gaussian design. This
proves the stated conservative coverage without identifying the unknown
curvature from a pretest. With an estimated covariance and an asymptotic
normal approximation, a corresponding conclusion needs the approximation
and covariance consistency; finite exact coverage does not follow merely
from inserting a standard error.

For the three-period average, stipulate the diagonal covariance
$\Sigma=\operatorname{diag}(1,4,4,4)$ in squared dollars. This is a
separate hypothetical sampling law, not a covariance inferred from the
15-cell event-study table. Then
\begin{equation}
 h=(2,1/3,1/3,1/3)',\qquad s_m^2=16/3.
 \label{exec17:s6-normal-design}
\end{equation}
Ignoring the noisy slope would retain only variance $4/3$, one quarter
of the correct variance under this particular law. At a 95\% normal
critical value of about 1.96, the sampling radius is about USD 4.53.
For an observed center of 15 and recurring cost 10, a positive lower
net endpoint would require
$5-1.96\sqrt{16/3}-(10/3)M>0$, approximately $M<0.142$.
Population robustness up to $M=1.5$ thus need not imply useful statistical
robustness with this sample precision.

The simple interval is deliberately conservative and uses only the last
pre-period slope. Additional pre-period information, joint restrictions and
better optimized inference can improve efficiency. The comparison with
HonestDiD should retain this distinction. Its documented smoothness and
relative-magnitude procedures use different restriction families; the original
paper discusses why fixed-length, conditional and hybrid approaches have
different domains. A result from one cannot be labelled as another because
their output endpoints happen to resemble each other.~\citep{rambachan2023,sw_honestdid}

There is a further selection issue. The proof fixes $M$, the target $a$
and the method before examining favorable post-period estimates. Searching
over horizons or restrictions and reporting the first positive endpoint is
a different random procedure. It requires a joint selection or simultaneous
inference argument. A sensitivity curve can be reported as a collection of
conditional conclusions without claiming that the most favorable point on
the curve has its nominal coverage after selection.

\section{Donor balance is a bias--noise decision}
\label{exec17:s6-donor}

An unaffected donor pool can supply a different counterfactual. Perfect
pre-period fit is attractive, but the future value of extrapolating beyond
that pool depends on noise and contamination as well as balance. To isolate
the tradeoff, take a single post-period untreated outcome law
\begin{equation}
 Y_T(0)=2f+\varepsilon_T,\qquad
 Y_0(0)=\varepsilon_0,\qquad
 Y_1(0)=f+\varepsilon_1.
 \label{exec17:s6-donor-law}
\end{equation}
Here the treated region has known loading 2, the two donors have known
loadings 0 and 1, and the common factor satisfies $|f|\leq B$.
The centered errors are independent with donor variances $\sigma^2$
and treated variance $\sigma_T^2$. The loadings and bound are conditional
premises of the example, not quantities learned exactly from a short panel.

Use donor weights $(1-a,a)$, which sum to one. Their sum removes a common
additive level, if present, but their weighted loading is $a$. The untreated
counterfactual residual is
\begin{equation}
 R(a)=Y_T(0)-(1-a)Y_0(0)-aY_1(0)
      =(2-a)f+\varepsilon_T-(1-a)\varepsilon_0-a\varepsilon_1.
 \label{exec17:s6-donor-residual}
\end{equation}
The same residual is the error in a treatment-effect estimate obtained by
subtracting this donor counterfactual from the treated observation. Conditional
on $f$, its expectation is $(2-a)f$, and its variance is
$\sigma_T^2+\sigma^2\{(1-a)^2+a^2\}$. Taking the worst mean squared
error over the declared factor interval gives
\begin{equation}
 {\cal R}(a)=\sigma_T^2+\sigma^2\{(1-a)^2+a^2\}
                         +B^2(2-a)^2.
 \label{exec17:s6-donor-risk}
\end{equation}
The derivative is
$2\sigma^2(2a-1)+2B^2(a-2)$, so the unique unrestricted minimizer,
when $2\sigma^2+B^2>0$, is
\begin{equation}
 a^*=\frac{\sigma^2+2B^2}{2\sigma^2+B^2}.
 \label{exec17:s6-donor-optimum}
\end{equation}
It lies between $1/2$ and 2. As $B$ approaches zero with positive
noise, the two donors are averaged to reduce variance. As the factor bound
dominates the noise, the weight approaches 2 and exactly matches the
treated loading, at the cost of a negative first-donor weight.

With $B=2$ and $\sigma^2=\sigma_T^2=1$, the complete risk is
\begin{equation}
 {\cal R}(a)=6a^2-18a+18
            =\frac92+6(a-3/2)^2.
 \label{exec17:s6-donor-square}
\end{equation}
Nonnegative weights restrict $a$ to $[0,1]$; the best value on that
interval is $a=1$ with risk 6. Exact loading balance uses $a=2$
and also has risk 6. The intermediate weights $(-1/2,3/2)$ have
risk $9/2$. Their squared bias at the worst factor endpoint is 1,
while their noise variance is $7/2$. The exact fit eliminates this bias
but raises the noise variance to 6. In this law, exact fit spends too much
variance to remove the remaining imbalance.

This calculation is a conditional minimax comparison over a known factor
bound and fixed weights. It is not a general risk guarantee for ridge
augmentation, and $a^*$ is not claimed to be the output of a particular
package. Ben-Michael, Feller and Rothstein develop augmented synthetic control,
including a ridge representation that adjusts balance and can produce
negative weights. Their inspected July 2020 manuscript explains how a penalty
controls departures from the original synthetic-control weights. That
mechanism motivates comparison with the tradeoff above; it does not identify
the bank's $B$ or noise law.~\citep{benmichael2021}

Diagnostics should therefore report more than a pre-fit score. The
estimated loading or outcome imbalance, the norm of the weights, and the
validation distribution under which a penalty is selected all matter.
Selecting a penalty on post-treatment outcomes can absorb the effect
being estimated. Selecting it on held-out pre-periods avoids that specific
leakage but still requires an argument that the pre-period prediction task
is informative about the relevant untreated post-period task. A factor that
becomes economically important only after rollout can defeat that transfer.

\subsection{Signed weights also amplify donor contamination}

Suppose the donor observations instead contain unknown mean shifts $c_0,c_1$
with $|c_i|\leq C$. Such shifts may represent spillovers from the program
or an unrecorded intervention in the comparison region. Keeping
\eqref{exec17:s6-donor-law} otherwise fixed, the worst absolute mean error is
\begin{equation}
 B|2-a|+C\{|1-a|+|a|\}.
 \label{exec17:s6-contamination-radius}
\end{equation}
The bound is sharp over the rectangular uncertainty set: choose the sign of
$f$ and each donor shift to align the terms in the residual. If scientific
knowledge links the shifts or restricts their signs, that uncertainty set
should be changed instead of asserting this rectangle is inevitable.

The weight sums are all one, but their absolute sums are 1, 2 and 3
at $a=1,3/2,2$. At a contamination bound $C=3$, the donor component
alone can therefore contribute USD 3, 6 or 9 of bias. Exact loading fit
does not protect against this distinct source of error. In a regional
deposit rollout, transfers of customers or marketing exposure across region
boundaries can be more consequential for a signed counterfactual than for
a nonnegative one, even though both use the same named donors.

If the desired criterion is worst mean squared error with both sources of
bounded bias, the correct objective is
\begin{equation}
 {\cal R}_C(a)=\sigma_T^2+\sigma^2\{(1-a)^2+a^2\}
       +\left[B|2-a|+C\{|1-a|+|a|\}\right]^2.
 \label{exec17:s6-contaminated-risk}
\end{equation}
The two bias budgets add before squaring because an adverse law can align
them. Adding their squared radii would discard the possible cross term.
This distinction is consequential even when every error variance is known.

For $B=2,C=1$ and both variances equal to one, the bias radius is
$5-4a$ for $a\leq0$, $5-2a$ for $0\leq a\leq1$, 3 for
$1\leq a\leq2$, and $4a-5$ for $a\geq2$. The corresponding risks are
\begin{equation}
 {\cal R}_1(a)=
 \begin{cases}
 18a^2-42a+27,&a\leq0,\\
 6a^2-22a+27,&0\leq a\leq1,\\
 2a^2-2a+11,&1\leq a\leq2,\\
 18a^2-42a+27,&a\geq2.
 \end{cases}
 \label{exec17:s6-contaminated-pieces}
\end{equation}
The risk decreases up to the kink at $a=1$ and increases after it.
More explicitly, the stationary points of the four pieces lie outside
their respective intervals, and checking the three finite kinks gives
risks $27,11,15$. The global minimizer is now the nonnegative choice
$a=1$. The previously optimal $a=3/2$ has risk $25/2$, while
exact fit has risk 15.

This reversal is not a demonstration that signed weights should always be
rejected. It demonstrates that a conclusion about regularization depends
on the uncertainty class. A comparison claiming robustness to spillovers
must price that class explicitly. The numerical reference also checks the
three candidate risks by enumerating every endpoint of the factor and
contamination rectangle and every combination of independent centered
unit-variance two-point errors. That direct expectation calculation agrees
with \eqref{exec17:s6-contaminated-risk}.~\citep{exec17_s6_math}

\section{Synthetic DiD needs both axes of balance}
\label{exec17:s6-sdid}

The preceding donor example isolates the unit axis. A block rollout also
permits weighting the pre-treatment calendar. To see what this adds, suppose
the untreated mean satisfies
\begin{equation}
 E\{Y_{it}(0)\}=\alpha_i+\gamma_t+\ell_i'f_t,
 \label{exec17:s6-factor-law}
\end{equation}
where $\ell_i$ and $f_t$ are vectors of the same dimension. This is a
conditional factor representation, with additive unit and calendar effects
written separately. We do not infer its rank or predictive validity from a
small pre-fit residual. Let $u_i$ be normalized nonnegative weights on
treated units, $w_i$ normalized donor weights, $v_t$ normalized
nonnegative weights on post periods, and $\lambda_t$ normalized
pre-period weights. Each of the four vectors sums to one over its own
support. Donor and pre-period weights may be signed in the algebra below;
allowing such weights is a comparison exercise, not a description of the
default constraints in an upstream estimator.

Form the fixed-weight double difference
\begin{equation}
 \widehat\tau_{w,\lambda}
 =\sum_{i\in T}u_i\left(\sum_{t\in P}v_tY_{it}
                         -\sum_{t\in B}\lambda_tY_{it}\right)
 -\sum_{i\in C}w_i\left(\sum_{t\in P}v_tY_{it}
                         -\sum_{t\in B}\lambda_tY_{it}\right).
 \label{exec17:s6-double-difference}
\end{equation}
Here $T,C$ are the treated and control sets, while $P,B$ are post and
pre calendars. In a block design, only cells in $T$ by $P$ receive the
treatment. The target causal effect is the treated-post weighted mean
$\tau_{u,v}=\sum_{i\in T,t\in P}u_iv_t\tau_{it}$. This explicit target
allows heterogeneous effects; it should not be silently exchanged for a
simple unweighted ATT when the exposure weights differ.

Normalization cancels the unit effects because the post-minus-pre weights
sum to zero for each unit. It cancels the time effects because the
treated-minus-donor weights sum to zero at each time. The remaining
factor contribution is exactly
\begin{align}
 E(\widehat\tau_{w,\lambda})-\tau_{u,v}
 &=\Delta_\ell'\Delta_f,\label{exec17:s6-imbalance-product}\\
 \Delta_\ell&=\sum_{i\in T}u_i\ell_i-\sum_{i\in C}w_i\ell_i,\nonumber\\
 \Delta_f&=\sum_{t\in P}v_tf_t-\sum_{t\in B}\lambda_tf_t.\nonumber
\end{align}
The proof expands the four rectangles of the factor mean. Every term is
bilinear, so they collect into the product of the two imbalances.
This gives a precise reason to combine unit and time weighting. Eliminating
either imbalance eliminates factor bias under this law. Reducing both can
make their product small without making either exactly zero.

A Euclidean bound is
$|\Delta_\ell'\Delta_f|\leq\|\Delta_\ell\|_2\|\Delta_f\|_2$.
But the separate Euclidean norms depend on how factors are scaled. Replacing
$\ell_i$ by $A\ell_i$ and $f_t$ by $A^{-T}f_t$ for an invertible $A$
leaves every outcome mean and the product unchanged. It need not preserve
the two Euclidean norms. Consequently numerical certificates for separate
loading and factor bounds must declare their normalization; the observed
outcome fit alone does not provide a unique latent scale.

For fixed weights and independent cell errors with common variance
$\sigma^2$, the noise variance of the contrast is
\begin{equation}
 \operatorname{Var}(\widehat\tau_{w,\lambda})
 =\sigma^2(\|u\|_2^2+\|w\|_2^2)
             (\|v\|_2^2+\|\lambda\|_2^2).
 \label{exec17:s6-sdid-noise}
\end{equation}
Each cell's coefficient is the product of its signed unit and time
coefficient. Squaring and summing these products separates into the two
squared norms. This proof requires independent cells and weights fixed
relative to the errors. Serial correlation, shared calendar shocks or
weights fitted to those errors change the variance problem.

For an exact comparison, take one treated unit with loading 2, two donors
with loadings 0 and 1, pre-period factors 0 and 1, and one post-period
factor $3/2$. Begin with donor weights $(0,1)$ and pre-period weights
$(1/4,3/4)$. The loading imbalance is 1 and the factor imbalance is
$3/4$, giving bias $3/4$ and noise variance $(13/4)\sigma^2$.
Changing only the time weights to $(-1/2,3/2)$ matches the post factor
exactly and eliminates the factor bias. It raises the noise variance to
$7\sigma^2$. Changing only the donor weights to $(-1,2)$ also
eliminates factor bias, but gives variance $(39/4)\sigma^2$.
All three calculations can be recovered by multiplying the actual
$3$ by $3$ factor-mean table by its signed unit and time coefficients.

Under these independent-cell premises, their mean squared errors are
\begin{equation}
 {\cal R}_{+}=9/16+(13/4)\sigma^2,\qquad
 {\cal R}_{\mathrm{time}}=7\sigma^2,\qquad
 {\cal R}_{\mathrm{unit}}=(39/4)\sigma^2.
 \label{exec17:s6-sdid-risk}
\end{equation}
The signed time extrapolation improves on the nonnegative comparison
only when $\sigma^2<3/20$. The signed unit comparison has a still
higher variance than the signed time comparison in this particular table.
At $\sigma^2=1$, leaving the small factor bias costs less in mean
squared error than either exact fit. At $\sigma^2=1/10$, signed time
matching has risk $7/10$, below the nonnegative risk $71/80$.
The comparison depends on the post factor, which was stipulated as known
here; an analyst cannot select the favorable matching rule using an
unobserved bank counterfactual.

There is also an important distinction between small mean squared error
and a valid confidence interval. If the remaining bias is not negligible
relative to the standard error, centering a nominal interval at the biased
comparison does not acquire coverage from having lower mean squared error.
A defended factor-bias envelope can instead be carried into a conservative
interval, or a different inferential procedure can be justified for a
specified factor and sampling regime. The risk table is a design
comparison, not an inference theorem.

Arkhangelsky and coauthors develop synthetic DiD with learned unit and
time weights for a block design. The inspected original's formal
normality result assumes independent identically distributed Gaussian
rows with a well-behaved within-row covariance, particular growth
conditions, spectral restrictions and sufficiently predictive, dispersed
oracle weights. It does not require independent calendar cells within a
row, but its independent-row assumption excludes an unrestricted shared
shock across units. The proof also includes a regularization-rate
condition. Neither the fixed-weight identity
\eqref{exec17:s6-imbalance-product} nor a good numerical pre-fit
establishes those premises.~\citep{synthdid2021}

The same original distinguishes unit bootstrap, a jackknife that holds
estimated weights fixed, and placebo-based variance estimation. Its
jackknife result uses the theorem's conditions and additional restrictions;
validity is not inferred simply because the estimator has an influence
representation. The single-treated-unit case also requires separate care:
the described jackknife is not defined there, while the placebo
variance argument relies on comparability of untreated noise across
units. This is a substantive reason to retain the inference method in
the output, rather than treating bootstrap, jackknife and placebo as
interchangeable settings.~\citep{synthdid2021}

The inspected \texttt{synthdid} documentation requires simultaneous
treatment onset and shows a matrix interface with the number of control
units and pre periods. It also describes the saved interface as beta.
A proposed integration can wrap that declared block route. It should
reject direct admission of staggered adoption, then either choose a
separately justified cohort design or state that the block route does not
answer the request. A mathematical extension discussed in a paper does
not prove that a saved software interface implements it.~\citep{sw_synthdid}

\section{Uncertainty belongs to the signed calendar contrast}
\label{exec17:s6-calendar}

A bank can have millions of account rows and few independent regional
funding shocks. At the same time, the direction of serial correlation's
effect depends on the signs in the requested contrast. Consider fixed
unit coefficients $q_i$ and calendar coefficients $r_t$, each summing
to zero. A double difference has the scalar form
$\widehat\theta=\sum_iq_i\sum_tr_tY_{it}$. Suppose its untreated
noise is
\begin{equation}
 e_{it}=\ell_i b_t+\epsilon_{it},\qquad
 \operatorname{Cov}(b_t,b_s)=\Gamma_{ts},
 \label{exec17:s6-shared-law}
\end{equation}
with independent, centered cell errors $\epsilon_{it}$ of common
variance $\sigma^2$, independent of $b$. The loadings $\ell_i$ are
fixed. Let $\kappa=\sum_iq_i\ell_i$. The shared component of the
contrast is $\kappa\sum_tr_tb_t$, so
\begin{equation}
 \operatorname{Var}(\widehat\theta)
   =\kappa^2 r'\Gamma r+\sigma^2\|q\|_2^2\|r\|_2^2.
 \label{exec17:s6-signed-variance}
\end{equation}
This is a finite-calendar covariance calculation, with no large-sample
claim hidden in the notation. If the unit contrast balances the common
loading exactly, $\kappa=0$, this particular common shock cancels.
The normalization $\sum_iq_i=0$ cancels only equal loadings; it does
not imply $\kappa=0$ when regions have different exposure.

Let $b_t$ be stationary AR(1), with marginal variance $\sigma_b^2$
and correlation $\rho^{|t-s|}$, $|\rho|<1$. For the gap-$h$
difference $r=(-1,0,\ldots,0,1)'$,
\begin{equation}
 r'\Gamma r
  =\operatorname{Var}(b_h-b_0)
  =2\sigma_b^2(1-\rho^h).
 \label{exec17:s6-gap-variance}
\end{equation}
The covariance enters with a minus sign. For a short difference and
positive persistence, nearby common shocks partially cancel. For a
longer gap they cancel less. This observation does not support ignoring
serial correlation; it shows why transferring a variance-of-a-level-mean
intuition to a signed treatment contrast can be wrong in either direction.

An independent check starts from
$b_h=\rho^hb_0+\sum_{j=1}^h\rho^{h-j}\eta_j$, with innovations
independent of $b_0$ and variance
$\sigma_b^2(1-\rho^2)$. Then
\begin{align}
 \operatorname{Var}(b_h-b_0)
 &=\sigma_b^2(\rho^h-1)^2
   +\sigma_b^2(1-\rho^2)\sum_{j=0}^{h-1}\rho^{2j}\nonumber\\
 &=2\sigma_b^2(1-\rho^h).
 \label{exec17:s6-innovation-check}
\end{align}
This derivation verifies the finite covariance result without constructing
a Toeplitz covariance matrix. It also exposes why the stationary
initial-state variance matters. Starting the process at a known fixed
zero state would give a different finite-calendar law.

For one treated region and one donor, set $q=(1,-1)'$ and
$\kappa=1$. With $\sigma_b^2=100$, $\rho=4/5$ and $\sigma^2=1$,
the adjacent difference has common variance 40 and idiosyncratic variance
4, for total 44. The gap-three difference has common variance
$488/5=97.6$ and total $508/5=101.6$. Both use the same two
regions and the same number of observed endpoints. Their uncertainty
differs because of the interval they span, not because one has more
customer records.~\citep{exec17_s6_math}

To examine customer replication explicitly, suppose each region-period
outcome is the mean of $n$ customers with independent unit-variance
idiosyncratic errors and the same regional loading. The idiosyncratic
variance of that mean is $1/n$, so the two total variances become
$40+4/n$ and $97.6+4/n$. As $n$ increases, the limits are 40
and 97.6. This calculation gives an economic reason to acquire more
independent regions or periods when feasible: an additional customer in
an already observed shock cell cannot remove the common component.
It does not prove that any given additional region is independent,
uncontaminated or comparable.

The covariance also interacts with weighting. In the double difference,
$q$ contains the treated and signed donor weights, and $r$ contains
post and signed pre weights. Choosing a weight to balance a factor can
reduce $\kappa$ while increasing $\|q\|_2$ or $\|r\|_2$.
The independent-cell variance formula
\eqref{exec17:s6-sdid-noise} is therefore a special case of a more
specific noise law, rather than a generic cost attached to every
synthetic comparison. If weights are fitted from the same outcome
errors, their randomness adds a further issue beyond the covariance
of a fixed contrast.

An inferential route must specify which dimension supplies independent
replication and what remains conditional. Resampling customers does not
reproduce a region-level common shock. Resampling whole units respects
within-unit trajectories but does not automatically reproduce shocks
shared across all units. A time-block route needs a defended dependence
regime and enough calendar information; a short fixed calendar cannot
be made asymptotically long by changing the resampling option. The
earlier covariance and dependence chapters provide possible methods.
Here the requirement is to evaluate them for the actual signed target,
including its loading imbalance and its information clock.

\section{A response under feedback is a policy-specific response}
\label{exec17:s6-feedback}

A regional time series supplies yet another comparison: forecast the outcome
under the existing operating rule and perturb an action. The intervention must
be specified together with the rule that will react to its consequences.
Consider the structural system
\begin{equation}
 Y_t=\rho Y_{t-1}+\beta A_t+\epsilon_t,\qquad
 A_t=kY_{t-1}+u_t.
 \label{exec17:s6-feedback-law}
\end{equation}
The state $Y_t$ might be a standardized funding position, and $A_t$ a
signed adjustment in servicing or funding intensity. The parameters and
innovation interpretation are supplied structural premises. A fitted
autoregression does not by itself establish that its residual is an
exogenous, manipulable action shock. We compare differences between paths
with the same initial state and structural innovations, so the additive
$\epsilon_t$ cancels in the response calculation.

Introduce a unit action innovation at time zero, $u_0=1$, and set later
innovations to zero. If the operating feedback remains $k$, substitution
into the state equation gives the response
\begin{equation}
 R_h^{\mathrm{maint}}=\beta(\rho+\beta k)^h,
 \qquad h=0,1,\ldots.
 \label{exec17:s6-maintained-response}
\end{equation}
The first state response is $\beta$. Its effect on the next action is
$k\beta$, which adds $\beta k\beta$ to the next state response in
addition to its inherited persistence $\rho\beta$. Repeating this
recursion proves the power in \eqref{exec17:s6-maintained-response}.
The later action responses are $k\beta(\rho+\beta k)^{h-1}$ for
$h\geq1$. A unit innovation is therefore followed by a sequence of
endogenous adjustments; it is not a one-period action path.

If the bank instead forces $A_0=1$ and every later action difference to
zero, the response is
\begin{equation}
 R_h^{\mathrm{open}}=\beta\rho^h.
 \label{exec17:s6-open-response}
\end{equation}
If it replaces the feedback coefficient by $k'$ and then introduces the
same unit innovation, the corresponding response is
\begin{equation}
 R_h^{\mathrm{replace}}=\beta(\rho+\beta k')^h.
 \label{exec17:s6-replacement-response}
\end{equation}
This last calculation assumes that the structural outcome law and innovation
meaning survive the replacement. A rule change that also changes the outcome
mechanism needs a different structural analysis. Moreover, it is the
incremental response to an innovation \emph{within} the replacement regime;
it is not the entire value of switching regimes from an arbitrary nonzero
initial portfolio.

Take $\rho=1/2$, $\beta=1$, $k=2/5$ and $k'=-1/5$. At horizons
zero through three the three state paths are
\begin{align}
 R^{\mathrm{maint}}&=(1,0.9,0.81,0.729),\nonumber\\
 R^{\mathrm{open}}&=(1,0.5,0.25,0.125),\nonumber\\
 R^{\mathrm{replace}}&=(1,0.3,0.09,0.027).
 \label{exec17:s6-feedback-paths}
\end{align}
All three have the same impact response. Their four-period sums are
$3.439$, $1.875$ and $1.417$. At USD 100 of margin per state-period
and an upfront USD 200 implementation cost, the resulting net amounts,
before charging later action intensity, are USD 143.90, $-12.50$ and
$-58.30$. A positive maintained-rule calculation can consequently coexist
with a negative value for the forced one-period action path. The discrepancy
is a difference between interventions, rather than an estimator's rounding
error.~\citep{exec17_s6_math}

For this numerical example the action can be positive or negative without a
capacity bound. A real servicing action may be nonnegative, integer-valued or
capped. Enforcing those restrictions changes the recursion when they bind.
Neither the powers above nor a linear software response should then be
presented as the response of the constrained policy without solving that
policy's state evolution.

\subsection{Closed-loop fit can leave the replacement response unidentified}

Even a perfect fit to the old feedback system need not identify these
interventions. Suppose the observed action is always $A_t=(2/5)Y_{t-1}$,
with no action innovation. Consider two structural worlds:
\begin{equation}
 (\rho,\beta)=(1/2,1)\quad\hbox{or}\quad(1/10,2).
 \label{exec17:s6-two-worlds}
\end{equation}
Both imply $Y_t=(9/10)Y_{t-1}+\epsilon_t$ under the observed rule.
With the same initial law and innovation law they generate the same joint
distribution of all observed states and actions. Yet replacing the feedback
by $k'=-1/5$ gives state coefficients $3/10$ and $-3/10$ respectively.
The observed reduced-form persistence and the known old action rule do not
separate the state persistence from the action effect.

For a known initial state $Y_{-1}=1$ and zero subsequent innovations, the
four replacement-rule states are $(0.3,0.09,0.027,0.0081)$ in the first
world and $(-0.3,0.09,-0.027,0.0081)$ in the second. Their sums are
$0.4251$ and $-0.2289$, or USD 42.51 and $-22.89$ of gross margin at
the same conversion. These are state paths after a controller replacement,
whereas \eqref{exec17:s6-feedback-paths} describes innovation responses
within each regime. Keeping that distinction prevents another target
substitution.

The witness is deliberately a lack-of-excitation case. Exogenous variation
in actions, a defended instrument, a structural restriction on $\rho$ or
$\beta$, or a new experiment can provide additional information. A high
forecast score under the old controller cannot supply it. An estimator
request should therefore carry both the action-support evidence and the
intervention being asked of the time series. When these are absent, a
set-valued answer or refusal is more faithful to the observed information
than a unique replacement-policy response.

\section{Value includes the action path and the unobserved tail}
\label{exec17:s6-value}

The same horizon discipline used for the trend budget applies to feedback
value. For a deterministic zero-baseline example, let $m_h$ be dollars per
state unit, $d\in(0,1]$ a per-period discount factor, $q\geq0$ the price
of squared action intensity, and $K$ an upfront fixed cost. The finite
incremental value is
\begin{equation}
 V_H=\sum_{h=0}^{H}d^h\{m_hR_h-q(\Delta A_h)^2\}-K.
 \label{exec17:s6-finite-value}
\end{equation}
This cost is a stipulated capacity penalty, not a market price. Its
quadratic form is useful for exposing the extra future actions hidden in
a maintained-feedback response. At a nonzero baseline action $A_h^0$,
the change in a quadratic cost is instead
$q\{2A_h^0\Delta A_h+(\Delta A_h)^2\}$. Ignoring the cross term
would generally misprice the incremental policy. Likewise, uncertainty
about actual action paths cannot be priced by squaring only their mean.

For the zero-baseline paths above, the four action responses are
\begin{align}
 \Delta A^{\mathrm{maint}}&=(1,0.4,0.36,0.324),\nonumber\\
 \Delta A^{\mathrm{open}}&=(1,0,0,0),\nonumber\\
 \Delta A^{\mathrm{replace}}&=(1,-0.2,-0.06,-0.018).
 \label{exec17:s6-action-paths}
\end{align}
At $d=1$, $q=10$, $m_h=100$ and $K=200$, their squared action sums
are $1.394576$, $1$ and $1.043924$. Subtracting the corresponding
penalties changes the four-period net values to USD 129.95424,
$-22.50$ and $-68.73924$. The values follow both from summing the
displayed paths and from a separate structural state/action recursion.
They establish an arithmetic comparison under the declared law, not the
economic superiority of a bank policy.~\citep{exec17_s6_math}

For a stable response $R_h=\beta a^h$ with maintained coefficient $k$,
constant margin $m$ and $H\geq0$, geometric sums give
\begin{align}
 V_H={}&m\beta\frac{1-(da)^{H+1}}{1-da}-K\nonumber\\
 &-q\left[1+k^2\beta^2d\,
                   \frac{1-(da^2)^H}{1-da^2}\right].
 \label{exec17:s6-geometric-value}
\end{align}
Here the impact action is one, and $\Delta A_h=k\beta a^{h-1}$
thereafter. The expression is interpreted by continuity when a denominator
vanishes, using the corresponding finite sum. Setting $k=0$ and
$a=\rho$ recovers the open-loop action path; setting $k=k'$ and
$a=\rho+\beta k'$ gives the replacement-regime innovation calculation.
Thus a common formula can compare the three interventions without
quietly carrying the old feedback's actions into the new one.

If $|da|<1$ and $da^2<1$, the infinite discounted value is
\begin{equation}
 V_\infty=\frac{m\beta}{1-da}-K
       -q\left[1+\frac{k^2\beta^2d}{1-da^2}\right].
 \label{exec17:s6-infinite-value}
\end{equation}
These are different convergence requirements. Discounting a state response
can make its benefit series converge while the squared action cost series
still diverges. For example, $d=1/2$ and $a=3/2$ give $da=3/4$ but
$da^2=9/8$. With positive $q$ and nonzero $k\beta$, finite discounted
benefit then coexists with an infinite action cost in this unconstrained
law. A reporting engine should not declare a finite net present value
after checking only the response's first-moment discount condition.

At $d=0.95$, the maintained, open and replacement example gives infinite
net values of approximately USD 473.06, $-19.52$ and $-70.56$.
These calculations use the same known parameters forever. Their large
difference from four-period values is a warning about the economic weight
placed on persistence, not evidence that the far future is predictable.
A structural break, capacity constraint or an unstable parameter in a
confidence region invalidates the corresponding plug-in extrapolation.

\subsection{A conditional truncation budget}

An infinite-horizon claim needs a defended tail restriction. Suppose
uniformly over a declared admissible regime
$|m_hR_h|\leq C r^h$ with known $C\geq0$, $r\geq0$ and $dr<1$.
Then the absolute omitted gross benefit after horizon $H$ is bounded by
\begin{equation}
 T_B(H)=C\frac{(dr)^{H+1}}{1-dr}.
 \label{exec17:s6-tail-budget}
\end{equation}
This follows by applying the triangle inequality to each omitted term
and summing the geometric series. It is an approximation budget
conditional on a uniform structural envelope, rather than a sampling
standard error. With $C=100$, $d=0.95$, $r=0.9$ and $H=3$, it is
about USD 368.55. Reporting the four observed periods to cents does
not reduce this stipulated omitted-benefit budget.

For the linear action envelope
$|\Delta A_h|\leq |k\beta|r^{h-1}$ at $h\geq1$, the omitted
quadratic cost after $H\geq0$ is nonnegative and bounded, if $dr^2<1$,
by
\begin{equation}
 T_A(H)=qk^2\beta^2\frac{d^{H+1}r^{2H}}{1-dr^2}.
 \label{exec17:s6-action-tail}
\end{equation}
Consequently, absent further sign information on the margin response,
the omitted net value belongs to
$[-T_B(H)-T_A(H),\ T_B(H)]$. This uses the known nonnegative
sign of the cost instead of adding a symmetric cost allowance to both
ends. If benefits are also known nonnegative, the lower endpoint can
be tightened to $-T_A(H)$. Such sign information is another premise;
it is not supplied by the word ``benefit.''

If $C=0$, the gross tail is zero at every horizon. For $C>0$, an ex ante
gross-tail tolerance $\varepsilon>0$ and $0<dr<1$,
the smallest admissible nonnegative integer horizon satisfies
\begin{equation}
 H+1\ \geq
 \frac{\log\{\varepsilon(1-dr)/C\}}{\log(dr)}.
 \label{exec17:s6-horizon-budget}
\end{equation}
The reversed inequality follows because $\log(dr)<0$. In the numerical
envelope, a USD 10 gross-tail tolerance needs $H=27$: $H=26$
still has a bound above USD 10, whereas $H=27$ is below it.
This is a requirement for the chosen envelope, not an instruction to
forecast 28 periods (indices $0$ through $27$) from an uninformative short history. More credible
long-horizon restrictions, a finite contractual decision horizon, or
an explicit unresolved value set may be preferable.

In statistical use, a simultaneous confidence region for structural
parameters can be propagated through finite-horizon value and a tail
bound valid uniformly on that region. If the region includes a parameter
with an unbounded adverse tail, its corresponding lower value cannot
be replaced by the value at the stable point estimate. A claim about
the region's coverage remains conditional on identification, dependence
and estimation assumptions; the deterministic tail inequality supplies
none of them. This separation makes it possible to report parameter,
trend, donor, numerical and horizon limitations in their correct units
without calling every budget a confidence interval.

\section{Software routes must preserve the comparison}
\label{exec17:s6-software}

The proposal should reuse a mature estimator where its target and
sampling conditions fit. The relevant choice is more specific than
``run a panel package.'' The required state includes the rollout
calendar, the permitted untreated comparisons, the population and
horizon weights, and the observation or innovation law. The selected
primary documentation supports three complementary routes below.
It is a bounded comparison of inspected versions; the wider software
survey, complete dependency rights and runtime qualification remain
separate obligations.~\citep{exec17_s6_paths}

\subsection{Cohort estimates followed by trend sensitivity}

The first route uses a cohort-aware event-study estimator and passes
a correctly normalized estimate vector with its joint covariance to
a trend-sensitivity procedure. The inspected HonestDiD example for
\texttt{did} requires a dynamic event study, a universal base period,
consecutive relative times and a reference period of $-1$. It
constructs covariance from the dynamic influence vectors and removes
the normalized zero coefficient. Its example target vector selects
one post period. A discounted banking total needs an explicitly
changed target vector and matching outcome units, rather than adding
the endpoints of several separately reported intervals.
\citep{sw_honestdid,exec17_s6_paths}

The bridge's influence-product covariance is not a universal remedy
for shared regional calendar shocks. Its validity depends on the
sampling units represented by the input influence rows and the
underlying estimator's inferential regime. A proposed wrapper must
either justify the joint covariance for the bank's actual design or
decline the requested inferential claim. Retaining the whole matrix
also permits the noisy-slope contrast in
\eqref{exec17:s6-noisy-slope}; exporting only marginal standard
errors would lose that possibility.

The selected \texttt{eventstudyinteract} route is a community Stata
estimator. Its saved \texttt{e(b\_iw)} and \texttt{e(V\_iw)} describe
the weighted estimates; \texttt{e(b)} and \texttt{e(V)} retain the
underlying interacted regression. A subprocess wrapper should export
the intended pair, its cohort weights and support, rather than rely
on whichever object a generic post-estimation command reads. The
manual says to exclude periods after the last cohort becomes treated
when that cohort supplies the controls. It also warns that the cited
validity result is for balanced panels without covariates, although
the implementation accepts unbalanced data. Admitting the broader
input format does not establish the broader causal or inferential
claim.~\citep{v12_sw_eventstudyinteract}

We propose wrapping these existing routes when their regimes and
rights are satisfied. Reimplementing their inference in another
language has no demonstrated benefit at this stage and would create
new numerical and semantic obligations. Stata runtime rights and the
community package's terms must remain distinct. The direct source review in
Section~\ref{sec:exec24-eventstudy-reuse} resolves the earlier missing wrapper
license: the pinned EventStudyInteract snapshot contains an MIT declaration
and text. Runtime entitlement, dependency rights and the intended distribution
still require qualification. These remaining obligations are reasons to
condition reuse, not evidence that the estimator is scientifically inferior.

\subsection{Augmented donors and block synthetic DiD}

The inspected \texttt{augsynth} dispatcher has separate single-outcome,
multiple-outcome and staggered paths. It exposes outcome-model,
balance and regularization choices, with some options specific to
the dispatched path. The saved declaration is version 0.2.0 at its
recorded source commit. We propose evaluating the original R path
and preserving the dispatcher result, donor identities, penalty,
weight signs and chosen uncertainty method. The entry-point manual
does not establish every downstream implementation or inferential
detail. Section~\ref{sec:exec24-augsynth-reuse} develops the selected ridge,
tuning and uncertainty paths, including their remaining numerical limits.
Native admission still needs a qualified environment and matched tests.
In particular, the oracle minimizer in \eqref{exec17:s6-donor-optimum}
should not be used as a claim about the output of its ridge option.
\citep{v12_sw_augsynth,exec17_s6_paths}

The proposed \texttt{synthdid} route is narrower: the inspected
matrix interface uses an untreated-unit count and a pre-period
count with simultaneous treatment onset. The transformation from
long banking observations to that matrix must preserve the meaning
of both axes. Dropping early-treated cells to create a rectangle
changes the observed comparison and may change the population
target. For a staggered request, a collection of separately justified
cohort comparisons is an alternative design, rather than a hidden
preprocessing trick that makes the block estimator universal.
Its beta-interface status also argues for a version-bound wrapper
before considering a port.~\citep{sw_synthdid}

These routes solve related but different balance problems. Augmentation
can spend extrapolation to improve unit balance; synthetic DiD also
changes the pre-period comparison. The product identity
\eqref{exec17:s6-imbalance-product}, norm penalty
\eqref{exec17:s6-sdid-noise} and contamination radius
\eqref{exec17:s6-contamination-radius} explain concrete reasons their
outputs can differ. They do not provide a ranking until the uncertainty
class, sampling law, common target and actual estimator behavior are
specified. We therefore propose retaining both original engines for
matched evaluation. We reject automatic substitution between their
outputs and defer a native port until a measured requirement justifies
its maintenance and verification cost.

\subsection{A prototype can expose semantic errors before a bank trial}
\label{exec17:s6-prototype}

A useful first prototype compares meanings before comparing leaderboard
scores. The exact event-study table supplies a known false-lead
reference, including the full design rank and projection matrix.
The curvature laws supply sharp endpoints and a noisy-slope variance
check. The donor and double-weight tables supply known bias, variance
and contamination tradeoffs. The signed-calendar and feedback laws
then expose two failures that a purely tabular fit score can miss:
the wrong source of replication and the wrong intervention after
a rule change. These are proposed admission tasks, supported by
local authored algebra; no upstream estimator has been run here.
\citep{exec17_s6_math,exec17_s6_design}

Permitted public panels can make the exercise more realistic by supplying
calendar spacing, exposure sizes, missingness and covariate patterns.
Their license, revision date, units and local identity would be retained.
They do not by themselves reveal the untreated outcome of a treated
region. One defensible design overlays explicitly generated potential
outcomes on a permitted public covariate structure, disclosing that
the causal truth is constructed. Another compares descriptive behavior
on a real policy episode while preserving unresolved causal truth.
Those designs answer different questions and should be reported
separately. Existing public-data survey entries can be screened for
these purposes without treating a dataset's availability as scientific
validation of the rollout design.

After semantic admission, a prospective statistical comparison would
freeze each engine and its dependency environment, the common
estimand, a fresh generating law and seed, sample sizes, repetitions,
inference methods, stopping rule and failure accounting. Scenarios
would vary treatment-effect heterogeneity, factor extrapolation,
donor contamination, calendar dependence, action excitation and
post-change stability. Methods should be compared only where their
targets coincide and their claimed regimes are declared; a justified
refusal is an observable outcome. No favorable preliminary result
should be used to revise the final scenario denominator.

The resulting report would retain target error, interval coverage
under the constructed law, availability, width, monetary decision
loss and observed end-to-end cost. It would also retain unsupported
requests, convergence failures, nonfinite outputs and unresolved
rights. Statistical qualification of one frozen regime would remain
different from acceptance of arbitrary panels or a bank's actual
assignment and shock law. Only after such a prototype demonstrates
useful behavior is a costly authentic-bank trial justified by evidence
from this program.~\citep{exec17_s6_design}

For the proposed library, CLI and MCP interface, a single result should
therefore carry the same target vector, reference period, supported
cells, signed weights, uncertainty method, sensitivity restrictions
and horizon budget. A scalar ``effect'' field would erase the
distinctions developed here. The transport should also preserve a
set-valued result or refusal when assumptions or support do not
admit a scalar claim. This is a proposed semantic contract for later
integration; the documented estimator interfaces and these exact
references do not establish interoperability with another repository.



% SOURCE docs/release/2026-10-29/source/docs/writing-repair/2026-10-29/src/1d5bc6fe3595023a.tex CHAR 0
\section{Choosing the unit and calendar weights together}
\label{exec17:s6-joint}

The block comparison above evaluates particular choices of weights. A bank
designing a comparison faces a further question: should it improve the donor
balance, the calendar balance, or both? The two choices interact. A residual
factor mismatch is a product of their imbalances, while independent-cell
noise is multiplied by both squared weight norms. Optimizing one axis while
holding the other fixed need not reveal the best joint choice. We can solve
one small reference problem completely, including its global alternatives,
without claiming that an estimated synthetic-control routine solves that
oracle problem.~\citep{exec17_s6_extension}

Use a new symmetric factor table, rather than change the inputs of the earlier
comparison. There is one treated region with loading 2, two donors with
loadings 0 and 1, one post period with factor 2, and two pre periods with
factors 0 and 1. The untreated mean in each cell is the product of its loading
and factor, plus unrestricted additive unit and time effects. The outcome
errors in the nine cells are independent, centered and have common known
variance $s>0$. A constant treatment effect $\tau$ enters only the treated
post cell. All choices below therefore estimate that same $\tau$; they do
not change the population or horizon to improve a score.

Put donor weights $(1-a,a)$ on the two donors and pre-period weights
$(1-b,b)$ on the two pre periods. Both pairs sum to one, even when a weight
is negative. In row order treated, donor zero, donor one, and column order
post, pre zero, pre one, the signed contrasts are
\begin{equation}
 q=(1,a-1,-a)',\qquad r=(1,b-1,-b)',\qquad
 \widehat\tau=\sum_{i=1}^3\sum_{j=1}^3q_i r_jY_{ij}.
 \label{exec17:s6-joint-contrast}
\end{equation}
Each contrast sums to zero, so the additive unit and time effects cancel.
The coefficient of the treated post cell is one. The remaining factor bias
is $(2-a)(2-b)$. Define $Q(x)=1+(1-x)^2+x^2$, the squared norm of one
axis's contrast. Independence of the nine cell errors gives the exact risk
\begin{equation}
 {\cal R}_s(a,b)
  =(2-a)^2(2-b)^2+sQ(a)Q(b).
 \label{exec17:s6-joint-risk}
\end{equation}
The first term is squared bias for $\tau$, and the second is its variance.
This expression is conditional on the declared factors and fixed weights.
It is a mean squared error, with squared outcome units. It is neither a
confidence radius nor a realized profit. Monetary scaling of an estimation
error would need the exposure and decision loss to be specified separately.

For any fixed $b$, this is a strictly convex quadratic in $a$. Its derivative
and unique conditional minimizer are
\begin{align}
 \frac{1}{2}\frac{\partial{\cal R}_s}{\partial a}
  &=(a-2)(2-b)^2+sQ(b)(2a-1),\nonumber\\
 a^*(b)&=\frac{2(2-b)^2+sQ(b)}{(2-b)^2+2sQ(b)}.
 \label{exec17:s6-joint-conditional}
\end{align}
The denominator is positive for $s>0$. There is a corresponding expression
for $b^*(a)$ by symmetry. When the other axis is already exact, for example
$b=2$, the optimum is $a=1/2$: factor bias is already zero, so the remaining
axis minimizes its noise norm. Matching the factor exactly on both axes
spends variance twice without reducing bias below zero. Conversely, when
the other axis retains imbalance, moving $a$ toward the extrapolating value
2 can have a benefit. Separate coordinatewise convexity does not establish
joint convexity or the uniqueness of a joint solution.

\subsection{A global certificate and two equally good comparisons}

Introduce $x=a-1$, $y=b-1$, $p=x+y$ and $t=xy$. The factor bias is
$1-p+t$. Direct multiplication gives
$Q(a)Q(b)=4(1+p+p^2-t+pt+t^2)$. Completing squares in these symmetric
coordinates yields
\begin{equation}
 \begin{split}
 {\cal R}_s(a,b)
  ={}&\frac{9s}{1+s}+3s(p+t)^2\\
     &+(1+s)\left\{p-t-\frac{1-2s}{1+s}\right\}^2.
 \end{split}
 \label{exec17:s6-joint-certificate}
\end{equation}
This identity is a lower-bound certificate for every real $a,b$. It does
not rely on a numerical search or on convergence of alternating updates.
For $0<s\leq1/2$, both squares can vanish with
\begin{equation}
 p=\frac{1-2s}{2(1+s)},\qquad t=-p,\qquad
 x,y=\frac{p\mathbin{\pm}\sqrt{p^2+4p}}{2}.
 \label{exec17:s6-joint-global}
\end{equation}
Here $p\geq0$, so the two roots are real. Substituting them produces actual
unit and calendar weights attaining the lower bound $9s/(1+s)$. Feasibility
of these roots matters: completing squares in $p,t$ alone would not prove
attainment by real $x,y$ for every noise level. No conclusion for $s>1/2$
is inferred from this zero-square construction.

At $s=1/8$, $p=1/3$ and $t=-1/3$. There are two distinct global solutions:
\begin{equation}
 (a,b)=\left(\frac{7+\sqrt{13}}6,\frac{7-\sqrt{13}}6\right)
 \quad\hbox{or the reversed pair},\qquad
 {\cal R}_{1/8}=1.
 \label{exec17:s6-joint-numeric}
\end{equation}
One coordinate is about 1.7676 and the other about 0.5657. The first
solution extrapolates across donor regions while using positive pre-period
weights; the second extrapolates across calendar periods while using
positive donor weights. Neither axis is exactly balanced. Equal risk comes
from the reference table's symmetry and common cell variance, rather than
from a general equivalence between borrowing information across regions
and borrowing it across time. A bank with common calendar shocks or unequal
exposures would usually break that symmetry.

At the boundary $s=1/2$, the certificate has only one real minimizer,
$a=b=1$, with risk 3. To see uniqueness directly, at $a=1+x,b=1+y$,
\begin{equation}
 {\cal R}_{1/2}(1+x,1+y)-3
   =3(x+y)^2+3x^2y^2.
 \label{exec17:s6-joint-boundary}
\end{equation}
Equality requires both $x+y=0$ and $xy=0$, hence $x=y=0$. The pair of
asymmetric optima at smaller positive noise therefore merges at this
boundary. This is a property of this finite reference objective. It is
not an asymptotic phase transition established for a fitted estimator.

\subsection{What constraints and exact matching cost}

Suppose policy requires nonnegative weights on both axes. Then
$(a,b)\in[0,1]^2$. At $s=1/8$, the first derivative in
\eqref{exec17:s6-joint-conditional} is strictly negative throughout the
square. Indeed $(a-2)(2-b)^2\leq-1$, while
$sQ(b)(2a-1)\leq1/4$, since $Q(b)\leq2$ on that interval. Thus
$\partial{\cal R}_{1/8}/\partial a\leq-3/2$. The same bound holds for
the other derivative. Increasing either coordinate until it reaches one
strictly lowers risk, proving the constrained global optimum $(1,1)$.
It is not merely the best point found on a grid.

The four comparisons on the same nine-cell law are
\begin{center}
\begin{tabular}{lrrr}
\toprule
Choice & $a$ & $b$ & Mean squared error\\
\midrule
Best nonnegative pair & $1$ & $1$ & $3/2$\\
Exact unit balance, minimum time norm & $2$ & $1/2$ & $9/8$\\
Exact balance on both axes & $2$ & $2$ & $9/2$\\
Unrestricted joint optimum & \multicolumn{2}{c}{pair in \eqref{exec17:s6-joint-numeric}} & $1$\\
\bottomrule
\end{tabular}
\end{center}
The exact-unit solution is also the best choice with $a$ constrained to 2:
its zero bias leaves $6sQ(b)$, minimized at $b=1/2$. Requiring both axes
to fit exactly multiplies two norms of 6 and gives $36s=9/2$. For the
nonnegative pair the residual bias is one, and the variance is $4s=1/2$.
The joint optimum accepts some residual bias to reduce the total risk below
the exact-unit alternative. These identities explain the comparisons;
observing a small pre-period residual alone would not do so.

There is a separate direct check. Each of the nine noise coefficients is
$q_i r_j$. Enumerating all $2^9=512$ independent signs for standardized
cell errors gives zero mean and second moment
$\sum_{i,j}q_i^2r_j^2$. Adding the squared factor-table contrast recovers
the displayed risks. For the irrational weights in
\eqref{exec17:s6-joint-numeric}, multiplying their exact cell coefficients
and summing squared coefficients also gives risk one. Thus neither a
floating-point optimizer nor a rounded square root determines the ranking.
\citep{exec17_s6_extension}

The result suggests a useful prototype diagnostic: retain both weight
vectors and the declared noise law when investigating a discrepancy between
two synthetic comparisons. It does not prescribe the oracle weights for an
actual bank. Its post factor and loadings are known here; a real untreated
post factor is counterfactual. Learning the weights from noisy pre outcomes
can correlate them with those errors. A contamination class, common shocks,
missing cells or a constraint on extreme weights changes the objective.
The earlier signed-contamination and calendar calculations show concrete
ways that the apparent gain can disappear. An estimated procedure needs its
own target-preserving selection rule and uncertainty argument. The global
certificate supplies an exact reference for a deliberately restricted law,
with which such a future implementation can be compared.

\section{Valuing replacement of the entire feedback rule}
\label{exec17:s6-full-rule}

The response to a new action innovation within a regime leaves a different
investment question unresolved. A bank may propose replacing the old
feedback rule for an already occupied portfolio, without adding an action
innovation. We now compare those entire regimes from the same nonzero
initial state, charging the actions each would actually take. An apparently
accurate model of the old path still leaves an interval of replacement
values, because the observed controller can hide the separate persistence
and action-effect parameters.

Retain the structural outcome law in \eqref{exec17:s6-feedback-law}, with
future structural disturbances and action innovations zero. Set $Y_{-1}=1$,
$k=2/5$ and the observed reduced coefficient $\lambda=9/10$. Supply the
continuous sensitivity restriction
\begin{equation}
 \beta\in[1,2],\qquad \rho=\lambda-k\beta.
 \label{exec17:s6-controller-class}
\end{equation}
Every world in this class gives exactly the same old state and action path,
for any length of observation under $A_t=kY_{t-1}$. The line
$\rho+k\beta=\lambda$ follows from the old deterministic observation law;
the numerical interval $[1,2]$ is additional supplied information. The
observations do not select that interval, supply probabilities over it or
make it a confidence region. Its use is a sensitivity analysis conditional
on a declared restriction and on structural invariance after replacement.

Replace $k$ by a fixed $k'$ at period zero. Define
$a(\beta)=\lambda+(k'-k)\beta$. Direct recursion gives
\begin{align}
 Y_t^{\rm new}&=a(\beta)^{t+1},&
 A_t^{\rm new}&=k'a(\beta)^t,\nonumber\\
 Y_t^{\rm old}&=\lambda^{t+1},&
 A_t^{\rm old}&=k\lambda^t,\qquad t=0,1,\ldots.
 \label{exec17:s6-full-paths}
\end{align}
Both first actions use the common initial state. There is no impact
innovation of size one and no factor $\beta$ multiplying the initial
state. Those features distinguish these paths from the within-regime
innovation responses developed earlier. If the initial state were zero
under these zero-disturbance premises, both paths would remain zero and
replacement would deliver only its setup loss.

Let $d\in(0,1)$ discount one calendar period, $m$ be dollars of margin per
state-period, $q>0$ be dollars per squared action-period, and $K$ be the
one-time replacement cost. The full paired-regime incremental value through
index $H$ is
\begin{equation}
 \begin{split}
 V_H(\beta;k')={}&\sum_{t=0}^H d^t
 \left[m\{a(\beta)^{t+1}-\lambda^{t+1}\}\right.\\
 &\left.\hspace{20mm}-q\{(k')^2a(\beta)^{2t}-k^2\lambda^{2t}\}\right]-K.
 \end{split}
 \label{exec17:s6-full-finite}
\end{equation}
The old action cost is subtracted from the new action cost, so a saving in
old operating expense raises incremental value. Omitting that old expense
would instead price the new policy against a costless comparator. The
summation contains $H+1$ periods; the setup cost is paid once. The units and
initial state have been normalized explicitly rather than extrapolated
from a local response to an unspecified portfolio.

For $G_H(z)=\sum_{t=0}^Hz^t$, equal to $(1-z^{H+1})/(1-z)$ when
$z\ne1$ and to $H+1$ when $z=1$, the finite formula becomes
\begin{equation}
 \begin{split}
 V_H={}&m\{aG_H(da)-\lambda G_H(d\lambda)\}\\
       &-q\{(k')^2G_H(da^2)-k^2G_H(d\lambda^2)\}-K.
 \end{split}
 \label{exec17:s6-full-geometric}
\end{equation}
This expression is finite for every finite $H$, including a rule whose
infinite discounted cost fails to converge. It makes a short explicit
planning horizon a substantive alternative to assuming a stationary
infinite future. The finite target still needs a defended terminal
treatment if the bank retains assets or liabilities beyond $H$.

\subsection{Robust stability and finite discounted cost are different}
\label{exec17:s6-rule-boundaries}

State stability requires $|a(\beta)|<1$. Since $a$ is affine in $\beta$,
its largest absolute value over $[1,2]$ is attained at an endpoint. Solving
both endpoint inequalities, with the old values $k=2/5$ and
$\lambda=9/10$, gives the robust stability interval
\begin{equation}
 -\frac{11}{20}<k'<\frac9{20}.
 \label{exec17:s6-robust-stability}
\end{equation}
For example, the upper endpoint comes from the $\beta=2$ world:
$9/10+2(k'-2/5)<1$. The corresponding $\beta=1$ bound is weaker.
The lower endpoint follows in the same way by bounding this coefficient
below by $-1$. Equality gives a nondecaying state in an admissible world,
so neither endpoint is included.

Convergence of discounted gross benefits only needs $|da|<1$.
For nonzero $k'$ and positive $q$, convergence of discounted quadratic
action cost instead needs $da^2<1$. The latter is stricter for
$d\in(0,1)$. Write $r_d=d^{-1/2}$. For this class, $r_d>\lambda>0$,
and the endpoint intersection for a uniformly finite discounted value is
\begin{equation}
 k-\frac{r_d+\lambda}{2}<k'<
 k+\frac{r_d-\lambda}{2}.
 \label{exec17:s6-robust-cost}
\end{equation}
The old policy already satisfies its convergence conditions. At
$d=19/20$, the new interval is approximately
$(-0.562989,0.462989)$. It is wider than the robust state-stability
interval. Discounting can make the sum finite while a state grows;
a finite valuation is consequently insufficient as an operational
admissibility criterion. Feasible capital, funding and action limits would
need to be imposed before calling such a rule implementable.

The cost condition cannot be replaced by convergence of the benefit
series. Consider $k'=47/100$. Then
\begin{equation}
 a(\beta)\in[97/100,26/25],\qquad
 d\max a=247/250<1,\qquad
 d(\max a)^2=3211/3125>1
 \label{exec17:s6-rule-divergence}
\end{equation}
at $d=19/20$. All admissible worlds have convergent discounted benefits,
but worlds with
$\beta\geq\{\sqrt{20/19}-9/10\}/(7/100)$ have a nonconvergent
positive action-cost series. In each of those worlds, $V_H$ tends to
$-\infty$ as $H$ grows: the benefit and old-cost totals have finite
limits, whereas the new quadratic expense does not. At equality the
discounted expense remains constant per period, which also makes its
sum diverge. A finite positive number obtained by formally substituting
into a geometric formula outside its convergence range would be invalid.

Where all four relevant series converge, the infinite value is instead
\begin{equation}
 \begin{split}
 V_\infty(\beta;k')={}&m\left\{\frac{a}{1-da}
                       -\frac{\lambda}{1-d\lambda}\right\}\\
 &-q\left\{\frac{(k')^2}{1-da^2}
                       -\frac{k^2}{1-d\lambda^2}\right\}-K.
 \end{split}
 \label{exec17:s6-full-infinite}
\end{equation}
The restriction must accompany the reported value. It is not supplied by
the fact that the old path is stable, since observationally equivalent
old models can cross the boundary differently after the same rule change.

\subsection{Locating the adverse world rather than assuming it}
\label{exec17:s6-rule-value}

Even inside the finite-value region, a quadratic cost can make net value
nonmonotone. Evaluating the two endpoint worlds alone therefore does not
generally give a robust value. For fixed $k'$, differentiating the
finite-valued infinite-horizon expression with respect to $a$ gives
\begin{equation}
 \frac{\partial V_\infty}{\partial a}
  =\frac{m}{(1-da)^2}
   -\frac{2dq(k')^2a}{(1-da^2)^2}.
 \label{exec17:s6-full-derivative}
\end{equation}
The old comparator and setup expense do not depend on $a$. Their absence
from this derivative does not justify omitting them from the level of
value. The negative term can dominate near a squared-cost boundary;
the worst case needs a sign argument or consideration of interior extrema.

For a concrete pair of stable rules, set $m=100$, $q=10$, $K=200$ and
$d=19/20$, and compare $k'=21/50$ with $k'=11/25$. Across both rules
and every $\beta\in[1,2]$, $23/25\leq a\leq49/50$. A single analytic
lower bound is
\begin{equation}
 \begin{split}
 \frac{\partial V_\infty}{\partial a}
 &\geq\frac{100}{\{1-(19/20)(23/25)\}^2}\\
 &\quad-\frac{2(19/20)10(11/25)^2(49/50)}
                   {\{1-(19/20)(49/50)^2\}^2}
 >5829.
 \end{split}
 \label{exec17:s6-controller-monotone}
\end{equation}
The first term uses its smallest value over the interval. The magnitude
of the second term uses its largest value: for $a\geq0$ below the
cost boundary, $a/(1-da^2)^2$ is increasing. Both rules have $k'>k$,
so $a$ also increases with $\beta$. These facts prove that the smallest
infinite value is at $\beta=1$ for each rule, rather than assume that
the largest action effect is always the adverse world.

The same endpoint is adverse at every finite $H$ in this example. For
$t\geq1$, the derivative of the new period's discounted contribution is
\begin{equation}
 d^t a^t\{m(t+1)-2tq(k')^2a^{t-1}\}>0.
 \label{exec17:s6-finite-monotone}
\end{equation}
Here $0<a\leq1$ and $m>2q(k')^2$ suffice; at $t=0$ the derivative
is $m>0$. Thus each summand, and their finite sum, is increasing in
$a$. This finite argument does not depend on an exchange of a derivative
and an infinite series. With a larger cost, a different action range or
a state above one, its sufficient bound would need to be reassessed.

The exact rational values, rounded to cents only for display, are
\begin{center}
\begin{tabular}{rrrrrr}
\toprule
$k'$ & $\beta$ & $a$ & $V_{20}$ & $V_{80}$ & $V_\infty$\\
\midrule
$0.42$ & $1$ & $0.92$ & $-112.57$ & $-92.60$ & $-92.59$\\
$0.42$ & $2$ & $0.94$ & $-4.44$ & $53.68$ & $53.77$\\
$0.44$ & $1$ & $0.94$ & $-5.48$ & $52.61$ & $52.70$\\
$0.44$ & $2$ & $0.98$ & $294.31$ & $580.12$ & $584.45$\\
\bottomrule
\end{tabular}
\end{center}
Values are dollars per normalized initial exposure, including both regimes'
operating expenses and the USD 200 setup cost. Under the supplied class,
rule 0.42 cannot be declared profitable in every world even at an infinite
horizon. Rule 0.44 has positive worst-case infinite value, but its
worst-case value through index 20 remains negative. The table's endpoints
are rigorous extrema for these rules because of the preceding derivative
bounds. They are not percentiles or probabilities of gain.

Notice that the two middle rows share $a=0.94$, but their values differ.
The rule coefficients are different, so their action expenses differ even
when their new state paths coincide. This is a concrete reason to retain
the action path when translating a predicted state improvement into money.
A model returning only a state forecast could hide the operating-cost
distinction. Similarly, a comparison of $V_{20}$ with $V_\infty$ changes
the investment target; it is not a numerical refinement of the same
finite-horizon quantity.

\subsection{A tail certificate for the whole paired comparison}

The earlier innovation tail cannot be reused unchanged for this question.
Subtracting \eqref{exec17:s6-full-geometric} from
\eqref{exec17:s6-full-infinite} leaves four tails:
\begin{equation}
 \begin{split}
 V_\infty-V_H={}&m\left\{
 \frac{a(da)^{H+1}}{1-da}
 -\frac{\lambda(d\lambda)^{H+1}}{1-d\lambda}\right\}\\
 &-q\left\{\frac{(k')^2(da^2)^{H+1}}{1-da^2}
 -\frac{k^2(d\lambda^2)^{H+1}}{1-d\lambda^2}\right\}.
 \end{split}
 \label{exec17:s6-full-tail}
\end{equation}
For the positive coefficients in the two-rule example, the triangle
inequality gives an absolute bound by adding these four positive
magnitudes. Maximizing the new-state and action terms at $a=0.98$ and
$k'=0.44$ yields a uniform bound of USD 4.3532, below USD 4.36, at $H=80$.
The actual largest endpoint tail is about USD 4.3231; cancellation makes
it smaller than the four-term bound. A uniformly safe truncation allowance
should use the defended envelope, rather than an observed cancellation
at a convenient world.~\citep{exec17_s6_extension}

Directly evolving old and new states, taking their separate actions and
accumulating discounted gains reproduces the finite geometric formula.
Adding the exact remaining tails reproduces the infinite value, for both
endpoint worlds and an interior world with $\beta=3/2$. These checks also
retain the first action and its cost. They would fail if the calculation
were silently replaced by a within-regime innovation response or if the
initial state were set to zero.

The decision implication remains conditional. A positive worst-case value
for rule 0.44 is warranted only for this supplied structural class, action
feasibility, initial exposure, cost law and horizon. The old observations
cannot certify the class or structural invariance. An exogenous action
perturbation, a credible restriction on persistence, or a justified
controlled pilot could narrow the class; a public panel can supply
realistic clocks and exposures without supplying that missing causal
truth. These are concrete next information choices. Until they are
resolved, the proposal should return the class of values and the
convergence restrictions alongside any chosen policy, so that an
attractive forecast is not mistaken for an identified bank investment.


% SOURCE docs/release/2026-10-29/source/docs/writing-repair/2026-10-29/src/47033306ecbad01f.tex CHAR 60215


\section{From the comparison laws to implementation choices}
\label{exec17:s6-end}

A panel or time series supplies a comparison only after the target,
counterfactual information and sampling law are specified. The exact
rollout regression produces positive leads without anticipation.
A target-specific trend budget becomes more demanding with the
horizon and with uncertainty in the pre-period slope. Exact donor
balance can lose to a regularized comparison, while an explicit
contamination class can reverse that ranking. Synthetic DiD's two
axes of balance interact with noise and shared shocks. A maintained
feedback response can value a different intervention from the one
the bank proposes to implement.

These results motivate concrete reuse and prototype routes. They
also identify the information still needed: credible counterfactual
restrictions, appropriate independent replication, action excitation,
stable intervention meaning, feasible costs and a defended horizon.
An integrated proposal should preserve those requirements through
the estimator, uncertainty calculation and monetary decision. A
successful software call cannot replace them.

The following studies apply that requirement to choices that the opening
examples held fixed: the baseline, a competing intervention, the units that
survive estimation and the covariance actually returned. This extends the
comparison to the implementation boundary without treating a source finding
as an executed package result.


% SOURCE docs/release/2026-10-29/source/docs/writing-repair/2026-10-29/src/6f605269c7484f80.tex CHAR 0
\section{Panel comparisons when the intervention, baseline and population differ}
\label{sec:exec28-panel-designs}

A branch rollout is not a single causal contrast until its comparison population and information clock are fixed.
Suppose a bank introduces a servicing platform across regions while a fee concession arrives on a different calendar.
The sponsor wants the platform's incremental contribution to customer cash flow. An event-study graph may instead describe
a change relative to a moving baseline, a combined platform-and-concession effect, or an average over a changing set of regions.
These are different questions even when every output column is called an ATT. The design therefore precedes the choice of R,
Stata or a language-neutral adapter. Three primary treatments make the distinction concrete: event-study normalization,
multiple events, and balanced stacked comparisons.~\citep{exec28_panel_roth,exec28_panel_tsai,exec28_panel_wing}

\subsection{What the event-study vertical axis measures}

Let treatment start at calendar time $t=1$, let $r=t-1$ denote event time, and let
$\Delta_t=E[Y_t\mid D=1]-E[Y_t\mid D=0]$ be the observed treated-minus-control mean gap.
In the common-adoption setting, a dynamic two-way fixed-effects comparison with period zero as its reference has
$\beta_r=\Delta_{r+1}-\Delta_0$. Roth's selected version shows why this interpretation cannot automatically be applied
to every pre-treatment coefficient from another estimator. A varying-base construction uses $\Delta_{r+1}-\Delta_r$
before treatment and $\Delta_{r+1}-\Delta_0$ afterwards. An imputation construction in his example uses the earliest
pre-period for pre-treatment coefficients but the mean pre-treatment outcome for post-treatment coefficients.
Different baselines create different graphical shapes even without a treatment effect.~\citep{exec28_panel_roth}

Consider an explicitly synthetic banking example, measured in dollars per customer per period. Treatment has zero effect,
but the untreated group gap is $\Delta_t=50t$ for $t=-4,\ldots,2$. This violates parallel trends deliberately. The last
pre-period is zero, the earliest gap is $-200$, and the mean pre-period gap is $-100$. At calendar times $-2,-1,0,1,2$,
the fixed-base contrasts are $-100,-50,0,50,100$. The varying-base contrasts are $50,50,50,50,100$ where the adjacent
pre-period exists. The earliest-base/mean-base imputation contrasts are $100,150,200,150,200$.
The apparent downward jump at adoption is a change in reference, not a negative platform effect. All these values
are independently recomputed from exact fractions.~\citep{exec28_panel_reference}

The practical remedy is to publish the comparison formula with the graph. In Roth's common-adoption example, a universal
pre-period base makes the Callaway--Sant'Anna construction numerically match the fixed-base comparison. Using the mean
pre-period on both sides gives $\Delta_t+100$ here, a constant vertical shift of the fixed-base series.
That pre-period series averages to zero by construction, so its average cannot independently establish parallel trends.
Roth also distinguishes in-sample pre-treatment residuals from post-treatment prediction errors: in his common-adoption
analysis the former can shrink toward zero. A small pre-period residual is not a general validation of the missing
counterfactual.~\citep{exec28_panel_roth}

This does not invalidate the recent estimators when their identifying assumptions hold. It addresses their diagnostic
behavior when those assumptions fail. The selected Roth paper records installations on 31 October 2025 and distinguishes
an earlier from an updated multiple-groups implementation. The separately pinned January 2024 example script has unpinned
package imports and cannot establish that later environment. A proposal should retain the diagnostic construction and
environment separately, and require a dependency lock before a future native replay. For the bank, the useful output is
the cohort path, its baseline and uncertainty, not merely an attractive common plot.~\citep{exec28_panel_roth,exec28_panel_software}

\subsection{A second event changes both the controls and the target}

Write $G^1$ for platform adoption and $G^2$ for concession adoption. Following the selected two-event framework, each
treatment is absorbing; reversal is outside this design. Let $Y_t(g^1,g^2)$ denote the potential outcome under both adoption
dates, with zero in an argument meaning no adoption of that intervention. For the joint cohort $G=(g^1,g^2)$ define
\begin{align*}
 C_G(t)&=E[Y_t(g^1,g^2)-Y_t(0,0)\mid G],\\
 A_G(t)&=E[Y_t(g^1,0)-Y_t(0,0)\mid G].
\end{align*}
The combined effect $C_G$ and platform-only effect $A_G$ coincide before the concession arrives, under no anticipation.
They need not coincide afterwards. Tsai's first stage uses controls not yet exposed to either event, with baseline
$b=\min(g^1,g^2)-1$. Under an \emph{unconditional} parallel untreated-change restriction for those groups,
\[
 C_G(t)=E[Y_t-Y_b\mid G]-E[Y_t-Y_b\mid G^1>t,G^2>t].
\]
With conditional parallel trends, one must instead adjust and standardize over the target covariate distribution; a raw
unconditional difference is not justified merely by a conditional assumption. Positivity must hold for the selected joint
cohort and comparison observations. The paper develops outcome-regression, inverse-probability and doubly robust moments.
Its sampling setup permits serial dependence within a unit but assumes independent sampling across units, a substantive
limitation for common banking shocks.~\citep{exec28_panel_tsai}

The source's homogeneous-effects decomposition explains why ignoring the concession can matter. If the platform effect is
$\tau_1$, the concession effect is $\tau_2$, and the joint interaction is $\kappa$, the ordinary target-event contrast,
under the additional untreated-change condition needed for that contrast, has the form
\[
 \tau_1+\tau_2(p_T-p_C)+\kappa q_T.
\]
Here $p_T,p_C$ are the fractions newly exposed to the concession between baseline and outcome in the treated and control
groups; $q_T$ is the fraction of the treated group exposed at the outcome date. For synthetic values $\tau_1=40$,
$\tau_2=-30$, $p_T=3/4$ and $p_C=1/4$, the ordinary comparison is $25$, not $40$. If $\kappa=20$ and $q_T=3/4$,
the contrast returns to $40$ by cancellation. Agreement with a favored answer therefore does not diagnose correct
specification. This is a homogeneous illustrative submodel, not identification of unrestricted interactions.
\citep{exec28_panel_tsai,exec28_panel_reference}

Removing contaminated controls identifies the combined effect; it does not automatically separate the interventions.
For $g^1<g^2$ and $t\geq g^2$, a derivation based on Tsai's Appendix A.4 uses the platform effect immediately before the
second event and its subsequent change among units with the same platform cohort but a later concession:
\[
 A_G(t)=C_G(b)+\sum_{s>t}\pi_s\{C_{(g^1,s)}(t)-C_{(g^1,s)}(b)\},\qquad b=g^2-1.
\]
The weights $\pi_s$ are the eligible later-concession cohort shares within the same platform cohort and sum to one.
No anticipation identifies the anchor and donor platform effects; parallel platform-effect changes across concession
cohorts justify transporting the increment. If the anchor is $40$ and the donor effect rises from $45$ to $55$, the
target is $50$. A current combined effect of $85$ must not replace the anchor: doing so gives $95$.
This comparison follows the appendix's explicit anchor and is independently derived here.~\citep{exec28_panel_tsai,exec28_panel_reference}

When $g^2<g^1$, set $b=g^1-1$ and use units with the same concession cohort but a later platform date. Under additive
effects and parallel concession-effect changes across platform cohorts,
\[
 A_G(t)=C_G(t)-C_G(b)-\sum_{s>t}\rho_s\{C_{(s,g^2)}(t)-C_{(s,g^2)}(b)\}.
\]
For synthetic combined effects $95$ and $20$ and a donor concession change from $20$ to $55$, this yields $40$.
The nuisance increment cancels. If concession trajectories differ between the selected cohorts, that cancellation fails.
If adoption dates coincide, this design lacks a period displaying just one effect and does not identify the separate
platform effect. A combined-policy estimand may still be useful, but must be named accordingly.~\citep{exec28_panel_tsai,exec28_panel_reference}

Source precision matters before porting. The selected v3 has inconsistent ordering language near its second stage, and
equation (11)'s target-period entry differs from the pre-concession anchor in Appendix equation (17).
The formulas here follow the appendix's explicit argument and event order; they do not silently declare every printed
expression interchangeable. Section~\ref{sec:exec28-panel-delegates} subsequently inspects the actual
\texttt{fastdid} multiple-event weight path and finds the pre-confounding anchor used in this derivation.
That resolves the selected branch's anchor, while a later-edition comparison, the broader theorem and native
implementation equivalence remain separate questions. The paper names \texttt{fastdid}; pinned repository metadata
independently establishes the author and multiple-event scope. Speed and additional-event claims remain documentation
claims here, not measured performance or a proof for unrestricted treatment histories.~\citep{exec28_panel_tsai,exec28_panel_software}

Inference must retain the full two-stage transformation. For fixed aggregation weights $w$ and first-stage covariance
$\Sigma$, the aggregate variance is $w^{\mathsf T}\Sigma w$, not a sum that discards overlap. Estimated cohort weights
add their own first-order contribution. Tsai invokes smooth parametric nuisance models, linear representations, finite
moments and a multiplier bootstrap. These are conditions to examine, not an automatic justification for dependent branches,
arbitrary machine-learning nuisance fits or rare joint cohorts. The adapter should expose first-stage cells, both event
dates, donor eligibility, aggregation weights and influence functions so the meaning survives serialization.~\citep{exec28_panel_tsai}

\subsection{Balancing an event window does not itself balance stack weights}

Let $a$ index an adoption cohort and $\mathcal A$ retain cohorts for a prespecified event window with adequate clean
controls. Let $n_a^T,n_a^C$ count treated and control units in each sub-experiment, including repeated control appearances
in different stacks. The trimmed treated-share target is
\[
 \theta_e=\sum_{a\in\mathcal A}p_a^T\tau_{a,e},\qquad p_a^T=\frac{n_a^T}{\sum_j n_j^T}.
\]
Keeping $\mathcal A$ and its weights fixed across event time prevents changing composition from masquerading as dynamics.
Constant cohort effects $20$ and $60$ average to $40$ when both are equally represented, but to $20$ when only the early
cohort remains. No effect has decayed. Trimming changes the target population and should be reported as a design choice.
\citep{exec28_panel_wing,exec28_panel_reference}

Wing, Freedman and Hollingsworth show a second problem in the basic pooled stacked comparison. Within each stack, let
the untreated change be $d_a$ in both arms and define $p_a^C=n_a^C/\sum_jn_j^C$. The pooled change contrast is
\[
 \sum_a p_a^T(\tau_{a,e}+d_a)-\sum_a p_a^C d_a
 =\theta_e+\sum_a(p_a^T-p_a^C)d_a.
\]
Parallel trends within every sub-experiment do not cancel the last term when the arms average sub-experiments differently.
This concerns the specified pooled construction; it does not equate every regression described as stacked DiD.
\citep{exec28_panel_wing}

In a synthetic branch example, two stacks contain treated counts $(10,30)$, control counts $(90,10)$, effects $(20,60)$,
and common within-stack untreated changes $(0,100)$. Treated shares are $(1/4,3/4)$ and control shares $(9/10,1/10)$.
The intended effect is $50$. The treated mean change is $125$ and the control change is $10$, so the unweighted pooled
estimate is $115$. Its $65$ excess comes entirely from mismatched trend weights. Even zero causal effects would leave
a false $65$ signal.~\citep{exec28_panel_reference}

Corrective sample weights are one for treated rows and $q_a=p_a^T/p_a^C$ for control rows. Because
$\sum_a p_a^Cq_a=1$, weighted controls have the treated cohort shares. Here control weights are $5/18$ and $15/2$,
the corrected control change is $75$, and the estimate is $50$. Exact row-level means agree with the aggregate algebra.
Large ratios expose weak support and precision costs; they cannot invent controls. A future implementation must reject
zero denominators and distinguish group-count, customer-population and other targets.~\citep{exec28_panel_wing,exec28_panel_reference}

Repeated controls also preserve dependence across stacks. Treating copies of the same branch as independent creates
information absent from the original data. The paper discusses original-unit and unit--stack clustering, and reports a
limited Monte Carlo comparison. This does not establish universal coverage for cross-branch shocks or few policy clusters.
Original IDs must survive stacking, and the covariance specification must state which common shocks it permits.
Correcting point-estimate weights alone does not qualify uncertainty.~\citep{exec28_panel_wing}

\subsection{Conditional reuse across the R and Stata routes}

The inspected stacked tutorial contains R and Stata examples, sample-share weights, and regression calls clustered on
the original state identifier. It is a useful prototype template subject to a pinned environment and membership/weight
tests, not a complete adapter. Its R inclusion expression also contains an \texttt{adoptionTime == TRUE} branch alongside
the intended treated, late and missing-date branches. Numeric-date encoding is therefore a specific review obligation.
The Stata weight example sums within-group row indices, proportional to counts only under appropriate balance.
Balanced and unbalanced fixtures are needed before translation. The complete selected tree has no named license file;
that bounded observation is not a legal conclusion about all possible permissions.~\citep{exec28_panel_software}

The nonabsorbing route addresses a different question. Pinned \texttt{nonabsdid} 0.4.1 metadata describes a wrapper around
\texttt{DIDmultiplegtDYN}, \texttt{PanelMatch} and \texttt{fect}, with an optional \texttt{fixest} reference.
A shared plotting schema is useful, but does not equate targets, histories or uncertainty procedures. Consider adopting
output conventions while inspecting every delegate and retaining method-specific fields. Direct adoption of the comparison
as a common estimand is premature. The metadata declares MIT plus a file identifying the copyright holder; dependency
rights and complete behavior remain separate checks.~\citep{exec28_panel_software}

Pinned \texttt{pampe} 1.1.2 documentation describes a single affected unit, donor outcome matrix, subset regression selection
and counterfactual prediction, with control and time placebos. Its documented defaults use AICc and a \texttt{leaps}
dependency. This suggests a comparator for a regional intervention with a credible stable donor relation, not a cohort ATT
obtained by renaming output. The manual's model-size discussion and actual delegate need reconciliation; the original
panel-approach argument and package paper require separate inspection. Metadata declares GPL-2. An isolated callable
comparison is a proposal option, not an approved port or redistribution decision.~\citep{exec28_panel_software}

These choices imply a concrete CLI/MCP contract: original unit IDs, outcome and event clocks, competing interventions,
reversal policy, baseline, clean-control and cohort rules, aggregation population, covariance assumptions and package identity.
A result must expose its supported target and excluded cells as well as estimates. An LLM can help elicit intervention
calendars and explain why a proposed control is already exposed; it cannot infer unobserved parallel trends or certify a
causal diagram from a plot. No acquired statistical engine or sibling repository was executed for these comparisons.

The source boundary is explicit. Roth is exact arXiv v2 submitted 16 January 2026 with a 21 January cover; Tsai is exact
v3 submitted 21 January 2025 with a 22 January cover. The stacked original is the retained January 2024 NBER working paper
from a mutable endpoint; its cover does not prove historical byte identity. The cutoff remains 7 October 2026.
Later-edition clarification, remaining originals, delegate inspection, dependency rights and rendered visual review remain
local obligations. The reviewed panel additions do not close the overall software or proposal program.


% SOURCE docs/release/2026-10-29/source/docs/writing-repair/2026-10-29/src/bdc867023b273723.tex CHAR 0
\section{Keeping the population and comparison fixed}
\label{sec:exec28-panel-delegates}

The preceding design comparison becomes actionable only when the selected software carries the same baseline, support
and event order. Ten further pinned files expose those decisions in \texttt{fastdid}, \texttt{nonabsdid} and \texttt{pampe}.
Their value is not that a function returns a number, but that its observable selection and transformation rules can be
compared with the banking question. This is a static source comparison with independent arithmetic, not execution of the
acquired packages.~\citep{exec28_panel_delegates}

\subsection{The multiple-event anchor is explicit in the implementation}

In the selected \texttt{fastdid} commit, the event-specific weight routine separates a direct case before the competing
event, an imputation case when the target arrives first, and a difference-in-differences case when it arrives later.
For target-first imputation it assigns the deterministic weight to the target cohort at the pre-confounding base period,
then adds the eligible donor change. Thus this inspected branch agrees with the explicit Appendix A.4 anchor developed
earlier, rather than equation (11)'s inconsistent target-period entry in the selected v3. That resolves one concrete
implementation question; it does not resolve the different edition named by the repository or verify its general
multiple-event theorem.~\citep{exec28_panel_delegates}

The routine also intersects donor cohort identities available at both periods. This matters independently of the anchor.
Suppose baseline donors A and B have effects 10 and 20, while outcome donors B and C have effects 25 and 60. Separately
normalizing available donors gives a change of $55/2$; the common donor B changes by only 5. Transporting $55/2$ would
mix a change in composition with a change in effects. The inspected code restricts both sides to the common identities
and warns and skips the cell if none remain. Its target-later branch matches already-arrived competing event dates and
requires other events to remain sufficiently far away. At equal first/earliest-competing event dates, no separating branch
supplies a nonzero weight. These are useful refusal semantics to retain in a bank-facing response.~\citep{exec28_panel_delegates,exec28_panel_delegate_reference}

The aggregation routine transforms both effects and influence functions, including estimated cohort-share contributions
in the second-stage and final aggregation. That is stronger than merely multiplying reported standard errors by weights.
Nevertheless its covariance depends on downstream bootstrap and estimation delegates not executed here. Selected validation
rejects doubly robust estimation with an unbalanced panel, time-varying covariates with an unbalanced panel, and additional
clustering or uniform bands without bootstrap. It may drop units with missing competing-event dates, incomplete panels or
first-period exposure. The first-stage driver can skip failed or very small cells. An adapter must report those exclusions
and the resulting target population; it must not quietly substitute zeros for refused cells.~\citep{exec28_panel_delegates}

Two additional checks remain important. Calendar normalization requires uniformly spaced periods and shifts every event
date consistently. The dynamic balance selector checks cohort extrema, but that alone is not proof that every interior
event-time cell survived earlier estimation failures. Full membership checks are therefore required for a promised common
window. The two-stage control routine uses both \texttt{double\_control\_option} and \texttt{control\_option} in different
branches; their combinations need an explicit fixture matrix before the options are exposed as interchangeable concepts.
The proposed reuse is a pinned reference integration with visible support and exclusions, conditional on delegate, test,
rights and native-environment verification.~\citep{exec28_panel_delegates}

\subsection{A common wrapper still carries method-specific choices}

The selected \texttt{nonabsdid} front door dispatches to DCDH, PanelMatch and three \texttt{fect} modes. Preflight verifies
named columns and binary treatment, can encode nonnumeric IDs, rejects entirely missing control columns, and explicitly
warns that partial missingness may yield different effective samples. It does not itself impose a common sample or prove
common estimands. Its PanelMatch path allows reversals and chooses propensity-score matching when controls are supplied,
otherwise no refinement. The \texttt{fect} path uses method-specific tuning defaults: five cross-validation rounds,
an IFE factor cap of two and a matrix-completion lambda grid of five when those values are not supplied. Such changes are
substantive analysis settings and belong in the result provenance.~\citep{exec28_panel_delegates}

The DCDH bridge constructs a parsed call with literal horizon arguments and binds data in an evaluation environment. Its
comments attribute this to a particular upstream version's call inspection; that claimed upstream behavior has not been
independently replayed here. A CLI/MCP wrapper should admit known option names and values and pin the upstream interface,
not assume an arbitrary parameter bag is safe or semantically stable. Output conversion also needs care: the effect-cell
helper creates normal intervals from standard errors when bounds are missing, and infers a standard error from interval
width when one is absent. That reverse calculation has the intended meaning for symmetric normal intervals, not arbitrary
bootstrap or simultaneous intervals. The product should retain native bounds and their method, plus label any derived
normal approximation.~\citep{exec28_panel_delegates}

\subsection{Stack-count balance and a preserved donor timeline}

The earlier stacked tutorial supplies a useful caution for ports. Its Stata helper sums within-unit row indices. For
$L$ observations of a unit, that sum is $L(L+1)/2$. If every unit in every stack has the same $L$, the factor cancels
from each normalized share: the earlier weights $5/18$ and $15/2$ are unchanged. If two treated units in different stacks
have lengths two and three, while both controls have length two, triangular-count weights are $(2/3,4/3)$; pooled-row
weights are $(4/5,6/5)$; at a common event period with one unit per arm in each stack, weights are $(1,1)$.
These calculations establish that the conventions differ outside balance. They do not establish which unbalanced-panel
estimand is identified. Enforce the declared balanced design or derive and test the alternative explicitly.
\citep{exec28_panel_delegate_reference}

The selected R tutorial's additional comparison to \texttt{TRUE} can include numeric adoption date 1 even if the focal
date is 5 and the clean-control cutoff is 7. A transcribed predicate check admits 1 whereas the intended treated-or-late
rule excludes it. This is a source-level encoding witness, not a native R execution report or an assertion that the
tutorial's calendar-year dataset triggers the branch. Calendar origins and sentinel values need explicit normalization
and rejection tests before the example becomes a general adapter.~\citep{exec28_panel_delegate_reference}

For \texttt{pampe}, the selected estimation function fits subset regressions on supplied pre-treatment rows and forms the
counterfactual from the selected OLS model. It rejects missing values and a nonsingle treated-unit request. Its AICc helper
includes the variance parameter in $k$, giving the penalty $2k+2k(k+1)/(n-k-1)$; for $n=8,k=4$ this is $64/3$.
The denominator must remain positive. The top-level path restricts model size, while placebo paths have their own local
choices and require separate checks. Metadata about AICc does not settle all small-sample paths.~\citep{exec28_panel_delegates,exec28_panel_delegate_reference}

The leave-one-control-out robustness helper reconstructs pre-period rows as one through the fitted-sample length instead
of retaining the originally supplied indices. A fit on rows 5--12 therefore leads the helper to use rows 1--8. It also refits
the retained donor subset rather than repeating the original full model-selection procedure. These are different
diagnostics. The repair contract is to store exact pre/post row indices and information dates in the result, validate
their reuse, and name fixed-subset sensitivity separately from selection-aware sensitivity. The selected main return
object does not itself retain those original index vectors. An isolated comparison adapter is plausible after these
contracts are repaired and checked; unrestricted adoption or a cross-language port remains premature.
\citep{exec28_panel_delegates,exec28_panel_delegate_reference}

For the bank, these details govern whether a servicing-platform result still refers to the approved cohort, timeline,
target and uncertainty after moving through an agent or another repository. Preserve options, excluded cells, original
IDs, interval provenance and selected model rows. The inspected source permits a concrete integration plan and precise
tests. It does not establish native equivalence, complete dependency rights, empirical coverage or bank acceptance.


% SOURCE docs/release/2026-10-29/source/docs/writing-repair/2026-10-29/src/a63777ec26daeadf.tex CHAR 0
\section{What an event-study estimate represents}
\label{sec:exec28-panel-estimator-contracts}

A branch-policy comparison can survive a file conversion while losing its
scientific meaning. A bank that moves a servicing-policy panel from Stata into
an R wrapper must preserve treatment codes, missingness reasons, observation
clocks, the population contributing to each contrast and the origin of every
interval. A common column layout is useful only after these distinctions are
made explicit. The following comparison uses nine additional files at the
previously fixed \texttt{fastdid} and \texttt{nonabsdid} commits. It is a static
implementation review, accompanied by independent arithmetic, and does not
claim that either package or its imported numerical dependencies was executed
\citep{exec28_panel_estimator_sources,exec28_panel_estimator_reference}.

\paragraph{Effective support and nuisance estimation.}
In the selected balanced \texttt{fastdid} estimator, the control contrast uses
normalized propensity odds when covariates and the chosen IPW or doubly robust
option activate that branch. Outcome adjustment fits a weighted linear model
to the control units' outcome changes. The code includes nuisance-estimation
terms in its influence calculation; thus an adapter that exports only a
coefficient and discards the fitted weighting law cannot reconstruct the same
uncertainty calculation. The unbalanced branch instead forms separately
normalized pre- and post-period means, weights propensity fitting by presence
in those periods and fits separate outcome regressions. Those branches need
distinct declared observation laws and tests. Their shared function name is
not a proof of equivalent targets or calibration.

The selected balanced branch warns about extreme fitted propensities and caps
them at $1-10^{-10}$. The corresponding control odds can still be
$9{,}999{,}999{,}999$ before sampling-weight multiplication. A warning and a
finite numerical value therefore provide no useful-overlap guarantee. That
branch also replaces missing or magnitude-above-$10^{10}$ entries of its
scaled covariance matrix by zero. The proposal must expose this alteration,
retain the original diagnostic and require a declared refusal or regularization
policy; it cannot quietly interpret zeroed entries as established certainty.
The selected unbalanced branch does not contain the same replacement at its
corresponding covariance step. Dependency versions, convergence and singular
design behavior remain separate qualification questions.

For weights $w_i$, the source calculates the effective support
\[
 n_{\mathrm{eff}}=\frac{(\sum_i w_i)^2}{\sum_i w_i^2}.
\]
It rejects a group, or an unbalanced group-period, when this quantity is
nonfinite or below two. Two equally weighted observations give
$n_{\mathrm{eff}}=2$; weights $(1,2)$ give $9/5$ and trigger the selected
refusal despite there being two physical observations. The residual component
is multiplied by $\sqrt{n_{\mathrm{eff}}/(n_{\mathrm{eff}}-1)}$ after this
check. For two equal weights its squared multiplier is two. This algebra
describes the implemented adjustment, not a general finite-sample coverage
theorem for estimated weights, correlated branches or selected cells. A public
comparison should report physical counts, effective counts, refused cells and
surviving aggregation weights separately. The upstream propensity solver,
weighted least-squares behavior and multiplier bootstrap still require their
own pinned review and matched execution contract.

\paragraph{Calendar and argument conversion.}
The auxiliary \texttt{fastdid} path normalizes uniformly spaced times and
rejects irregular spacing. Calendar-valued group and time outputs are mapped
back, while dynamic event-time coordinates remain in normalized period units.
Sampling weights and time-invariant inputs are taken from the first unit row
in the unbalanced auxiliary construction. A bank must therefore define which
period's exposure weight represents a unit, and must not silently treat a
one-period event coordinate as one calendar day. Custom aggregation expressions
also require an explicitly controlled interface. The proposed CLI/MCP should
expose named, validated aggregation choices rather than pass arbitrary business
text into the package's expression option.

The \texttt{nonabsdid} alias layer rejects simultaneous alias and canonical
arguments. It maps \texttt{effects} to \texttt{leads=effects-1} and
\texttt{placebo} to \texttt{lags}; its selected validation checks numeric range
but then uses integer coercion. At that layer, an input of $2.9$ effects becomes
two effects and one lead, while $1.9$ placebos becomes one lag. These are
source-level conversion witnesses, not native invocation results. The proposed
boundary should require finite integer horizons before dispatch and return
both requested and resolved options. The DCDH tidier subtracts one from the
upstream time coordinate and reads estimates from embedded plot data. A call
that omits the plot is refused at the inspected extraction boundary. Preserve
the native clock alongside the display clock and qualify the exact upstream
object version instead of treating a visible plot slot as a permanent API.

\paragraph{Intervals, reference rows and file round trips.}
The \texttt{fect} tidier carries native interval bounds and takes standard
errors only when the designated matrix has a matching number of rows. This
length check does not independently establish that its row identities match
the event-time vector. Its method-label helper maps missing and unrecognized
labels to \texttt{IFE}; an integration should retain the original label and
flag an unknown method rather than infer that a named estimator ran. The
PanelMatch tidier requires named estimate coordinates, handles separate
post-treatment and placebo fields and can append a zero-estimate, zero-error
row at time $-1$. That row is a plotting reference, not an observed exact-zero
effect, and must be excluded from empirical coverage and precision denominators.
Existing native bounds, normal approximations and simultaneous bands must keep
distinct provenance through the common constructor; that constructor and
upstream versions remain explicit follow-up inspections.

The selected Stata reader defaults to converting labelled values into factors
and collapsing tagged missing values to ordinary missingness. Those defaults
may be convenient for presentation but are not a reversible preservation of a
numeric treatment code or a reason-for-missingness taxonomy. Preprocessing must
retain those mappings before fitting. The writer sanitizes names, truncates
them to 32 characters and then makes them unique. Two names sharing the same
32-character prefix consequently produce a 34-character second name after
the selected suffix operation. This is a pre-writer boundary witness; whether
the exact imported writer rejects or transforms it has not been executed.
A portable schema should allocate collision-safe names within the length bound
and preserve a reverse mapping.

These findings support a conditional adapter design: preserve raw and resolved
arguments, native and displayed clocks, treatment labels, missingness tags,
effective samples, influence provenance and reference-row status. They do not
support an unconditional port or adoption decision. Full upstream numerical
and constructor inspection, exact dependency and rights profiles, the broader
multiple-event paper edition, and native public fixtures remain prerequisites
to stronger claims. Unknown counterfactual trends remain scientific assumptions
that a faithful file or API conversion cannot establish.


% SOURCE docs/release/2026-10-29/source/docs/writing-repair/2026-10-29/src/cc7a4cc6470c77b8.tex CHAR 0
\section{Commercial difference-in-differences methods for staggered adoption}
\label{exec18:gp-panel}

A staggered branch rollout poses a different question from the cross-sectional
graphical query. Branches can adopt a policy at different dates, the policy's
effect can change with exposure, and calendar shocks can move all outcomes.
A commercial panel command should therefore be compared on the treatment
history, cohort-time target, admissible controls and uncertainty procedure.
A numerical coefficient without these meanings is insufficient for a banking
decision. The inspected Stata manuals distinguish scalar linear DID and
heterogeneous cohort-time effects; the EViews introduction supplies a more
limited panel-DID capability statement.
\citep{exec18_gp_stata_did,exec18_gp_stata_hdid,exec18_gp_stata_xthdid,exec18_gp_eviews}

The scalar Stata pair is \texttt{didregress} for repeated cross-sections and
\texttt{xtdidregress} for panels. The first fits time and group fixed effects;
the second fits time and panel fixed effects. Their heterogeneous counterparts
are \texttt{hdidregress} and \texttt{xthdidregress}, respectively. The latter
manuals describe regression adjustment, inverse-probability weighting,
augmented weighting and an extended two-way fixed-effects approach. Listing
these alternatives does not make their assumptions or output interchangeable.
The selected heterogeneous panel formulas assume that treatment remains on
after adoption. A history \((0,1,0)\) therefore requires refusal under that
profile or a separately justified switching-treatment design.
\citep{exec18_gp_stata_did,exec18_gp_stata_hdid,exec18_gp_stata_xthdid}

Let \(G_i\) denote branch \(i\)'s first adoption date, with infinity for a
never-treated branch. Write \(Y_{it}(g)\) for its outcome if adoption begins
at \(g\), and \(Y_{it}(\infty)\) for its untreated outcome. A cohort-time
target and one possible aggregation are
\begin{equation}
\label{exec18:gp-cohort-target}
 \theta(g,t)=\mathbb E\{Y_{it}(g)-Y_{it}(\infty)\mid G_i=g\},
 \qquad \Theta_{\omega}=\sum_{g,t\geq g}\omega_{gt}\theta(g,t),
 \quad \omega_{gt}\geq0,\quad\sum_{g,t\geq g}\omega_{gt}=1.
\end{equation}
This notation is our explicit target convention. Its identification needs
the selected comparison assumptions, including consistency, an appropriate
no-anticipation condition, overlap and untreated parallel trends for the
declared comparison population. The inspected manuals allow never-treated
or not-yet-treated controls and document base-time choices for preperiod
effects. Those choices alter which contrasts and support are available.
Not-yet-treated units cease to be eligible untreated comparisons after
adoption. The request must therefore retain the control rule, anticipation
window, base period and the branch/time population before aggregation.
\citep{exec18_gp_stata_hdid,exec18_gp_stata_xthdid}

\subsection{An exact positive-effect rollout can have a negative projection}
\label{exec18:gp-panel-example}

Consider five equally weighted hypothetical branches, \(E_1,E_2,L_1,L_2,N\),
observed at dates \(0,1,2\). Number them \(i=0,\ldots,4\) in that order.
Their untreated outcomes, in thousands of stipulated USD per branch-period,
are \(Y_{it}(\infty)=10i+2t\). The early cohort adopts at date 1, the late
cohort at date 2, and \(N\) never adopts. There is no anticipation, spillover,
attrition or missing period. Treatment adds 1 at the early cohort's first
treated date, 100 at its second, and 1 at the late cohort's treated date.
These are deliberately extreme finite values that make a target error
visible, rather than forecasts of a bank rollout.
\citep{exec18_gp_panel_reference}

\begin{center}
\begin{tabular}{lrrrr}
\hline
Cohort and date & Branches & Effect & ATT weight & TWFE weight\\
\hline
Early, date 1 & 2 & 1 & \(1/3\) & \(5/8\)\\
Early, date 2 & 2 & 100 & \(1/3\) & \(-1/8\)\\
Late, date 2 & 2 & 1 & \(1/3\) & \(1/2\)\\
\hline
\end{tabular}
\end{center}
Here ATT means the average over the six treated branch-periods, so it is
\((2+200+2)/6=34\) thousand USD. Never-treated DID contrasts recover the
three effects exactly: compare each early date with date 0 and the late
treated date with date 1, subtracting the corresponding change in \(N\).
These are finite arithmetic equalities under the stipulated untreated law.
They carry no sampling confidence statement.
\citep{exec18_gp_panel_reference}

Now fit the scalar least-squares projection
\(Y_{it}=\alpha_i+\lambda_t+\beta D_{it}+e_{it}\) to the same 15 rows,
where \(D_{it}=1\{t\geq G_i\}\). In this balanced panel, let bars denote
branch, date and grand means. Removing branch and date fixed effects gives
\begin{equation}
\label{exec18:gp-fwl}
 \widetilde D_{it}=D_{it}-\overline D_i-\overline D_t+\overline D,
 \qquad
 \widehat\beta=\frac{\sum_{i,t}\widetilde D_{it}Y_{it}}
                         {\sum_{i,t}\widetilde D_{it}^{\,2}},
 \qquad \sum_{i,t}\widetilde D_{it}^{\,2}=\frac{16}{15}.
\end{equation}
The treatment functional weights the three cohort-time effects as
\begin{equation}
\label{exec18:gp-negative-weight}
 \widehat\beta=\frac58(1)-\frac18(100)+\frac12(1)
              =-\frac{91}{8}=-11.375,
 \qquad \operatorname{ATT}=34.
\end{equation}
The weights sum to one, but one is negative. The negative projection and
positive treated-period average answer different questions. Exact rational
residualization and a separate solution of the full normal equations agree;
the treatment row of the inverse Gram matrix independently reproduces these
weights. No Stata or EViews engine was used. This example proves a discrepancy
for its specified panel, rather than a universal failure of fixed effects
or a claim about a vendor's undocumented default aggregation.
\citep{exec18_gp_panel_reference}

If the cost is stipulated to be USD 10,000 per treated period, the true
retrospective mean net increment under this law is USD 24,000. Mistaking
the projection for that ATT would instead produce minus USD 21,375.
The latter is an explicitly invalid value calculation. Neither quantity
evaluates a new rollout policy, estimates bank prices or measures ROI.
An investment comparison still needs the applicable population, treatment
costs, feasible alternative and uncertainty about its transported outcomes.
\citep{exec18_gp_panel_reference}

\subsection{Aggregation, trends and uncertainty belong to the request}

The heterogeneous postestimation manual distinguishes aggregation by cohort,
calendar time, exposure time, and cohort/time, and distinguishes
\texttt{weights(timecohort)} from \texttt{weights(cohort)}. The former builds
weights using the subsets used for the individual ATETs; the latter uses
cohort sizes. A branch closure or a changing observed sample can therefore
matter to the interpretation of an aggregate. The adapter should return
the supported cohort-time cells and their denominators first, then the
chosen normalized weights and the resulting estimand. Replacing those
weights by a convenient average can change the banking population.
\citep{exec18_gp_stata_post}

Preperiod diagnostics are useful checks, but cannot validate the unobserved
postperiod counterfactual. For an exact counterexample, let two groups have
identical paths at dates \(-2,-1\). At adoption, suppose the treated group's
true treatment effect is 2, and an otherwise unobserved group-specific
untreated-outcome shock is 3. The observed postperiod DID is 5. A world
with effect 5 and no differential shock has the same observed contrast.
Thus even exact preperiod equality cannot distinguish those worlds without
an additional restriction on the untreated postperiod path. A nonrejected
pretrend test is weaker still. Sensitivity to differential trends and
anticipation must remain part of the business interpretation.
\citep{exec18_gp_panel_reference}

Inference is likewise command-specific. The inspected \texttt{xthdidregress}
manual equates \texttt{vce(robust)} with clustering at the declared treatment
group. That is not a generic promise of independent branch-level rows or
a universal interpretation of every command's robust option. The
postestimation manual disallows its multiplier simultaneous intervals after
TWFE and disallows \texttt{estat aggregation} after bootstrap or jackknife
estimation. A requested interval may target one cell, one fixed aggregate
or simultaneous coverage of many cells; those events differ. A wrapper
must check the selected estimator, clustering, interval target and option
compatibility before returning uncertainty. The five-branch arithmetic
example supplies no cluster asymptotics, multiplier qualification or
confidence coverage. A future generated multi-cluster trial must declare
its sampling law and coverage target separately.
\citep{exec18_gp_stata_xthdid,exec18_gp_stata_post,exec18_gp_panel_reference}

\subsection{What the commercial evidence supports and what to build first}

The acquired Stata originals contain July/August 2026 creation metadata and
StataNow feature flags. They were acquired on 18 October 2026; neither their
metadata nor a displayed copyright authenticates the precise bytes or
entitlements available at the 7 October cutoff. An edition-specific
comparison needs the release and options actually under consideration.
The EViews page, displaying a March 2025 update, states that built-in DID is
provided for panel workfiles and introduces a TWFE and diagnostics chapter.
That introductory page is insufficient evidence of every referenced
algorithm, complete covariance implementation or licensed product edition.
These source limits qualify the historical comparison; they do not imply
that an uninspected vendor lacks a feature.
\citep{exec18_gp_stata_did,exec18_gp_stata_hdid,exec18_gp_stata_xthdid,exec18_gp_stata_post,exec18_gp_eviews}

Our bounded candidate search also retains SAS, IBM SPSS, MATLAB and
commercial frontends to open-source tools as unresolved families.
DID-specific primary originals and exact editions are missing for those
comparisons. Generic panel regression is not sufficient evidence of a
causal DID workflow. The immediate comparison is consequently a documented
Stata baseline and a limited EViews lead, with explicit omissions rather
than a claim of a saturated commercial market. Embedding rights, prices,
support and executable compatibility remain separate selection questions.

The saved local dispatcher probe returned an unsupported disposition for
its dependent-panel DID request. This is an observed limitation of that
request and revision, not evidence that every panel research module is absent.
Chapter~\ref{gap29:public} separately exercises fitted-summary inference through
library, CLI and MCP; it does not qualify a raw-panel DID estimator or erase
the probe's refusal. A proposed extension should admit a
fully declared cohort-time target, preserve the treatment and information
history, report cells and weights, and reject unsupported reversals or
inference options. A licensed vendor comparison can then reproduce the
15-row reference, add unbalanced support and alternative controls, and
examine an independently designed uncertainty trial. Completing this
bounded path would provide a meaningful prototype comparison; it would
not by itself establish general trend validity or bank qualification.
\citep{exec18_gp_dispatch,exec18_gp_panel_reference}


% SOURCE docs/release/2026-10-29/source/docs/writing-repair/2026-10-29/src/7dda824a5d83b3d7.tex CHAR 0
\section{Choosing a fixed-effect engine without changing the bank's question}
\label{exec20:fx-scope}
A bank comparing outreach across branches and months needs more than a fast
regression. It must know which account-periods remain in the calculation, whether
the treatment varies after the fixed effects are removed, what variation supports
uncertainty, and which customers an aggregated event-study effect describes.
These requirements motivate conditional reuse of the original \texttt{fixest}
engine through a typed R worker. The scope initially proposed here is linear
fixed-effect estimation; an IV formula, a nonlinear family or an event-study
expansion needs its own scientific and numerical qualification.

The inspected revision is commit \texttt{789f888416e9} of 26 June 2026, whose
\texttt{DESCRIPTION} declares version 0.14.2. This identifies the selected source,
without establishing the latest or a distributed release. Its R dependency is
at least 3.5.0, with Rcpp linking and imports including \texttt{numDeriv},
\texttt{nlme}, \texttt{sandwich}, \texttt{dreamerr} and \texttt{stringmagic}.
The declaration is GPL-3; the associated standard license text is available
separately. A root \texttt{LICENSE} file is absent from the inspected tree.
That absence does not erase the declaration or establish complete rights for
an assembled distribution. A separate process is an engineering choice and
does not settle distribution obligations. The selected change log records
covariance-class, IV-collinearity and degrees-of-freedom fixes; a workflow
specifies checks on three operating systems. These are maintenance evidence,
not observed passing checks or a support commitment.~\citep{exec20_fixest_source}

\subsection{Projection, causal meaning and the retained population}
\label{exec20:fx-projection}
Let $Y$ be the vector of contribution margins before the incremental outreach
cost, $D$ the binary outreach indicator and $Z$ the full-rank matrix of branch
indicators. With equal observation weights, write
\begin{equation}
 M_Z=I-Z(Z^{\mathsf T}Z)^{-1}Z^{\mathsf T},\qquad
 \widehat\beta=(D^{\mathsf T}M_ZD)^{-1}D^{\mathsf T}M_ZY.
 \label{exec20:fx-within}
\end{equation}
The expression requires $D^{\mathsf T}M_ZD>0$. It is the coefficient obtained
by ordinary least squares on the explicit dummy matrix $[Z,D]$: minimizing
over branch intercepts first gives the residual objective
$\|M_Z(Y-D\beta)\|^2$, whose scalar normal equation yields
\eqref{exec20:fx-within}. For strictly positive weights $W$, the corresponding
argument applies to $W^{1/2}Y$, $W^{1/2}D$ and $W^{1/2}Z$. Consequently a
conversion that changes weights changes the projection being requested.

This algebra alone identifies no intervention effect. For example, a
constant-effect panel interpretation additionally needs a credible restriction
such as $\mathbb E[\varepsilon_{it}\mid D_{i1},\ldots,D_{iT},Z]=0$ in the
specified model, as well as consistent treatment, an appropriate interference
assumption and the required sampling conditions. Outreach triggered by an
unobserved deterioration in customer finances can violate that restriction.
An IV extension substitutes instrument relevance and substantive exclusion
and exogeneity assumptions; the formula syntax cannot supply them. Likewise,
adding time effects changes $Z$ and may remove the variation that previously
identified the projection. These are scientific choices that the proposed
business-to-model interface must state before submitting a regression.

Consider the following exact authored example, with two periods per branch.
The outcomes are hypothetical USD contributions per account-period.
\begin{center}
\begin{tabular}{c|cc|cc}
Branch & $D_1$ & $D_2$ & $Y_1$ & $Y_2$\\\hline
A & 0 & 1 & 10 & 12\\
B & 1 & 0 & 20 & 16\\
C & 0 & 1 & 30 & 36
\end{tabular}
\end{center}
Within-branch residualization gives $D^{\mathsf T}M_ZD=3/2$ and
$\widehat\beta=4$. Solving the explicit dummy normal equations gives the
same result, with intercepts $(9,16,31)$. Reversing input row order leaves
the coefficient and group scores unchanged when row identifiers travel with
their data. If instead treatment is one for every observation in B and zero
elsewhere, $D$ lies in the span of $Z$ and $M_ZD=0$. The coefficient is
unavailable; reporting zero would invent an answer to an unidentified
regression target. These are rational-arithmetic reference cases,
not results from running R.~\citep{exec20_fixest_reference}

The inspected \texttt{feols} documentation describes recursive removal of
singletons and perfect-fit observations under its default setting. Its stated
slope invariance in the relevant removal setting does not imply invariant
inference or population reporting. The source also tracks collinear variables
and observation selections. For finite weights it rejects negative values
and removes zero weights; missing or infinite weights enter sample-selection
handling. Our initial adapter should prevalidate all weights and require an
explicit policy for exclusions, rather than rely on a message printed after
estimation. Every retained and removed row needs its original identifier and
reason. A factor level disappearing after exclusions can change rank, and a
dropped treatment must remain a typed unavailable result.~\citep{exec20_fixest_source}

\subsection{A covariance request is part of the estimand service}
\label{exec20:fx-covariance}
The selected startup code sets the default covariance family to iid. Later
calls can read session options. Thus a formula containing branch fixed effects
does not by itself request branch-clustered errors. A reproducible worker
must set the covariance family, cluster identifiers, small-sample correction
and reference degrees of freedom explicitly. The source's \texttt{ssc}
interface separates $K$ adjustment, fixed-effect counting, cluster adjustment
and the degrees of freedom used for the Student reference distribution.
The default fixed-effect count is \texttt{nonnested}; \texttt{none},
\texttt{full} and optional exact counting express different conventions.
Its upstream tests describe an irregular two-way example with counts 6, 1,
8 and 7 under these choices. Those tests are useful future conformance
specifications and have not been executed here.~\citep{exec20_fixest_source}

For the three-branch example let $\widetilde D=M_ZD$,
$\widehat u=M_ZY-\widetilde D\widehat\beta$ and
$s_g=\sum_{i\in g}\widetilde D_i\widehat u_i$. The uncorrected one-way
cluster sandwich for this scalar coefficient is
\begin{equation}
 \widehat V_0=\frac{\sum_gs_g^2}{(\widetilde D^{\mathsf T}\widetilde D)^2}
 =\frac{(-1)^2+0^2+1^2}{(3/2)^2}=\frac89.
 \label{exec20:fx-cluster}
\end{equation}
An explicitly chosen correction
$\widehat V_K=\widehat V_0\{(n-1)/(n-K)\}\{G/(G-1)\}$ gives
$4/3$ for $n=6,G=3,K=1$ and $10/3$ for $K=4$. The coefficient remains 4.
Choosing $G-1$ gives two reference degrees of freedom; it is distinct from
either variance scaling. This calculation isolates the choices an adapter
must preserve, without asserting that every native option or weighting law
uses these exact formulas. Three clusters provide no evidence of calibrated
coverage. More generally, dependence within a branch and between branches
must be matched to the uncertainty procedure.~\citep{exec20_fixest_reference}

Numerical covariance repair adds another distinction. In the inspected code,
\texttt{vcov\_fix} defaults to true. A nonpositive eigenvalue triggers a
candidate repair; assignment of the repaired covariance occurs inside
\texttt{if(vcov\_fix)}. An adjacent comment says repair always occurs even
when the flag is false, but the executable branch has the narrower behavior.
The warning decision uses a maximum change greater than $10^{-8}$, so an
absence of a warning is not proof that nothing was changed. A proposed
adapter must preserve the requested policy and the original and repaired
matrices, or refuse to certify a repair state it cannot observe. Positive
definiteness is a numerical property; it does not validate a causal model,
cluster law or interval coverage.~\citep{exec20_fixest_source}

\subsection{Event-time support and the meaning of an aggregate}
\label{exec20:fx-eventtime}
Suppose branches adopt a servicing policy in different calendar months.
The cohort is the adoption month and relative time is the observed month
minus that cohort. The inspected \texttt{sunab} documentation warns that,
when relative time is inferred from calendar values, a cohort absent from
the supplied period support is treated as never treated. Its example uses
cohort 1994 and observed years 1995--2005: an already-treated cohort can be
misclassified as never treated. The proposed conversion must therefore
construct treatment status from the business history before invoking the
helper. It must distinguish never treated, treated before observation and
not yet treated, and preserve reference periods, binning and unsupported
cells. A valid design further needs the substantive comparison assumptions,
including appropriate parallel trends and no anticipation; a matrix expansion
does not verify them.~\citep{exec20_fixest_source}

The selected aggregation routine constructs shares from the signs of design
columns, using observation weights when requested, normalizes those shares
and averages the corresponding coefficients. It calculates uncertainty from
the full coefficient covariance. For a fixed vector of shares $a$ and
coefficient vector $\widehat b$, the target and conditional variance are
\begin{equation}
 \widehat\theta=a^{\mathsf T}\widehat b,\qquad
 \widehat{\operatorname{Var}}(\widehat\theta)=a^{\mathsf T}\widehat V a.
 \label{exec20:fx-aggregate}
\end{equation}
The second identity follows by expanding the covariance of a linear
combination. If shares themselves are estimated for a different population
target, their uncertainty needs an additional argument; treating the observed
shares as fixed does not settle that question.~\citep{exec20_fixest_source}

For an authored contrast take $\widehat b=(2,8)^{\mathsf T}$,
$a=(3/4,1/4)^{\mathsf T}$ and
$\widehat V=\left(\begin{smallmatrix}4&1\\1&9\end{smallmatrix}\right)$.
Then $\widehat\theta=7/2$ and the variance is $51/16$. Dropping the covariance
cross term gives $45/16$; averaging the two standard errors gives $9/4$,
which is not $\sqrt{51}/4$. Equal coefficient shares would instead give an
effect of 5. At a stipulated incremental cost of USD 4 per account, these two
different targets imply net margins of USD $-0.50$ and USD 1. Neither is an
observed bank return; the contrast shows why silently changing aggregation
weights can reverse a proposal's recommendation.~\citep{exec20_fixest_reference}

\subsection{A conditional reuse choice and a public qualification path}
\label{exec20:fx-choice}
The original R engine is the preferred proposal route because its sample,
design and covariance behavior can be kept together behind one typed boundary.
\texttt{feols}, \texttt{feglm} and \texttt{femlm} are distinct model surfaces
within that package; qualifying one linear profile does not qualify all of
them. \texttt{pyfixest} remains a separately assessed implementation.
An explicit-dummy reference is attractive for small transparent examples,
while a direct port of numerical routines would also have to reproduce
formula coding, row selection, fixed-effect rank, covariance defaults,
aggregation and failures. No measured cost or target-platform requirement
currently justifies that duplication.~\citep{exec20_fixest_source,exec20_fixest_reference}

The proposed request carries the immutable data snapshot, units, row keys,
calendar, design, target population and uncertainty choices. The response
carries coefficient identities, covariance, retained sample, exclusions,
rank and convergence diagnostics, aggregation shares and any repair status.
Library, CLI and MCP surfaces must preserve those meanings. A missing target,
invalid request, unsupported model, numerical failure and cancelled job
remain different outcomes. An agent may propose a reviewed design, but an
initial tool should accept bounded formula structures rather than arbitrary
generated R expressions. These are proposed contracts; no such fixest worker
has been admitted by this source review.

Public and generated panels can test the merit of that design before a bank
experiment. Start with exact dummy/within agreement and an absorbed treatment,
then test row permutations, zero/missing weights, singleton cascades,
near-collinearity, covariance defaults and correction choices. Event-time
cases must include adoption outside the observed calendar and aggregation
with nonzero covariance. Native runs should compare outputs, masks and typed
failures against the authored references and independently inspected upstream
specifications. A complete pinned distribution, compiler/runtime environment
and dependency rights review precede executable admission. Subsequent public
data studies may assess numerical stability and workflow value; they cannot
establish a private bank's assignment mechanism or actual return.

The scientific owner must own identifying assumptions and dependence; the
data owner must own calendars, weights and exclusions; the adapter maintainer
must own conversions and upgrade conformance. Operational costs to measure
include worker startup, serialization, memory, rank and covariance computation,
dependency maintenance and diagnostics. This conditional proposal selection
can be completed now while native qualification and deployment remain open.
It converts a broad intention to reuse panel software into a testable design
without treating an available coefficient as a validated causal conclusion.


% SOURCE docs/release/2026-10-29/source/docs/writing-repair/2026-10-29/src/2247b08365fedf14.tex CHAR 0
\section{Panel software: timing, weighting and uncertainty}
\label{exec23:panel-scope}

A bank comparing staggered branch interventions with a simultaneous regional
launch faces two different designs. A Python fixed-effect regression and an R
synthetic difference-in-differences fit can both return a scalar measured in
dollars, but that coincidence does not establish that they estimate the same
policy value. The reusable object must describe the population, adoption clock,
control eligibility, outcome horizon, weighting and uncertainty. This section
examines concrete interfaces in PyFixest and the original R synthdid package,
building on the separate R fixest discussion. Source inspection establishes
what the selected code says and how a proposed adapter must represent it; it
does not establish that an installed package has passed our trials.
\citep{exec23_pyfixest_reuse_source,exec23_synthdid_reuse_source}

The PyFixest snapshot is commit
\texttt{45d52b6e368cefc7321239921f2547b4e08c8f20}, dated 6 October 2026.
Its Python version is dynamic through Maturin; the Rust core declares 0.60.0.
The original synthdid snapshot is commit
\texttt{70c1ce3eac58e28c30b67435ca377bb48baa9b8a}, dated 15 January 2024,
and declares version 0.0.9. The different dates matter for maintenance planning:
the latter README describes a beta interface, while the former has changing
Python and native numerical boundaries. A cutoff-selected commit and a package
version are identities, not promises of future support. Neither Python syntax
similarity to R fixest nor the use of the name SDID by other implementations
transfers a qualified capability.
\citep{exec23_pyfixest_reuse_source,exec23_synthdid_reuse_source}

\subsection{A fixed-effect coefficient can reverse a policy conclusion}
\label{exec23:panel-projection}
Let $D_{it}$ indicate treatment and $Y_{it}$ be a monetary outcome for unit $i$
at period $t$. In a balanced panel, residualizing $D$ against unit and time
effects gives
\begin{equation}
 \widetilde D_{it}=D_{it}-\bar D_{i\cdot}-\bar D_{\cdot t}+\bar D,
 \qquad
 \widehat\beta_{\rm TWFE}=
 \frac{\sum_{it}\widetilde D_{it}Y_{it}}
      {\sum_{it}\widetilde D_{it}^{\,2}}.
 \label{exec23:panel-twfe}
\end{equation}
This is a projection identity. To interpret it as the intended average causal
effect requires additional restrictions on treatment effects, counterfactual
trends and comparisons. The fixed-effect labels in a formula do not supply
those restrictions.

An exact constructed example makes the distinction visible without estimating
anything in PyFixest. Stipulate zero untreated potential outcomes and use
\begin{equation}
 D=\begin{pmatrix}0&1&1\\0&0&1\end{pmatrix},
 \qquad
 Y=\begin{pmatrix}0&1&10\\0&0&1\end{pmatrix},
 \qquad
 \widetilde D=\frac16
 \begin{pmatrix}-1&2&-1\\1&-2&1\end{pmatrix}.
 \label{exec23:panel-witness}
\end{equation}
All three treated-cell effects are positive and their average is $4$. The
denominator in (\ref{exec23:panel-twfe}) is $1/3$, so the projection is $-7/2$.
Its weights on the six outcomes are
$(-1/2,1,-1/2;1/2,-1,1/2)$. In particular, the large positive late effect
receives a negative weight. The example intentionally has no never-treated
unit and no untreated comparison in its final period; it demonstrates the
projection's meaning, not a successful causal design. With a stipulated cost
of $2$ dollars per treated cell, the target net is $2$ dollars but treating the
projection as that target would report $-11/2$ dollars. These numbers reproduce
with rational arithmetic and are neither native package outputs nor measured
bank returns. \citep{exec23_pyfixest_reference}

In the inspected PyFixest event-study dispatcher, the concrete TWFE branch
constructs a treatment indicator, optional covariates and unit/time effects.
The stored ATT flag does not itself generate a lead--lag formula in that
branch. DID2S and the saturated implementation follow different fitting and
aggregation paths. The product therefore needs an explicit estimator profile
and a displayed estimand, rather than an event-study label that conceals the
formula actually fitted. \citep{exec23_pyfixest_reuse_source}

\subsection{Calendar conversion, first-stage support and aggregation}
\label{exec23:panel-clock}
PyFixest's common DID base accepts numeric time and cohort fields but constructs
relative time by subtraction. A calendar code is not necessarily a period
index: $202601-202512=89$, whereas the elapsed month count is one. The proposed
adapter should retain real calendar dates and an explicit consecutive index,
including a policy for missing observations. The base constructs absorbing
treatment from the declared cohort and uses zero for never treated. Its event
origin is the first observed treated period. If declared adoption is period
two and the first observed post-treatment record is period three, the declared
event time is one while recentering on the observed record makes it zero. Both
dates must survive conversion so that an observation gap cannot silently
change the scientific question.
\citep{exec23_pyfixest_reuse_source,exec23_pyfixest_reference}

The DID2S helper verifies two unique numeric treatment values and then casts
them to Boolean. Codes two and three both map to true; the adapter must
validate the actual binary convention before that conversion. Its first stage
fits untreated observations. Missing fixed-effect levels in that sample lead
to missing predictions and an explicit error. This is meaningful support
information: a bank should retain which branches or calendars could not be
represented, instead of deleting the error from the request denominator. The
GMM covariance uses residuals from both stages. The wrapper retains the
second-stage inference degrees of freedom and disables specified
post-estimation capabilities; availability of a generic fitted-object method
does not establish that it is valid for this two-stage result.
\citep{exec23_pyfixest_reuse_source}

The saturated branch also needs a narrower contract than its common public
surface might suggest. The inspected dispatcher does not pass the covariate
formula to its saturated helper. The actual aggregation implementation
accepts period and shares choices, despite broader options in its docstring;
its period-weight code includes relative period zero despite contrary prose.
Its linear-combination interval uses the normal multiplier $1.96$. This is
not automatically a few-cluster correction or a simultaneous event-time band.
Those source discrepancies are reasons to test and constrain the chosen
profile, not empirical claims that a newly executed native request failed.
\citep{exec23_pyfixest_reuse_source}

Aggregation is itself part of the estimand. For effects
$\widehat\theta=(10,-2)'$, weights $w=(1/4,3/4)'$ and the stated covariance
\begin{equation}
 V=\begin{pmatrix}4&1\\1&9\end{pmatrix},
 \qquad w'\widehat\theta=1,
 \qquad w'Vw=\frac{91}{16},
 \label{exec23:panel-aggregation}
\end{equation}
the equal-cohort effect would instead be $4$. Omitting cross-covariance gives
$85/16$, a different variance even with unchanged marginal standard errors.
The example assumes fixed aggregation weights. If the weights are estimated,
their additional uncertainty must be justified for the chosen procedure.
The reporting schema should retain cohort sizes, eligible populations, full
covariance and coefficient order so an analyst can see why the summaries
differ. \citep{exec23_pyfixest_reference}

PyFixest offers useful numerical machinery: its demeaning wrapper checks the
success flag and raises on nonconvergence. But a profile must also retain
singleton removal, analytic-weight defaults, rank tolerance, solver choice,
backend-specific preconditioner fallback and warnings. Rust, within, optional
Numba and Torch paths are not one interchangeable execution guarantee.
The earlier weighted-IV source discussion remains relevant to its IV branch;
fixed-effect fitting and first-stage diagnostics do not by themselves provide
weak-instrument confidence-set inversion.
\citep{exec23_pyfixest_reuse_source}

\subsection{The original synthetic difference-in-differences block}
\label{exec23:panel-sdid}
The original synthdid package addresses a different layout. There are $N_0$
control units followed by $N_1$ treated units, and $T_0$ pre-treatment periods
followed by $T_1$ post-treatment periods. Its treatment indicator equals one
only in the bottom-right block. This is the block design identified in the
original paper's inspected discussion; broader extensions in the literature
are not evidence that this package supports staggered adoption.
\citep{synthdid2021,exec23_synthdid_reuse_source}

Let $Z=Y-X\beta$ denote outcomes after the selected covariate adjustment, and
let $\omega$ and $\lambda$ be control-unit and pre-period weights. The code
computes the bilinear contrast
\begin{equation}
 \widehat\tau=
 (-\omega',N_1^{-1}\mathbf1_{N_1}')
 Z
 (-\lambda',T_1^{-1}\mathbf1_{T_1}')',
 \label{exec23:panel-contrast}
\end{equation}
so treated units and post periods receive equal weights within that block.
A branch-period effect is not automatically a customer-weighted portfolio
effect. If customer populations differ, the proposal must either defend this
target or derive and qualify an alternative instead of changing weights
after seeing the answer. The code's scalar result alone cannot establish
no anticipation, no interference, or the counterfactual trend, factor and
noise restrictions needed for a causal interpretation.
\citep{exec23_synthdid_reuse_source}

The panel converter rejects missing values, nonbinary treatment and anything
other than one observation at every unit--time pair. It then checks
simultaneous adoption. It converts factor and date fields to character before
sorting. ISO dates may preserve calendar order, but arbitrary period labels
need an explicit ordering map. The option that leaves treated units unsorted
does not change the estimator's expectation that controls occupy the first
$N_0$ rows. Nor does the converter retain other columns as a covariate array;
that alignment remains an adapter obligation. The direct matrix entrypoint
has a different, narrower set of checks. The proposed worker must enforce
finite data, nonempty treatment/control and pre/post groups, valid noise
estimation support, and row identity on either entry route.
\citep{exec23_synthdid_reuse_source}

For an independent arithmetic reference, use
\begin{equation}
 Y=\begin{pmatrix}
 10&12&14&16\\20&22&26&28\\15&17&25&27\\18&20&28&30
 \end{pmatrix},
 \quad\omega=(3/4,1/4)',\quad\lambda=(1/4,3/4)',
 \label{exec23:panel-block-example}
\end{equation}
with two controls, two treated units and two periods on either side. The
treated-minus-weighted-control gap is four dollars in both pre-periods and
$19/2$ dollars in both post-periods. Thus (\ref{exec23:panel-contrast}) is
$11/2$ dollars; a stipulated four-dollar cost leaves $3/2$ dollars. These are
supplied weights, not weights fitted by synthdid. Their effective donor count
$1/\sum_i\omega_i^2=8/5$ describes concentration, not the number of
independent observations or guaranteed inferential degrees of freedom.
\citep{exec23_synthdid_reference}

The same example supplies a useful conversion oracle. Valid unit-sum control
weights cancel a common 100-dollar post-period shock. Replacing them by
$(1/2,1/4)'$, which sums to $3/4$, changes the contrast by 25 dollars. The
direct estimator checks supplied weight lengths but does not enforce the
simplex constraints for fixed weights. An SDID profile must check those
constraints explicitly. The synthetic-control wrapper deliberately uses
zero pre-period weights and no unit-weight intercept, so its convention must
remain a separate profile rather than being rejected by an indiscriminate
SDID validator. \citep{exec23_synthdid_reuse_source,exec23_synthdid_reference}

\subsection{Optimization and inference are separate reproducibility obligations}
\label{exec23:panel-inference}
The default unit penalty scales with estimated pre-treatment first-difference
noise and $(N_1T_1)^{1/4}$; the time penalty uses a small multiplier. The
no-covariate path uses Frank--Wolfe steps with exact line search, followed by
optional sparsification and another optimization. Its stopping condition
uses objective decrease and an iteration limit. The covariate path alternates
weight updates with a $1/t$ gradient step and tests absolute objective change.
Neither returns a separate certified error bound just because iteration
stopped. A worker should retain objective histories, limits, finite-value
checks and an independent numerical comparison when feasible. The upstream
CVXR/ECOS reference provides an explicit simplex objective; the inspected
test compares one California example at tolerance $0.03$. That source test
is useful for planning, but has not been run here and cannot establish
general parity or calibrated uncertainty.
\citep{exec23_synthdid_reuse_source}

Reproducibility also extends beyond the fitted point estimate. The attached
options retain penalties, intercepts, update flags, minimum decrease and the
main iteration limit, but omit the sparsification function and the
pre-sparsification iteration limit. Resampling refits reconstruct calls from
these saved options. When a caller changed an omitted option, subsequent
refits can return to its default. The proposed adapter must record the full
fit configuration independently and replay it deliberately; preserving the
scalar and its attached setup alone is insufficient. This is a static
call-path observation requiring native regression tests.
\citep{exec23_synthdid_reuse_source}

The package offers bootstrap, jackknife and placebo variance estimators with
different support requirements. Bootstrap and jackknife return missing
uncertainty for one treated unit. The default fixed-weight jackknife also
returns missing uncertainty when only one control has nonzero weight; it
fixes the unit and time weights rather than reoptimizing them. Placebo
requires more controls than treated units and constructs artificial treated
groups from controls. That construction needs its own scientific
comparability argument. The package vignette explicitly cautions about the
single-treated placebo illustration. A default refusal must not become a
zero standard error, and a placebo standard deviation is not by itself a
randomization-test p-value. \citep{exec23_synthdid_reuse_source}

The bootstrap retries until it has the requested count of nonmissing
estimates, with no attempt limit in the inspected loop. For two controls and
two treated units, independent sampling of four rows with replacement draws
only controls or only treated units with probability
$2(1/2)^4=1/8$. Even if every mixed draw computes successfully, the expected
attempt count per acceptable composition is $8/7$. Further numerical missing
values can increase work or prevent termination; exceptions have a distinct
failure path. This calculation quantifies composition rejection only. It is
not a native timing result or a claim that every problematic numerical fit
returns missing rather than throwing an error. Record attempts, exclusions,
accepted draws and deadlines separately.
\citep{exec23_synthdid_reuse_source,exec23_synthdid_reference}

Finally, a display method can change the inference presented to a reader.
The default summary calls bootstrap, a fast summary calls jackknife, and
the formatter uses the fast summary with a $1.96$ multiplier. A service that
requests placebo variance and later parses the formatted object could thus
report a different uncertainty convention. The output should carry the
selected variance method, seed, successful and failed draws, weights and
interval convention as explicit fields. Native output, statistical validity
and user-interface presentation need separate checks.
\citep{exec23_synthdid_reuse_source}

\subsection{A proposed reuse decision with accountable costs}
\label{exec23:panel-reuse}
For this proposal, both packages are retained as reference implementations
with conditional routes to a Python or R worker. PyFixest is attractive for
its formula interface and native numerical backends; synthdid exposes a
compact original block estimator and a separate optimization reference.
Direct adoption of either package as an unrestricted causal service is
deferred. Porting now would duplicate numerical, inferential and maintenance
obligations before a measured bottleneck establishes the benefit. A
cohort-time DID estimator may be preferable for staggered adoption, while
SC, ordinary DID and SDID must be compared under matched estimands and
appropriate assumptions rather than a shared scalar output.
\citep{exec23_pyfixest_reuse_source,exec23_synthdid_reuse_source}

PyFixest declares MIT licensing; synthdid declares GPL version two or later
or the supplied three-clause BSD alternative. These source declarations and
the retained notice texts do not clear every distributed dependency or
deployment artifact. The scientific owner must approve the target, timing,
support and uncertainty; the adapter maintainer must own identity conversion,
exact options, native oracles and upgrade regressions; the operator must
account for budgets, failure recovery and distribution requirements.
\citep{exec23_pyfixest_reuse_source,exec23_synthdid_reuse_source}

The first useful public or synthetic trial should therefore compare the
complete request: eligible sample construction, estimation, uncertainty,
failed requests, explanation and reproducible cost accounting. Measure
preparation time, model fitting, covariance or resampling, rejected draws,
worker recovery and maintenance separately. The exact examples here are
input/output oracles for such a trial, not evidence that it has already
succeeded. A successful source-based reuse decision can justify the
prototype experiment while leaving native qualification, external review
and actual bank economics explicitly unresolved.
\citep{exec23_pyfixest_reference,exec23_synthdid_reference}


% SOURCE docs/release/2026-10-29/source/docs/writing-repair/2026-10-29/src/c44dd23720277503.tex CHAR 0
\section{Choosing difference-in-differences software}
\label{sec:exec24-panel-reuse}
A bank introducing a collections policy in selected branches needs a statement
about whose outcome changes, over which dates, and under which counterfactual
trend. A two-period comparison and a staggered event study are different
contracts even when both return a coefficient labelled a treatment effect.
The earlier software survey establishes the documented roles of DRDID and
did2s. Inspection of the pinned implementations now makes the proposed reuse
choice more precise: retain their original R implementations as conditional
references, with explicit sample and uncertainty contracts. Neither package is
admitted as a verified native capability by this source review.
\citep{v12_sw_drdid,v12_sw_did2s,exec24_drdid_reuse_source,exec24_did2s_reuse_source}

\subsection{The target precedes the computation}
Let $D$ indicate membership in the group treated after the baseline date,
$Y_t(d)$ denote the outcome at date $t$ under treatment $d$, and $X$ collect
pre-treatment information. The two-period target is
\[
 \tau_T=\mathbb E[Y_1(1)-Y_1(0)\mid D=1].
\]
A causal reading requires a suitable consistency and interference regime,
no anticipation in the baseline outcome, support for comparison, and a
conditional parallel-trend restriction such as
\[
 \mathbb E[Y_1(0)-Y_0(0)\mid D=1,X]
 =\mathbb E[Y_1(0)-Y_0(0)\mid D=0,X].
\]
For an observed panel, writing $\Delta Y=Y_1-Y_0$, this gives the conditional
comparison $\mathbb E[\Delta Y\mid D=1,X]-\mathbb E[\Delta Y\mid D=0,X]$
then averaged over the treated population. Repeated cross-sections require
an additional account of population composition across dates; replacing
matched changes by independent snapshots is not a storage-only change.
The package's documented ATT and separate panel/cross-section modes preserve
these distinctions; software output does not certify the identifying
restrictions.~\citep{v12_sw_drdid,exec24_drdid_reuse_source}

An authored financial example shows why the target field is mandatory.
Suppose two equally prevalent borrower strata have gains of USD 20 and USD 100,
while treatment probabilities are $4/5$ and $1/5$. The population ATE is
USD 60, but the treated-population effect is
$(4/5)20+(1/5)100=36$ dollars. With a stipulated cost of USD 40 per treated
borrower, substituting the ATE reports a margin of USD 20 where the intended
ATT margin is USD $-4$. These are exact consequences of an invented law,
not measured collections benefits or native package estimates.
\citep{exec24_panel_reuse_reference}

\subsection{DRDID: modes, retained observations and nuisance contracts}
The selected DRDID 1.3.0 source is commit
\texttt{b63f66d7e35bb6940a139a03cb2d38988555986a}.
The top-level call dispatches separately to improved and traditional panel
and repeated-cross-section routines. It returns the ATT, uncertainty,
bootstrap draws when requested, and an influence function when requested;
the improved path also returns a propensity-fit flag. A wrapper should
retain these diagnostics and the actual resolved options. In the inspected
preprocessor, an unrecognised estimation mode warns and changes to the
improved mode, and an unsupported bootstrap mode can warn and change to the
weighted mode. An API that stores only the user's requested option would
therefore misdescribe the calculation.~\citep{exec24_drdid_reuse_source}

The preprocessor requires numeric time, treatment and supplied identifiers,
checks duplicate identifier--time cells and group constancy, sorts rows,
and accepts exactly two dates and two observed treatment groups. The
implementation uses the numeric treatment indicator algebraically; the
proposed adapter must explicitly validate the documented zero/one coding
rather than infer correctness from the count of distinct values. In the
panel path it requests a balanced panel and checks invariant covariates,
group membership and weights. It then warns about missing values, removes
incomplete rows and drops collinear covariates using a QR tolerance of
$10^{-6}$. The returned before/after outcome vectors no longer carry the
original keys in this routine. The sample contract must record each exclusion
and verify the paired keys after every filtering step, including asymmetric
missing outcomes. Static inspection identifies a qualification case; it does
not establish an observed native misalignment.~\citep{exec24_drdid_reuse_source}

The improved panel routine caps fitted propensities at $1-10^{-6}$ and
applies the default $0.995$ trimming threshold to controls. It constructs
\[
 \widehat\tau_{\rm imp}
 =\frac{\mathbb P_n[ w_i s_i
       \{1-(1-D_i)/(1-\widehat p_i)\}
       (\Delta Y_i-\widehat m_i)]}{\mathbb P_n[w_iD_i]},
\]
where $\mathbb P_n$ is the sample mean, $w_i$ is the normalized input weight,
$s_i$ is the retained indicator, $\widehat p_i$ the capped propensity, and
$\widehat m_i$ the fitted control change. The traditional panel path instead
normalizes its treated and odds-weighted control components separately and
includes nuisance-estimation terms in its influence calculation. The improved
repeated-cross-section path includes four group/date regressions and extra
adjustment terms; it is not the panel expression with the identifier deleted.
\citep{exec24_drdid_reuse_source}

A supplied-nuisance fixture makes the normalization distinction visible.
Take two treated residual changes $(10,14)$, two control residual changes
$(2,4)$, unit weights, and control propensities $(1/2,999/1000)$. The control
odds are $(1,999)$. Without trimming, the improved-style expression is
$(24-2-999\cdot4)/2=-1987$, whereas the separately normalized expression is
$12-(2+999\cdot4)/1000=4001/500$. Removing the high-propensity control gives
11 and 10 respectively. These nuisance values are supplied to a mathematical
fixture; they are not claimed to satisfy either routine's fitted calibration
equations. The example neither ranks the estimators nor refutes double
robustness. It requires an adapter to preserve the actual threshold,
normalization and nuisance method instead of silently interchanging them.
\citep{exec24_drdid_reuse_source,exec24_panel_reuse_reference}

The inspected analytic paths compute standard errors from the individual
influence vector, while the multiplier helper draws one multiplier per vector
entry. There is no branch-cluster argument in these inspected signatures.
A panel subject is not automatically a branch-independent sampling unit.
A bank using branch assignment must establish the appropriate uncertainty
procedure, or decline to report qualified uncertainty. The traditional path
also warns on propensity nonconvergence and stops on missing coefficients or
singular designs. Native qualification must distinguish a returned point
estimate from an accepted fit and from accepted uncertainty.
\citep{exec24_drdid_reuse_source}

\subsection{did2s: generated outcomes retain first-stage uncertainty}
The selected did2s 1.2.1 source is commit
\texttt{31febce3581304c3f6e8a8ef91d07e20175d7c41}.
It estimates the first-stage outcome relation on untreated observations,
predicts the untreated outcome for the full data, and regresses the resulting
residual outcome on the second-stage design. Let $X$ be the first-stage
matrix, $X_0$ its untreated-row version with treated rows zeroed, and $Z$ the
second-stage matrix. In an illustrative full-rank unweighted case,
\[
 \widehat\gamma=(X_0'X_0)^{-1}X_0'Y,
 \qquad
 \widehat\beta=(Z'Z)^{-1}Z'(Y-X\widehat\gamma).
\]
Consequently the response to an outcome perturbation $dY$ contains both
$(Z'Z)^{-1}Z'dY$ and the correction
$-(Z'Z)^{-1}Z'X(X_0'X_0)^{-1}X_0'dY$.
Treating the generated baseline as fixed deletes the second term.
The inspected analytic path combines first- and second-stage influences,
then sums them within the supplied cluster before forming outer products.
Its bootstrap path resamples clusters and refits both stages.
\citep{exec24_did2s_reuse_source}

For a minimal authored illustration, estimate a constant untreated baseline
from two independent unit-variance observations and subtract it from the mean
of two independent treated observations. Conditional on a fixed baseline,
the variance is $1/2$. With the estimated independent baseline, the variance
is $1/2+1/2=1$. This exact variance identity motivates preserving the first
stage; it is not a replication of did2s inference or a universal factor-of-two
correction. Actual clustered event-study inference depends on the design,
first-stage support and sampling law.~\citep{exec24_panel_reuse_reference}

The source also exposes conversion and numerical obligations. Rows removed
by the second-stage fit are used to align the data and first-stage residuals;
weights enter through square roots in the analytic covariance. The returned
object is a \texttt{fixest} summary with a replacement covariance, so exporting
only coefficients and ordinary regression standard errors loses the method's
uncertainty. Both fitting calls suppress selected warnings/notes, and the
final summary suppresses warnings; an operator must retain explicit support,
rank and failure diagnostics rather than equate silence with qualification.
\citep{exec24_did2s_reuse_source}

The helper named \texttt{robust\_solve\_XtX} is documented as QR based, but its
body forms $X'X$, attempts a direct solve and falls back on an SVD of that
cross-product after an error. Its matrix right-hand side is interpreted by
shape: the row-count branch precedes the column-count branch. If $X$ is a
square diagonal matrix with diagonal $(2,3)$ and an identity matrix is supplied
as an already formed right-hand side, the intended solution has diagonal
$(1/4,1/9)$; interpreting that matrix as raw responses gives $(1/2,1/3)$.
This algebraic witness motivates an explicit right-hand-side representation
and numerical qualification. Reachability from an actual full-rank did2s
fit has not been established; no native failure is claimed.
\citep{exec24_did2s_reuse_source,exec24_panel_reuse_reference}

The separate event-study convenience function provides a further reason to
qualify narrow entrypoints. Its did2s branch sets the cluster variable to the
unit identifier; its neighbouring estimator branches have their own calls,
sample handling and reference periods. An array of estimates under a common
plotting interface is therefore not automatically a matched comparison.
The proposed initial profile uses an explicit did2s call, declared cluster,
first-stage support, omitted periods and event aggregation, with simultaneous
path inference separately justified.~\citep{exec24_did2s_reuse_source}

\subsection{Reuse choices, distribution and qualification}
For both packages, an original R worker is the conditional integration route:
it preserves the inspected calculations and their ecosystem. A direct Python,
Java or C++ port would also inherit responsibility for formula expansion,
rank decisions, first-stage support, covariance and failure propagation. No
measured benefit currently justifies that full port. The existing did,
fixest, PyFixest and supplied-equation references remain alternatives selected
by target and assumptions; sharing an estimator label does not establish parity.
\citep{exec24_drdid_reuse_source,exec24_did2s_reuse_source,v12_sw_drdid,v12_sw_did2s}

DRDID declares GPL-3, R at least 3.5, and dependencies including trust, BMisc,
Rcpp and fastglm, with compilation required. The GPL text's source,
modification, notice and corresponding-source provisions are retained as
review obligations for any proposed distribution. did2s declares MIT plus its
package notice, R at least 3.5.0 and fixest at least 0.13.2, with further
imports including boot, Matrix, did and other estimators. Its own package
metadata says compilation is unnecessary; that is not a statement about the
entire dependency environment. The package notice and standard MIT text are
inspected together. Neither review clears every transitive artifact or decides
the legal status of a future combined distribution. A process boundary alone
is not treated as a rights determination.~\citep{exec24_drdid_reuse_source,exec24_did2s_reuse_source}

A method owner must approve the target, counterfactual trend, interference and
sampling regime. An integration maintainer must preserve keys, resolved
options, diagnostics and exact environment; a rights reviewer must resolve
distribution and dependency obligations; and an evaluation owner must register
matched public/generated trials. Those trials include duplicate and missing
panel cells, nonbinary coding, changing weights, unsupported overlap,
nonconvergence, rank loss, unseen fixed-effect support, cluster resampling,
failed replicates and event-path covariance. Complete costs include source
and dependency maintenance, data conversion, both fitting stages, covariance,
resampling, failed attempts and review. No price, speed advantage or bank ROI
is inferred here. The decisions close the proposal's reference disposition
only; native integration, empirical performance and broad survey completion
remain open.~\citep{exec24_drdid_reuse_source,exec24_did2s_reuse_source}


% SOURCE docs/release/2026-10-29/source/docs/writing-repair/2026-10-29/src/10da7269ea4877ce.tex CHAR 0
\section{EventStudyInteract: cohort support and covariance define the rollout claim}
\label{sec:exec24-eventstudy-reuse}

A staggered branch-policy rollout creates several clocks: calendar time,
initial adoption cohort, time since adoption, and the time at which the sponsor
requests a comparison. An apparent decline in an event-study path may reflect
a change in which cohorts remain observable, rather than a weakening effect
within any branch. The community Stata package EventStudyInteract is a candidate
reference for interaction-weighted estimates in this setting. Its selected
repository revision is dated 30 October 2023; the actual command header is
version 0.1 dated 24 January 2022, while the help header and distribution metadata
carry earlier version labels. We bind these separate identities instead of
presenting the newest-looking label as a verified installed version.
\citep{exec24_eventstudy_sources}

The authentic saved Sun--Abraham original is the September 2020 arXiv version
of the work with a 2021 journal bibliographic year. Its contribution is not a
universal cure for observational rollout bias. It develops cohort-specific
comparisons and their aggregation under stated parallel-trend, no-anticipation
and sampling conditions. We inspect its estimator construction and asymptotic
assumptions separately from the selected implementation. The older citation's
limited contamination claim remains preserved; the additional inspected pages
support the more specific comparison here.
\citep{sun2021,exec24_eventstudy_original}

\subsection{Cohort effects, aggregation and a banking decision}
Let $E_i=e$ be branch $i$'s first adoption period, $\ell=t-e$ its relative time,
and $\delta_{e\ell}$ the cohort-specific effect at that relative time. For the
cohorts supported at $\ell$, write $p_{e\ell}$ for their population shares.
The interaction-weighted target and its estimator have the forms
\[
 \theta_\ell=\sum_e p_{e\ell}\delta_{e\ell},\qquad
 \widehat\theta_\ell=\sum_e\widehat p_{e\ell}\widehat\delta_{e\ell}.
\]
The original constructs an interacted regression, estimates cohort shares and
aggregates the corresponding effects. With a never-treated comparison cohort,
its proposed specification may use all periods. With only a last-treated
comparison cohort, it excludes periods at or after that cohort's adoption.
Always-treated cohorts lack a pretreatment comparison and are excluded from
that construction. Validity depends on the chosen comparison population,
parallel trends and no anticipation, not on the existence of coefficient labels.
\citep{exec24_eventstudy_original}

For an authored example, early and late cohorts have constant effects two and
ten repayment units. Shares $(3/4,1/4)$ at relative time zero give an aggregate
of four. If only the early cohort supports the next horizon, its aggregate is
two, although neither cohort's effect has changed. Fixed equal shares would
give six where both effects are identified; they cannot reveal an unobserved
late-cohort outcome at the longer horizon. The API must return horizon-specific
cohort support and weights, and distinguish a changing-population path from a
fixed-population decision. At a stipulated 100 USD per repayment unit, the
three numbers imply 400, 200 and 600 USD before costs, but answer different
population questions.
\citep{exec24_eventstudy_contract}

Suppose the early cohort adopts at time one and the last cohort at time three.
The latter is eligible as untreated comparison only before time three, so the
early cohort's nonnegative supported lags are zero and one. A lag-two claim
using the now-treated comparison is invalid for that construction. With
stipulated early and late effects five and seven and no baseline trend, the
comparison at time three is $5-7=-2$, not the early cohort effect five. The
package help assigns the exclusion to the caller; zeroing the comparison
cohort's event indicators does not remove those calendar periods.
\citep{exec24_eventstudy_sources,exec24_eventstudy_contract}

\subsection{Sample membership and returned objects must survive the interface}
The inspected command marks an analysis sample, copies relative-time
indicators with the control cohort set to zero, and estimates cohort shares
by no-intercept regressions among non-control observations. It obtains share
covariance through \texttt{avar}, then calls \texttt{reghdfe} for the interacted
outcome regression. The share calculations precede that outcome regression.
The inspected wrapper does not recompute the shares on the backend's eventual
estimation sample. If a separately qualified backend drops additional rows,
the weight and outcome populations need to be reconciled explicitly.
This is a conditional integration concern; no such native drop was observed
in this review.
\citep{exec24_eventstudy_sources}

Counts $(3,1)$ across the two authored cohorts produce weights $(3/4,1/4)$ and
effect four. If a hypothetical outcome fit retains counts $(1,1)$, its
corresponding equal-share target is six. That example makes the necessary
contract concrete: retain the marked sample, actual estimation sample,
cohort-by-horizon counts, row identifiers and all exclusion reasons. The
caller must also validate absorbing treatment, first-treatment coding,
reference period, mutually coherent event indicators, controls and support.
Permitting a syntactic option does not prove its causal or numerical validity.
\citep{exec24_eventstudy_contract}

The correct exported estimate and covariance are \texttt{e(b\_iw)} and
\texttt{e(V\_iw)}. The underlying \texttt{e(b)} and \texttt{e(V)} still belong
to the interacted regression. The separately returned \texttt{V\_interact}
contains extracted coefficient variances, not the complete interacted
covariance matrix. An adapter must preserve the full covariance it actually
uses, row/column ordering, event-time labels and cohort shares; substituting
diagonal summaries changes subsequent inference.
\citep{exec24_eventstudy_sources}

\subsection{Uncertainty: coefficient covariance, estimated shares and joint claims}
For a fixed horizon, the first-order differential of the aggregate is
$d\theta=p^{\mathsf T}d\delta+\delta^{\mathsf T}dp$. Its general variance
contains the coefficient term, the share term and their cross-covariance.
Under a justified zero cross-covariance condition this reduces to
\[
 \operatorname{Var}(\widehat\theta)\simeq
 p^{\mathsf T}V_\delta p+\delta^{\mathsf T}V_p\delta.
\]
The source computes a transformation of the full interacted covariance and
adds a cohort-share covariance contribution. The original derives its limit
under explicit saturated-regression assumptions, including independent
cross-sectional units, finite fourth moments, rank conditions and within-unit
time dependence. These conditions do not assert validity for arbitrary
cross-branch dependence or selected sparse event cells.
\citep{exec24_eventstudy_sources,exec24_eventstudy_original}

With authored weights $(3/4,1/4)$ and coefficient covariance
\[
 V_\delta=\begin{pmatrix}4&1\\1&9\end{pmatrix},
\]
the full coefficient contribution is $51/16$, whereas discarding its
covariance gives $45/16$. If the early-cohort share has variance $0.01$, the
share gradient is $2-10=-8$, adding $16/25$ under this example's stipulated
zero cross-covariance. The resulting variance is $1531/400$. These are exact
propagation calculations, not measurements of package coverage.
\citep{exec24_eventstudy_contract}

The share-covariance call in the selected wrapper hardcodes a robust option;
the requested covariance option is passed to the outcome regression. The
source itself flags its normalization for unbalanced panels as unfinished.
The help also separates the balanced, no-covariate result it cites from the
broader data accepted by the interface. Thus a multiway or serial covariance
request must be audited through every component and dependency, not inferred
from the outcome-regression option. A future native benchmark needs explicit
balanced and unbalanced designs, singleton and missing-cell cases, requested
weights, dependence patterns and covariance alignment.
\citep{exec24_eventstudy_sources}

Pointwise intervals answer one relative-time question. They do not establish
simultaneous coverage of a ten-period path. If ten intervals had independent
95 percent coverage, their joint coverage would be $0.95^{10}\simeq0.5987$.
Independence is stipulated only for this illustration; actual event-study
estimates are correlated. A portfolio decision that scans the largest effect,
or accepts a rollout only if every horizon improves, needs an appropriate
joint procedure and its own qualification. Trend sensitivity and anticipation
also remain scientific questions after a covariance calculation succeeds.
\citep{exec24_eventstudy_original,exec24_eventstudy_contract}

\subsection{A permissive wrapper still requires an admissible execution path}
The selected repository includes an MIT license, and its distribution metadata
also declares MIT. This resolves the earlier inspection's unestablished
wrapper-license field for this pinned snapshot. It does not clear the Stata
runtime, all dependency distributions or a particular deployment. The dependency
script contains unpinned installation requests for \texttt{ftools}, \texttt{avar}
and \texttt{reghdfe}; we inspect these as text without executing them. Runtime
entitlement, dependency versions and rights, compilation, environment and future
maintenance remain qualification work. Help and syntax also differ on admitted
weight classes, so an adapter must declare and test a selected profile rather
than equating parser acceptance with supported inference.
\citep{exec24_eventstudy_sources}

Our disposition is reference, with a conditional original Stata worker after
rights, environment and scientific qualification. A standalone implementation
of cohort aggregation is useful as an independent oracle but is not a port of
the full regression and inference stack. A separately qualified R or Python
event-study route may be an alternative only after matching controls, support,
weights, covariance and target. Whole-package porting is deferred until a
specific business need and comparative cost justify it. The scientific owner
sets the rollout target and assumptions; the adapter owner preserves samples
and objects; the platform/rights owner establishes the executable environment;
and the evaluation owner preregisters public positive and adverse fixtures.
\citep{exec24_eventstudy_contract}

Illustrative cost of 1,800 USD over 90 requests, 75 conformant results and 60
useful results gives 20, 24 and 30 USD per denominator. These are transparent
planning inputs, not observed prices. Generated rollout panels can verify
cohort arithmetic, sample reconciliation, covariance and refusal behavior before
expensive bank experiments. They cannot establish the bank's untreated trend,
absence of anticipation or missing long-horizon outcomes. The reuse route is
known; target truth, native qualification and complete rights remain conditions.
\citep{exec24_eventstudy_contract}


% SOURCE docs/release/2026-10-29/source/docs/writing-repair/2026-10-29/src/e2865a7ddf66f706.tex CHAR 0
\section{R panel software: compare the target before comparing the estimator}
\label{sec:exec24-panels-r}

A branch rollout can ask about a particular adoption cohort, a calendar
quarter, the first year after adoption, or the average benefit across all
exposed branch-periods. These are different targets even before nuisance
models or standard errors enter. The R comparison therefore connects the
selected \texttt{did} implementation to the already inspected
\texttt{DRDID}, \texttt{did2s}, \texttt{fixest}, \texttt{HonestDiD},
\texttt{synthdid} and \texttt{gsynth}/\texttt{fect} paths. These are concrete
alternatives with different assumptions and interfaces, not interchangeable
labels for an automatic causal answer.
\citep{exec24_panelr_contract}

The selected \texttt{did} revision is dated 28 August 2026 and declares
version 2.5.1.902. Its previously retained \texttt{att\_gt} file authenticates
to the same complete cutoff-pinned tree. We additionally inspect aggregation,
preprocessing, cell estimation and multiplier-bootstrap delegates. The
revision's declared GPL-3 license, R requirement and dependencies are
observations about these bytes, not proof of an installed distribution,
commercial support or qualification of every imported estimator.
\citep{exec24_panelr_did}

\subsection{Four aggregate answers from the same cohort-time effects}
Write $\tau_{gt}$ for the effect in cohort $g$ at calendar time $t$ and $p_g$
for its population share. A treated-cell average weights every supported
post-adoption cell by its cohort share. A group average first averages within
each cohort over selected post periods and then weights cohorts. A dynamic
average first aggregates cohorts at each relative time and then averages the
selected nonnegative relative times. A calendar average first aggregates
cohorts at each calendar time and then averages those calendar effects.
The inspected implementation contains separate branches for all four targets;
its public aggregation function defaults to the group target. Where weights
are estimated, their influence contributions accompany the effect estimates.
\citep{exec24_panelr_did}

In an authored panel, cohort one has effects two and four at times one and
two; cohort two has effect ten at time two. Their population shares are
$3/4$ and $1/4$. The simple treated-cell target is
\[
 \frac{(3/4)2+(3/4)4+(1/4)10}{3/4+3/4+1/4}=4.
\]
The group target is $(3/4)(2+4)/2+(1/4)10=19/4$.
Dynamic effects at horizons zero and one are both four, giving an overall
value four. Calendar effects are two and $11/2$, averaging to $15/4$.
At a stipulated 100 USD per outcome unit, these planning quantities are 400,
475, 400 and 375 USD per respective target unit. The differing denominators
must be explicit before a sponsor treats them as comparable returns.
\citep{exec24_panelr_contract}

Balancing on one post-adoption lag retains only the early cohort in this
example, giving the path $(2,4)$ and mean three. That is a supported,
restricted population, not an estimate of the late cohort's absent longer
history. In the selected source, event balancing affects dynamic aggregation;
simple and group aggregation warn that they ignore that option. Calendar
aggregation warns that it ignores the event-window and balancing options.
The CLI/MCP result must report the resolved request and target population,
not merely repeat options the caller supplied.
\citep{exec24_panelr_did,exec24_panelr_contract}

\subsection{Missing cells and changing weights are population choices}
The selected cell wrapper catches a two-by-two estimation error, warns, and
leaves that cell unavailable. Aggregation rejects missing effects by default;
when removal is requested it filters cells and, in the group branch, removes
groups without retained post-treatment effects in the requested window. The
result therefore describes what remains. It must not be relabeled the original
portfolio without an argument for the missing contributions.
\citep{exec24_panelr_did}

If the late cohort's time-two cell is unavailable in the authored example,
the simple and group averages among remaining cells are both three. The
stipulated complete-data targets were four and $19/4$. Those complete values
are known only because this is a constructed example; a production analysis
must not impute the missing effect by copying them. Requested, eligible,
attempted, estimated and omitted cohort-time cells belong in separate counts.
\citep{exec24_panelr_contract}

The selected \texttt{att\_gt} documentation and cell code also distinguish
sampling-weight clocks. A balanced panel's default two-period comparison
uses earlier-period weights; explicitly varying weights switches to a stacked
repeated-cross-section route while holding covariates at the earlier period.
Base-period or first-period anchoring on an unbalanced panel can remove units
absent at that anchor. For effects $(0,10)$, weights $(9,1)$ give one while
weights $(1,9)$ give nine. If only the first unit is present at an anchor,
the retained target is zero. Neither a constant variable name nor a successful
fit establishes that the population stayed fixed.
\citep{exec24_panelr_did,exec24_panelr_contract}

The stacked route folds pre- and post-period influence contributions back to
the same unit with a factor of one half. Authored stacked contributions
$(-2,2,-4,4)$ become $(-3,3)$; summing without the half gives $(-6,6)$,
four times the variance and twice the standard error under the same variance
calculation. The selected source includes the half. This is a positive
semantic check and a proposed conversion oracle, not a reported defect.
\citep{exec24_panelr_did,exec24_panelr_contract}

\subsection{Controls, anticipation and uncertainty are resolved outputs}
The preprocessing path validates unit-time uniqueness and constant first
adoption within a unit, and treats the intervention as irreversible. A
first-adoption field cannot faithfully encode repeated switches between
policies. The not-yet-treated comparison requires a cohort to adopt strictly
after the comparison period plus anticipation. With time two and one period
of anticipation, a cohort adopting at three is not eligible; a cohort adopting
at four is. When no never-treated group exists, the selected preprocessing
can use the last cohort as comparison and removes unsupported later periods.
These transformations and support losses must appear in the result.
\citep{exec24_panelr_did,exec24_panelr_contract}

Disabling influence-function computation returns point estimates only and
prevents the documented aggregation and inference paths. With inference
retained, the selected multiplier bootstrap permits the individual sampling
unit and at most one additional coarser clustering variable. Its clustered
path aggregates influence contributions to cluster sums, with corresponding
normalization. It also removes degenerate dimensions and nonfinite bootstrap
maxima before forming its simultaneous critical value. A future worker must
return the actual dimensions and denominator used, and qualify sparse or
unstable cases rather than calling every finite critical value calibrated.
\citep{exec24_panelr_did}

An aggregation request cannot introduce a cluster variable absent from the
stored group-time object. The inspected code warns and removes unsupported
extra clustering. Thus a request for branch clustering on an object retaining
only region clustering cannot be reported as fulfilled. The adapter should
refuse or explicitly mark the result nonconformant and request recomputation
from the appropriate first-stage object. A simultaneous-band request has a
similar contract: invalid or unexpectedly small bootstrap critical values
trigger a pointwise fallback, and the returned status must say pointwise.
A large critical value triggers a warning rather than proving useful precision.
\citep{exec24_panelr_did,exec24_panelr_contract}

\subsection{Choose among complementary R routes with their assumptions intact}
The previously inspected \texttt{DRDID} dispatcher separates panel and
repeated-cross-section designs and improved from traditional nuisance routes.
Its selected multiplier helper does not itself accept arbitrary cluster
structure. \texttt{did} can organize such two-period comparisons, aggregation
and a separate influence-based covariance layer, but that composition needs
a qualified dependency version and identity alignment. The \texttt{did2s}
route fits an untreated first stage and a full-sample second stage, with
its own correction for generated outcomes and clustered inference. Ordinary
second-stage regression uncertainty is not automatically that correction.
\citep{exec24_drdid_reuse_source,exec24_did2s_reuse_source}

\texttt{fixest} supplies a broad projection and covariance engine. Its
selected defaults, singleton/perfect-fit removals, small-sample scaling,
covariance repair and event-cohort mapping require explicit decisions.
Fixed effects alone do not select clustered errors. Its event-study
aggregation uses full covariance, which a summary-table adapter must preserve.
\texttt{HonestDiD} adds sensitivity to specified trend deviations, with
relative-magnitude and smoothness restrictions of different meanings. It
requires correctly ordered coefficients and covariance and honest numerical
status; it cannot verify that the imposed restriction is substantively true.
\citep{exec20_fixest_source,exec22_honestdid_engine}

The original \texttt{synthdid} source reviewed here targets a simultaneous
adoption block. It should not inherit arbitrary staggered-adoption support
from another implementation's name. The selected \texttt{gsynth} compatibility
wrapper delegates to \texttt{fect}; the earlier inspection found that its
weight argument was not forwarded even though its documentation described a
weighted effect. The inspected \texttt{fect} engine distinguishes fitting
weights, aggregation weights, factor assumptions, reversal and inference
gates. A stable numerical factor fit does not establish counterfactual
stability or solve an unidentified missing trajectory.
\citep{exec23_synthdid_reuse_source,exec21_gsynth_wrapper,exec21_gsynth_engine}

For a branch rollout, the comparison sequence is therefore to specify target,
clock, treatment history, controls, support, sampling population and permissible
trend restrictions; select a concrete route; retain actual samples and
uncertainty; and compare against independent public or generated oracles.
The qualified scope may be a restricted cohort or a refusal. A software census
also needs explicit tails: alternative panel estimators, reversal designs,
spatial/network dependence and other R packages remain to be screened. The
selected seven families do not establish exhaustive coverage.
\citep{exec24_panelr_contract}

\subsection{Public qualification and costs before expensive bank experiments}
Known-law panels can compare cohort, event and calendar targets; anticipation,
missing cells, sampling-weight anchors, influence-vector ordering, cluster
membership and fallback behavior. A preregistered benchmark must preserve
requested and delivered targets, complete failures, dependence assumptions,
coverage uncertainty and compute costs. Illustrative total cost of 3,000 USD
across 150 requests, 125 conformant outputs and 100 useful outputs gives 20,
24 and 30 USD per denominator. These are planning values, not observed prices.
\citep{exec24_panelr_contract}

The scientific owner approves causal and sampling restrictions; the adapter
owner preserves transformations, objects and refusal semantics; the
platform/rights owner pins the installed dependencies and distribution rights;
and the evaluation owner supplies independent oracles and complete accounting.
No selected source inspection alone grants a native integration. Public
experiments can establish merit before bank experiments, but cannot manufacture
untreated trends, missing counterfactuals or universal validity under arbitrary
dependence. Coverage remains incomplete while search tails and those native
qualification contracts remain open.
\citep{exec24_panelr_contract}


% SOURCE docs/release/2026-10-29/source/docs/writing-repair/2026-10-29/src/95d9c9bafe56fb21.tex CHAR 0
\section{Python panels: preserve the estimand and the sampling unit across translations}
\label{sec:exec24-panels-python}

A bank comparing branch rollout effects across Python implementations needs
more than matching coefficient names. The analysis must preserve the treatment
clock, cohort support, population weights, influence-vector ordering and the
sampling unit used for uncertainty. We compare the selected \texttt{differences}
and \texttt{csdid} sources with the previously reviewed \texttt{PyFixest}
path. The inherited \texttt{linearmodels} IV review does not qualify its panel
engine. The separate direct panel review in
Section~\ref{sec:exec24-linearmodels-panel} supplies that source comparison,
including covariance clocks and numerical termination limits. Its qualification
must remain distinct from the IV reference when selecting an integration.
\citep{exec24_panelpy_contract,exec23_pyfixest_reuse_source,exec22_linearmodels_reuse}

The selected \texttt{differences} commit is dated 12 April 2026 and uses a
dynamic package version, Python at least 3.10 and declared GPLv3 terms.
The selected \texttt{csdid} commit is dated 2 July 2026 and declares version
0.4.2. Its actual license file and setup license field say MIT, while a setup
classifier says Apache. We preserve that discrepancy and do not infer a
second grant from the classifier. Neither commit is evidence of an installed
wheel, tested environment or cleared dependency graph. Twenty-one retained
files authenticate to the two pinned repository trees; technical conclusions
below refer to identified paths within those files.
\citep{exec24_panelpy_differences,exec24_panelpy_csdid}

\subsection{Matching the target precedes a parity comparison}
The \texttt{differences} implementation distinguishes simple, cohort,
event and calendar-time aggregations. Cohort effects average retained post
periods, the overall cohort target uses population shares, and event and
calendar summaries average their retained event or time effects. Its influence
aggregation includes contributions from estimated weights. The selected
\texttt{csdid} aggregator distinguishes simple, group, dynamic and calendar
targets, with separate support and missing-cell rules. Similar names therefore
invite a target-level comparison, not an assumption of identical defaults.
\citep{exec24_panelpy_differences,exec24_panelpy_csdid}

Reuse the authored three-cell rollout: early-cohort effects are two and four,
the late-cohort effect is ten, and cohort shares are $3/4$ and $1/4$.
The simple target is four, the cohort target is $19/4$, the event summary is
four, and the calendar summary is $15/4$. These are independent arithmetic
oracles for deliberately matched populations. They are not observed outputs
from either package. A parity trial must first freeze the retained cells,
normalization, weighting clock and treatment history, then compare both
estimates and influence contributions.
\citep{exec24_panelpy_contract}

Control encoding also changes meaning. The inspected \texttt{differences}
cell selector uses missing cohort values for never-treated units and requires
a not-yet-treated cohort to lie strictly beyond the later comparison period
plus anticipation. Its validator explicitly rejects one common zero-coded
never-treated input convention. In \texttt{csdid}, absence of a never-treated
group can instead cause the last cohort to be used as control and late
periods to be removed, with a warning. A last-cohort control is valid only
over the retained untreated support and under the stated causal restrictions.
The API must retain the resolved control population and information clock.
\citep{exec24_panelpy_differences,exec24_panelpy_csdid}

At comparison time two with anticipation one, cohort three fails the strict
not-yet-treated test and cohort four passes it. The arithmetic establishes
the selection rule; it cannot verify that anticipation is actually one.
Parallel trends, no interference and the relevant treatment history remain
scientific restrictions, with sensitivity or refusal when unjustified.
\citep{exec24_panelpy_contract}

\subsection{Unequal clusters expose a consequential normalization choice}
Let $\phi_i$ be a scalar influence contribution for a unit-average target,
$n$ the number of units, and $S_c=\sum_{i\in c}\phi_i$ the sum in cluster
$c$. Under independent unit-variance cluster multipliers $\xi_c$, the
authored perturbation for that target is
\[ Z_u=\frac{1}{n}\sum_{c=1}^{C}\xi_c S_c. \]
Replacing sums with cluster means $\bar\phi_c=S_c/n_c$ and averaging over
clusters instead gives
\[ Z_c=\frac{1}{C}\sum_{c=1}^{C}\xi_c\bar\phi_c. \]
These coincide for equal cluster sizes but need not for unequal sizes.
They are different conditional perturbation laws; a package adapter must
justify the law against the target and upstream influence normalization.
\citep{exec24_panelpy_contract}

For influences $(3,-1,-1,-1)$ and clusters of sizes one and three, the sums
are $(3,-3)$ and the means are $(3,-1)$. Enumerating the four sign vectors
gives $\operatorname{Var}(Z_u)=9/8$ and
$\operatorname{Var}(Z_c)=5/2$. Treating all four units as independent gives
$3/4$. This witness demonstrates why a branch-cluster request cannot be
silently reduced to an independent-unit or equal-cluster calculation.
These exact variances are not claims about either package's returned
interquartile-range standard error or empirical confidence coverage.
\citep{exec24_panelpy_contract}

In the selected \texttt{differences} bootstrap, the helper really does divide
group sums by group counts. The outer function multiplies its resampled
cluster-mean average by $\sqrt C$ and divides the resulting scale estimate
by $\sqrt C$. This establishes the selected normalization, while full
scientific qualification still requires the relevant upstream influence
definition and sampling law. The analytic inference branch calls the
unit influence cross-product helper without using the extra cluster-group
argument. Thus zero bootstrap iterations do not, at this boundary, fulfill
an extra-cluster request merely because that argument was supplied.
\citep{exec24_panelpy_differences}

The inspected \texttt{csdid} bootstrap instead forms cluster sums, generates
multiplier averages on those sums, and returns the scale multiplied by
$\sqrt C/n$. Its analytical aggregation helper uses
$\sqrt{\sum_c S_c^2}/n$. This is positive source evidence for a different
normalization. Correct labels, row alignment, cluster invariance, sufficiently
informative clusters and the assumptions behind the influence expansion
remain separate obligations; formula resemblance is not a coverage theorem.
\citep{exec24_panelpy_csdid}

\subsection{Requested uncertainty and delivered uncertainty can diverge}
The \texttt{csdid} aggregator turns a list of cluster names into its first
element before passing the resolved option onward. Consequently a request
ordered as unit then region can lose the extra region specification at this
boundary. The lower bootstrap helper rejects multiple extra names, but that
guard cannot recover names already removed. Missing retained cluster data
also causes a warning and a fallback. The proposal therefore requires
explicit requested and delivered cluster fields, with nonconformance or
recomputation when they differ.
\citep{exec24_panelpy_csdid,exec24_panelpy_contract}

For calendar bands, invalid or unusually small bootstrap critical values
can trigger a pointwise normal fallback and set the \texttt{cband} flag
false. The selected formatter chooses the word ``Simult.'' from the
\texttt{bstrap} flag instead. If bootstrap remains true after that fallback,
the displayed label and delivered uncertainty type can disagree. This is a
static path finding, not an observed native failure. A CLI/MCP adapter must
derive its uncertainty type from resolved inference and critical value,
preserve warnings, and reject an unmet simultaneous-band requirement.
\citep{exec24_panelpy_csdid,exec24_panelpy_contract}

The \texttt{differences} bootstrap uses a truth test for the supplied seed;
zero therefore takes the global-random-state branch. Its no-progress branch
constructs an array with dimensions $(1,2,\ldots,J)$, where $J$ is the
resolved number of jobs, before zipping that array with worker chunks.
For two jobs, the array has one outer element rather than two. With ten
requested iterations split into two chunks of five, the visible zip can
schedule only the first chunk, and a later consumer may reject the
array-valued progress flag. These are inspection and shape-arithmetic
findings; no downloaded joblib, NumPy or package code was run.
\citep{exec24_panelpy_differences,exec24_panelpy_contract}

Its selected critical-value helper divides by an interquartile scale without
a local finite-maxima filter. Conversely, the inspected \texttt{csdid}
helper has explicit degenerate-dimension and empty-finite-statistic branches.
The distinction informs adversarial public tests: all-zero influences,
single or sparse resamples, seed zero, unequal clusters and incomplete
worker output need declared statuses and complete request denominators.
A finite answer or successful process exit does not certify calibration.
\citep{exec24_panelpy_differences,exec24_panelpy_csdid}

\subsection{A bounded integration choice and its public qualification route}
The earlier \texttt{PyFixest} inspection supplies another concrete route,
including distinct event-study branches, two-stage covariance and solver
failure contracts. It also preserves source limitations around saturated
covariates, cohort-time mapping and full native solver inspection.
The separate \texttt{linearmodels} panel study in
Section~\ref{sec:exec24-linearmodels-panel} examines within transformation,
row selection, covariance normalization and solver diagnostics. It provides
a conditional fixed-effects alternative, with native qualification still open;
its findings do not follow from the earlier IV review.
\citep{exec23_pyfixest_reuse_source,exec22_linearmodels_reuse}

For the proposal, retain both newly inspected implementations as comparison
references. Do not yet admit either as a qualified native adapter or port its
code. The adapter owner should specify target, cohort/control encoding,
weights, input mutation, missing-cell accounting, cluster identity, seeds,
actual resample counts, warnings and uncertainty type. The scientific owner
must predeclare the sampling law and independent oracle. The platform and
rights owner must pin the complete environment and resolve distribution
terms, including borrowed helpers and dependencies. These are accountable
work roles, not assigned bank personnel.
\citep{exec24_panelpy_contract,exec24_panelpy_differences,exec24_panelpy_csdid}

A public qualification study should begin with known-law panels and the
exact witnesses above, then add prespecified heterogeneity, anticipation,
unbalanced observations, missing cells, unequal clusters and fallback cases.
Native trials require a separately approved execution environment; this
review executes only authored arithmetic. Matched failures and refusals
belong in the accounting. A stipulated 3,600 USD across 180 requests,
150 conformant outputs and 120 useful outputs gives 20, 24 and 30 USD
per respective denominator. These are planning costs, not vendor prices
or observed bank returns.
\citep{exec24_panelpy_contract}

Coverage remains bounded. The direct \texttt{linearmodels} panel comparison
that follows closes a specific source-inspection question. Additional
synthetic-control and panel packages, nonabsorbing or multiple treatments,
complete preprocessing delegates, nuisance estimators, compiled dependencies
and broader dependence regimes remain search or qualification tasks. More package inspection can resolve implementation
contracts; it cannot manufacture unobserved counterfactuals or validate
parallel trends. The selected Python coverage remains inconclusive, and
native, coverage, visual and final qualification contracts remain open.
\citep{exec24_panelpy_contract}


% SOURCE docs/release/2026-10-29/source/docs/writing-repair/2026-10-29/src/02ebd1c023b539a0.tex CHAR 0
\section{linearmodels panels: fixed effects, covariance clocks and output meaning}
\label{sec:exec24-linearmodels-panel}

A bank comparing regional responses over time needs separate decisions about
fixed effects, causal interpretation and sampling uncertainty. The earlier
linear-IV reference decision does not qualify the panel engine in the same
distribution. Direct inspection of the pinned \texttt{linearmodels} panel
source supplies a more specific comparison: transformations, row selection,
absorbed regressors, covariance construction and reported reference
distributions all affect the proposed result contract.
\citep{exec22_linearmodels_reuse,exec24_lmpanel_source,exec24_lmpanel_reference}

Eight selected panel and helper originals authenticate to repository commit
\texttt{13390eee\allowbreak{}49ea06fc\allowbreak{}42f83125\allowbreak{}5ff2163f\allowbreak{}2ad13e8b}, selected on
7 October 2026 before the client cutoff. Inherited original license and
build metadata retain their exact identities. The package version is dynamic;
equivalence to a particular release wheel or installed binary is unverified.
The inspected top-level NCSA terms do not resolve every dependency's rights.
This review reads original source and executes authored rational arithmetic;
it does not import the downloaded package or claim native calibration.
\citep{exec22_linearmodels_reuse,exec24_lmpanel_source}

\subsection{Within transformation and causal meaning are different steps}
For the simple additive panel model
\[
 Y_{it}=\alpha_i+X_{it}\beta+u_{it},
\]
demeaning within unit removes $\alpha_i$ algebraically. It does not remove
time-varying confounding or establish a policy interpretation for $\beta$.
Under an appropriate strict-exogeneity condition and a common linear
coefficient, within least squares has a specified statistical target.
Heterogeneous treatment responses, anticipation, feedback from prior
outcomes or unobserved common shocks require a different argument. The
proposed product must retain these assumptions rather than convert a
regression coefficient automatically into an intervention effect.
\citep{exec24_lmpanel_reference}

An authored two-region example makes the transformation auditable. Region A
has regressor values $(0,1)$ and revenue $(10,13)$ thousand USD; region B has
$(0,2)$ and $(20,22)$. With entity effects, no time effects and no separate
intercept, the demeaned vectors are
\[
 \widetilde x=(-1/2,1/2,-1,1)',\qquad
 \widetilde y=(-3/2,3/2,-1,1)'.
\]
The cross product $\widetilde x'\widetilde x=5/2$ gives
$\widehat\beta=7/5$ thousand USD per regressor unit. Residuals are
$(-4/5,4/5,2/5,-2/5)$ and scalar score contributions are
$(2/5,2/5,-2/5,-2/5)$. This is a reproducible regression calculation, not
native output, a coverage study or identified banking treatment evidence.
\citep{exec24_lmpanel_reference}

The inspected input path drops rows with missing outcome, regressors or
weights, requires strictly positive retained weights and normalizes them
by their mean. Weight shapes may represent time, entity or full-panel
information, with an explicit ambiguity error when dimensions coincide.
These are not interchangeable exposure and sampling weights. The adapter
must preserve original indices, weight roles, removed rows and final
analysis dimensions before reporting a result.
\citep{exec24_lmpanel_source}

The default \texttt{singletons=True} returns before the singleton-dropping
path; setting it false activates two-core removal. This behavior must be
read from the selected implementation rather than guessed from the option
name. Recursive removal can exceed a one-pass deletion. In an authored
entity-time graph with edges $(A,1),(A,2),(B,2)$, removing the first singleton
edges leaves $(A,2)$, which is then removed too. Absorbed-regressor checking
is another distinct step. With \texttt{drop\_absorbed=True}, the fit can
change the stored regressor matrix and coefficient names. A requested
treatment variable that disappears must be reported as unidentified in
that specification, not as a zero effect.
\citep{exec24_lmpanel_source,exec24_lmpanel_reference}

\subsection{Covariance scaling and reference degrees of freedom need separate fields}
The inspected PanelOLS fit defaults to unadjusted covariance with
\texttt{debiased=True} and \texttt{auto\_df=True}. Fixed effects therefore
do not select cluster-robust uncertainty automatically. For a single effect,
the automatic clustered adjustment can suppress counting an effect nested
in the supplied clusters; the multiple-effect route follows a different
default. A faithful adapter records the resolved effect count and scale,
the cluster identities, and any separate group-debias option.
\citep{exec24_lmpanel_source}

For the authored example, region score sums are $(4/5,-4/5)$. The unscaled
cluster sandwich is
\[
 V_{\rm cl}=\frac{(4/5)^2+(-4/5)^2}{(5/2)^2}
             =\frac{128}{625},
\]
whereas summing squared individual scores gives $64/625$. If the one-effect
nested-cluster route excludes entity effects from the covariance adjustment,
the debiased coefficient scale is $4/(4-1)=4/3$, giving $512/1875$.
The optional group scale $[2/(2-1)][(4-1)/4]=3/2$ gives $256/625$.
These values have squared slope units. They illustrate the declared
formula, not reliable small-sample inference with only two regions.
\citep{exec24_lmpanel_reference}

The regression itself still counts two entity effects and one slope, leaving
one residual degree of freedom. In the inspected results object, debiased
p-values and confidence intervals use that residual degree of freedom;
otherwise they use the normal distribution. The covariance denominator
three in this example is not the reference degree of freedom. Neither a
clustered label nor a scalar small-sample correction implies a particular
cluster-count reference distribution. These quantities must remain separate
and the validity of the chosen approximation needs a justified sampling law.
\citep{exec24_lmpanel_source,exec24_lmpanel_reference}

The nested-effect helper also iterates over supplied cluster columns.
Its single-effect result is assigned within that loop, so multi-column
nesting behavior requires explicit permutation and nesting fixtures before
relying on automatic adjustments. The proposed interface should prefer
explicitly reviewed settings when the scientific resampling unit differs
from an entity or calendar label.~\citep{exec24_lmpanel_source,exec24_lmpanel_reference}

\subsection{Two-way and temporal covariance retain finite-sample qualifications}
The two-way cluster source adds the two group covariance components and
subtracts the intersection component, optionally correcting each component
separately. It symmetrizes the final matrix without a local positive-
semidefinite projection. An abstract score witness isolates the algebra:
scores $(1,-1,-1,1)$ arranged in a two-by-two entity-time grid have zero
entity and time score sums, while intersection meat is one under division
by four. The combined scalar meat is therefore $-1$. This witness is not
claimed to be a fitted PanelOLS dataset. It shows why finite
inclusion-exclusion algebra alone cannot guarantee a positive covariance.
The results layer takes square roots of diagonal entries; a proposed worker
must test covariance validity before presenting ordinary intervals.
\citep{exec24_lmpanel_source,exec24_lmpanel_reference}

The inspected Driscoll--Kraay route sums scores by time code, sorts observed
time groups and applies a kernel to ordinal row lags. With no supplied
bandwidth it uses a count-based rule. Original elapsed calendar distances
are not the lags supplied to the kernel helper. For authored times
$(1,2,10)$ and scores $(1,1,-2)$, a lag-one weight $1/2$ produces meat
$5/3$ using neighboring observed rows. Using the same weight only for
calendar distance one gives $7/3$. Neither is universally correct; they
refer to different dependence constructions. A banking panel with missing
months must declare its time scale and preserve gaps rather than silently
interpret observed rows as equally spaced months.
\citep{exec24_lmpanel_source,exec24_lmpanel_reference}

Cluster counts, calendar counts, group sizes, bandwidth and score construction
therefore belong in the scientific output. A robust covariance does not
repair endogenous treatment, omitted trends or extrapolation outside the
retained population. Few clusters, leverage, shared shocks and irregular
sampling require targeted public experiments and sometimes additional
information or a different design.~\citep{exec24_lmpanel_reference}

\subsection{Numerical pathways need observable termination}
Dense least squares selects the SciPy rank-revealing \texttt{gelsy} route
with a dimension-and-machine-precision cutoff. The sparse effects path
calls LSMR with specified tolerances but extracts only its coefficient
component in the inspected local calls. This does not mean SciPy lacks
termination diagnostics; it means those diagnostics are not retained by
that path. Iterated general demeaning uses a scaled change tolerance and
has no explicit local iteration cap. A bounded worker needs resource,
termination and residual checks around these operations, together with
comparison against explicit dummy-variable or exact small references.
\citep{exec24_lmpanel_source}

An adapter must distinguish algebraic rank, numerical rank, absorbed
variables and scientific identification. Tightening a tolerance cannot
create missing treatment variation, and dropping columns changes the
question. Comparisons across sparse, dense, weighted and low-memory routes
should retain the resolved design, objective, coefficient mapping and
transformed residuals, not only compare printed coefficients.
\citep{exec24_lmpanel_reference}

\subsection{Scope of reuse and a public qualification route}
The direct panel review closes a specific source-inspection gap in the
earlier Python comparison. It does not close the broader search for
synthetic-control, reversal and multiple-event packages, qualify every
linearmodels panel estimator, or change the original inconclusive coverage
receipt into a passing exhaustive survey. The inherited PyFixest,
\texttt{differences} and \texttt{csdid} comparisons remain substantive
alternatives with distinct estimands and uncertainty contracts. A common
CLI/MCP schema must retain these distinctions.
\citep{exec24_lmpanel_reference}

For a bounded fixed-effects prototype, a Python worker is a reasonable
candidate after environment, rights and public trials are qualified.
The request should name the causal target, formula, effects, weights,
index clock, missingness policy, singleton and absorption settings,
covariance law and numerical limits. The result should expose removed rows
and columns, retained estimand, coefficients, covariance components,
effect and residual degrees of freedom, group corrections, bandwidth,
solver evidence, refusals and costs. Porting the full stack is deferred
until there is a measured need; a small rational oracle remains useful
for transformation and schema tests.~\citep{exec24_lmpanel_reference}

A scientific owner must justify the assignment and dependence laws; the
adapter owner must preserve transformations and result identities; the
platform and rights owner must qualify the distribution; the evaluation
owner must freeze tasks and acceptance rules. An illustrative 3,600 USD
workflow over 180 requests, 150 finite results and 120 accepted results
costs 20, 24 and 30 USD on the respective denominators. These are authored
prices, not measured economics. Native trials must count preprocessing,
every failed or retried fit, validation, storage and maintenance.
\citep{exec24_lmpanel_reference}

The concrete route is to freeze known-law public panels covering unequal
clusters, nested and crossed effects, absorption, iterative singleton
removal, weighted transformations, irregular calendars, sparse solver
termination and adverse covariance. Compare all-request coverage,
availability, interval width and cost against declared alternatives.
Unknown treatment truth and arbitrary dependence retain no universal
software remedy. Execution and independent validation remain outstanding.
\citep{exec24_lmpanel_reference}


% SOURCE docs/release/2026-10-29/source/docs/writing-repair/2026-10-29/src/e7a2c1165317e35d.tex CHAR 0

\section{Java regression components and the missing panel contract}
\label{sec:exec25-panels-java}
A bank standardizing its services on Java still needs to decide which causal
panel problem each numerical component solves. A branch identifier, a time
column and an ordinary regression interface do not specify a staggered-policy
estimand. We compare two dated Java regression implementations, Haifeng Li's
Statistical Machine Intelligence and Learning Engine (Smile) and Apache
Commons Math, with the distinct econometrics-oriented iMetrica candidate.
Smile here is not BayesFusion's SMILE Bayesian-network engine. The selected
source commits precede the research cutoff; they are source profiles, not
claims that a current binary release or a bank deployment has been tested.
\citep{exec25_paneljava_smile,exec25_paneljava_commons,exec25_paneljava_imetrica}

\subsection{Useful numerical machinery, different scientific obligations}
The selected Smile source declares version 6.3.0. Its ordinary least-squares
path expands a formula into a design matrix and rejects a system with no more
rows than columns. It selects QR or singular-value decomposition (SVD), and
the QR path catches a runtime exception before attempting SVD. If standard
errors or recursive updating are requested, the subsequent path still constructs
an inverse through a Cholesky route. Thus an SVD coefficient fallback should
not be advertised as proof that every rank-deficient uncertainty calculation
will succeed. This is a source-path qualification question, not an observed
execution failure. The reported coefficient standard error is the residual
scale times the square root of a diagonal inverse-Gram entry. The linear-model
constructor calculates the residual scale using residual sum of squares divided
by $n-p$. The inspected path does not consume cluster identifiers or calculate
a cluster sandwich.~\citep{exec25_paneljava_smile}

The Apache profile is a 4.0-SNAPSHOT source tree, not a stable 4.0 release
claim. Its legacy-module descriptor explicitly describes code that is unsupported
or not yet modularized. The inspected OLS implementation solves by QR, with
an intercept included by default. Its public parameter-variance method returns
the inverse Gram matrix; the standard-error method separately multiplies its
diagonal by the estimated residual variance before taking square roots. This
distinction matters to an adapter that expects a fully scaled covariance matrix.
Selected upstream tests cover numerical examples and residual or standard-error
comparisons, but their presence is not a locally observed test pass.
\citep{exec25_paneljava_commons}

For full-rank design matrix $X$, the conventional formulas can be expressed
as $B=(X^\mathsf{T}X)^{-1}$ and $\widehat V=\widehat\sigma^2 B$, with
$\widehat\sigma^2=\mathrm{RSS}/(n-p)$. In an authored one-parameter
example, $X^\mathsf{T}X=2$ and $\widehat\sigma^2=9$ give $B=1/2$ and
$\widehat V=9/2$. Treating the unscaled return as covariance loses a factor
of nine. Conversely, multiplying an already scaled matrix again is also wrong.
The integration contract must specify the returned object's scale, residual
degrees of freedom, intercept convention and rank handling, and test them
against an independent oracle. Matching method names is insufficient.
\citep{exec25_paneljava_reference}

The inspected iMetrica README describes interactive multivariate time-series
modeling, real-time filtering and signal extraction, with computational
libraries written in Fortran or C and a Java interface. It lists native
dependencies and a Linux-oriented build. This makes it relevant to an
econometrics survey, but the inspected README and root license do not
establish a staggered-adoption or cluster-inference implementation. We defer
it as a panel delegate while retaining its stated time-series role. This is
a bounded scope decision, not an ecosystem-wide claim that Java has no
panel-data software.~\citep{exec25_paneljava_imetrica}

\subsection{An exact staggered-policy counterexample}
Let $D_{gt}$ indicate an absorbing branch policy for group $g$ and period $t$.
Consider two equally weighted groups over three periods, with treatment rows
$(0,1,1)$ and $(0,0,1)$. Stipulate untreated outcomes of zero, and positive
treatment effects producing observed outcomes $(0,1,10)$ and $(0,0,1)$.
These are the same constructed outcomes used in
Section~\ref{exec23:panel-projection}, now carried into the Java comparison.
This is a matched target example, not independent empirical corroboration.
The equal-weight average across the three treated cells is $(1+10+1)/3=4$.
\citep{exec25_paneljava_reference}

Now fit a constant treatment coefficient with group and period effects. In
this balanced example, subtracting row and column means and adding the grand
mean gives the residualized treatment
\[
\widetilde D=
\begin{pmatrix}-1/6&1/3&-1/6\\1/6&-1/3&1/6\end{pmatrix},
\qquad \sum_{g,t}\widetilde D_{gt}^2=1/3.
\]
The treatment coefficient is $\sum\widetilde D_{gt}Y_{gt}/
\sum\widetilde D_{gt}^2$. The three treated-cell contributions have weights
$1,-1/2,1/2$, so the coefficient equals
$1-10/2+1/2=-7/2$. Every stipulated treatment effect is positive, but the
constant-coefficient regression is negative. The numerical regression can
solve the supplied least-squares problem exactly while answering a different
question from the average treated-cell effect. No defect in QR or SVD is
required for this discrepancy.~\citep{exec25_paneljava_reference}

At the last period, both groups are treated. There is no untreated comparison
group in these observed cells. Our construction knows the untreated outcomes
because it stipulated them; an analyst of ordinary observational data would
need assumptions or additional controls to know them. The public test should
therefore check the declared target, supported comparisons and unsupported
horizons, rather than insist that every package returns the oracle value 4.
It should retain cohort and event-time heterogeneity, anticipation windows,
treatment reversals and the control rule as explicit input contracts.
A routine designed for absorbing treatment must refuse or separately model
a reversed policy path; an ordinary design matrix does not decide that rule.
\citep{exec25_paneljava_reference}

\subsection{Sampling dependence and unidentified trends}
Suppose four observed rows consist of two copies from each of two independent
groups, with each group's shock independently equal to $-1$ or $1$ with
probability one half. The mean is $(Z_1+Z_2)/2$ and has variance $1/2$.
Four independent row shocks instead give variance $1/4$. Enumerating the
four group-shock assignments and sixteen independent-row assignments verifies
the distinction exactly. The example concerns two different known sampling
laws; it is not an empirical coverage estimate or an endorsement of a
two-cluster asymptotic confidence interval. A worker must preserve cluster
identity and the justified sampling law, and disclose when too few independent
units prevent its proposed uncertainty procedure.
\citep{exec25_paneljava_reference}

Under an explicitly justified independent-cluster asymptotic regime, a proposed
uncorrected sandwich would use $B[\sum_g X_g^\mathsf{T}\widehat u_g
\widehat u_g^\mathsf{T}X_g]B$, where $X_g$ and $\widehat u_g$ collect
group $g$'s rows. This is a different object from $\widehat\sigma^2B$.
Implementing that expression would still leave finite-cluster correction,
degrees of freedom, resampling, high-dimensional effects and the scientific
sampling assumptions to justify. Neither inspected OLS path establishes
those panel guarantees merely by returning coefficient tests.
\citep{exec25_paneljava_reference,exec25_paneljava_smile,exec25_paneljava_commons}

There is a separate identification limit that better covariance code cannot
repair. Suppose a treated branch's outcome moves from 10 to 15 while a control
branch moves from 10 to 12. The observed difference in differences is 3.
One compatible world has untreated treated-branch outcome 12 and effect 3;
another has untreated outcome 15 and effect 0. Both produce the same four
observed cells. A parallel-trends restriction or additional evidence is needed
to distinguish them. More accurate optimization, an agent-generated narrative
or a small standard error cannot supply the missing counterfactual information.
\citep{exec25_paneljava_reference}

\subsection{Reuse, rights and the public prototype route}
For a Java service, our proposed architecture first separates an estimand
contract from numerical execution. A narrow Java numerical delegate could
serve a justified regression subproblem, while a separately reviewed R or
Python panel worker retains its own group-time, control and uncertainty
semantics. The earlier differences and csdid inspections provide concrete
comparison routes, with their own unresolved normalization, dependency and
failure concerns. Delegation must preserve row IDs, panel keys, treatment
history, weights, missingness, cohort/time support and the exact returned
uncertainty type. A cross-language call is not evidence of statistical parity.
\citep{exec24_panelpy_differences,exec24_panelpy_csdid,exec25_paneljava_reference}

The selected Smile build declares Java 25 with preview features, and its
root notice distinguishes core-library GPLv3/commercial terms from separate
server and studio terms. These are inspected declarations, not a legal
clearance or a quoted commercial price. Apache's selected tree declares
Apache-2.0 material and third-party notices, a Java 8 compilation target
and multiple component dependencies. iMetrica includes root GPLv3 text and
native library requirements. A final choice needs the exact distribution,
transitive dependencies, intended deployment, rights review, maintenance
commitments and measured full costs. We do not transfer one module's terms
to another or infer a supported runtime from a build file alone.
\citep{exec25_paneljava_smile,exec25_paneljava_commons,exec25_paneljava_imetrica}

The immediate proposal decision is to retain Smile and Commons Math as
numerical references, and to defer direct panel admission and whole-library
porting. The public qualification route is concrete: freeze a source and
runtime profile; compare coefficients, covariance scaling, intercept and
rank-failure behavior on known linear problems; then test the declared causal
panel target on staggered, reversed, anticipated, unbalanced and dependent
fixtures with known intervention answers or required refusals. Test data
conversion and every error path, including missing groups and unsupported
post-treatment horizons. Keep numerical parity, scientific correctness and
operating acceptance as separate reported outcomes.
\citep{exec25_paneljava_reference}

Complete costs also need distinct denominators. An illustrative USD 2,400
experiment costs USD 24 per submitted case for 100 submissions, or USD 40
per scientifically usable case if only 60 qualify. Neither ratio is a
measured price or an ROI result. Report engineering, expert adjudication,
dependency and entitlement work, failed attempts and runtime in the numerator.
The bounded repository search and selected source paths leave further Java
and JVM panel, synthetic-control and bridge candidates to inspect. This
comparison narrows the implementation route while preserving the remaining
coverage and causal-identification obligations.~\citep{exec25_paneljava_reference}


% SOURCE docs/release/2026-10-29/source/docs/writing-repair/2026-10-29/src/585ea2e83bcea072.tex CHAR 0

\section{C++ panel delegates: preserving the statistical wrapper}
\label{sec:exec25-panels-cxx}
A bank moving a panel estimator into a C++ service must preserve more than
its coefficient solver. The R or Python wrapper may select observations,
encode time, remove singleton groups, define comparison cohorts, scale
covariances and decide which quantities are unavailable. We inspect selected
fixest C++ kernels together with their already reviewed R callers, and a
dated xhdfe C++ source profile with public interfaces and numerical test
definitions. Four further candidates supply explicit survey leads. This
comparison concerns retained source text and exact project constructions;
no downloaded implementation or upstream test was run.
\citep{exec25_panelcxx_fixest,exec20_fixest_source,exec25_panelcxx_xhdfe,exec25_panelcxx_reference}

\subsection{The residualization boundary}
The inspected fixest demeaning entry point accepts R objects, weights,
fixed-effect IDs and tables, varying-slope information, iteration limits,
tolerances and thread settings. It returns transformed variables and iteration
counts, with optional means and fixed-effect coefficients. Its headers and
generated exports use Rcpp and R objects. Copying the kernel therefore does
not produce an R-independent panel product automatically. The preceding
R source review remains necessary for sample selection, weights, covariance
options and treatment-time semantics. A typed R worker is the narrower
initial integration route; a native port must explicitly replace those
responsibilities.~\citep{exec25_panelcxx_fixest,exec20_fixest_source}

Weighted residualization illustrates the contract. For observations $(1,3)$
with weights $(1,3)$, the weighted mean is $5/2$. Subtracting it gives
residuals $(-3/2,1/2)$ with weighted sum zero. Subtracting the unweighted
mean instead gives $(-1,1)$, whose weighted sum is 2. A port that silently
changes the projection changes the supplied estimating equations. In a
bank's branch panel, the distinction can arise from exposure or portfolio
weights; the interpretation of those weights must be fixed before execution.
An exact arithmetic witness supplies a regression test, not a justification
for the statistical weight law.~\citep{exec25_panelcxx_reference}

The inspected Cholesky kernel tests successive diagonal remainders against
a tolerance, returns excluded-column indicators and a minimum-norm diagnostic,
and can return an all-removed status. Its R caller preserves unavailable
coefficients and handles the all-removed case. Consider treatment constant
within each of two branches, $(0,0,1,1)$ across their two observations.
Removing branch effects leaves a zero treatment column. The treatment effect
is not identified by that within-branch variation. A native adapter must
preserve this unavailable target rather than report a causal effect of zero.
Column ordering, tolerance and scale also belong to numerical comparison
fixtures.~\citep{exec25_panelcxx_fixest,exec20_fixest_source,exec25_panelcxx_reference}

\subsection{Time ordering is supplied by the caller}
The selected panel Newey--West kernel's comment says to sort by unit then
time, but its implementation forms contiguous time blocks and matches unit
IDs across those blocks. The inspected R caller sorts by time then unit,
reorders the score matrix consistently, rejects duplicate unit--time pairs,
and sets the zero-lag weight to one half because the kernel later adds the
transpose. The operative contract is established by the caller and body
together; a port based only on the comment could change it. This inspection
does not establish an observed failure in the original R path.
\citep{exec25_panelcxx_fixest,exec20_fixest_source}

For scalar scores $(1,3)$ for branch A and $(2,4)$ for branch B over two
periods, the sum of squares is 30 and the matched same-branch lag product
is 11. A lag-one Bartlett construction with off-zero weight $1/2$ and the
kernel's final symmetrization yields $30+11=41$. Passing zero-lag weight
one instead of one half would yield $60+11=71$ under the same algebra.
This is an authored witness about interface normalization, not a package
measurement. In fact, the selected kernel caps its weight count at $T-1$
after dropping a trailing zero. With two periods this leaves only lag zero,
so the lag-one algebra above must not be reported as that native return.
Appending a third, zero-score period retains the same sum of squares and
lag-one product while allowing the lag-one loop under the inspected cap.
Both short-panel and zero-padded cases belong in the public test inventory.
\citep{exec25_panelcxx_fixest,exec25_panelcxx_reference} The proposed tests also distinguish event time from calendar
time: mapping two observed months to consecutive integers can erase a
missing month. Whether such compression is appropriate depends on the
declared sampling and lag interpretation, not on storage convenience.
\citep{exec25_panelcxx_reference}

\subsection{A richer standalone profile still needs a causal target}
The inspected xhdfe tree declares a 2.29.3 development candidate, with CMake
version 2.29.3.20261007 and C++17. Its public configuration distinguishes
homoskedastic, robust and clustered covariance; several absorption and
tolerance modes; singleton removal; frequency-weight bookkeeping;
fixed-effect and cluster degrees-of-freedom adjustments; and intercept
handling. The result structure includes covariance, counts before and after
selection, cluster counts, convergence and a covariance-repair flag. These
are useful observability fields for a proposed delegate. They do not
establish a cohort-time causal estimand or a universal uncertainty guarantee.
\citep{exec25_panelcxx_xhdfe}

The inspected implementation follows cluster scaling, recovered-intercept
covariance and multiway covariance correction paths. It distinguishes
unavailable inference for fewer than two clusters and other unsupported
degrees-of-freedom cases. A postcondition checks finite outputs and shapes
for the quantities whose inference is available. A separate normal-equation
check can use the supplied Gram system; per-row certification is an
additional path, not an automatic consequence of every default call.
Upstream test definitions include empty-slope cluster cases, missing cluster
labels and full covariance comparisons. Their inspected existence motivates
public tests here; it is not a local reproduction of their results.
\citep{exec25_panelcxx_xhdfe}

The earlier constructed staggered-policy example applies equally to a more
capable numerical engine. Positive stipulated effects averaged over treated
cells can coexist with a negative constant two-way-effects coefficient.
Faster absorption, more precise arithmetic and cluster covariance do not
change the meaning of that regression target. The request must specify
whether it seeks that coefficient, a cohort-time effect, an event-time
aggregation, a synthetic-control contrast or a nonlinear structural
parameter. Unsupported target conversions require refusal or a separately
justified estimator.~\citep{exec25_paneljava_reference,exec25_panelcxx_reference}

\subsection{Checking the complete covariance object}
The selected xhdfe multiway path may project a standardized full covariance
matrix onto the positive-semidefinite cone and project the slope block
separately, then replace that block in the full matrix. Its ordinary-precision
helper returns false both when no correction is needed and when its
eigensolver reports failure; a wider-precision path instead throws on
eigensolver failure. The inspected output postcondition checks finiteness
and dimensions rather than a full-matrix eigenvalue condition. These
specific paths justify adversarial test requirements. They do not prove
that a bad covariance occurred in an upstream experiment.
\citep{exec25_panelcxx_xhdfe}

An exact construction shows why separately valid projections do not suffice
as a general full-matrix guarantee. Start with
\[
V=\begin{pmatrix}1&2\\2&1\end{pmatrix}.
\]
Its eigenvalues are 3 and $-1$. Clamping the negative one gives
$P(V)=\tfrac32\begin{pmatrix}1&1\\1&1\end{pmatrix}$. The original
one-dimensional first block is already positive, so its separate projection
remains 1. Replacing the first block in the full projection yields
\[
\overline V=\begin{pmatrix}1&3/2\\3/2&3/2\end{pmatrix},
\quad\det(\overline V)=-3/4,
\quad (3,-2)\overline V(3,-2)^\mathsf{T}=-3.
\]
This is a general algebraic counterexample to a block-replacement guarantee.
It is not proof that this matrix is attainable from xhdfe's data and
normalization paths, and it is not an observed native package failure.
A public fixture must establish reachability before making that stronger
claim.~\citep{exec25_panelcxx_reference}

For a bank requesting a portfolio contrast involving both slope and
intercept, negative quadratic variance is unusable even if every diagonal
entry is finite and nonnegative. The proposed result contract therefore
retains raw and corrected covariance, the normalization metric and repair
status, checks supported contrasts and full-matrix validity, and refuses
ambiguous or failed numerical validation. A correction that ensures a
nonnegative matrix still does not demonstrate frequentist coverage. The
sampling law, small-cluster restrictions and decision tolerance remain
scientific obligations.~\citep{exec25_panelcxx_reference}

\subsection{Concrete alternatives and bounded decisions}
The survey now records four additional dated leads separately. The elakkis
SDID repository declares an R package with Rcpp/RcppArmadillo and a
synthetic-difference-in-differences workflow; its kernel and inference
implementation remain unreviewed here. capybara 2.4.0 declares fixed-effect
linear and generalized linear models with clustered standard errors and
cpp4r/armadillo4r dependencies. Its README distinguishes an earlier published
implementation from a planned account of the 2.x changes. We do not cite
or inherit an underlying paper's technical claims from this README alone.
\citep{exec25_panelcxx_sdid,exec25_panelcxx_capybara}

The gfe descriptor declares grouped fixed-effects models, balanced and
unbalanced panels and group-heterogeneous coefficients, but its License field
contains an unfilled template while its root LICENSE contains MIT text.
That discrepancy is retained for resolution. NNRPanel declares nonlinear
interactive fixed-effects estimation, tuning and bias-correction routines,
with Rcpp/RcppEigen dependencies and a GPL-3 declaration; its mathematical
and implementation claims require the underlying originals and source-path
review. These four are scoped candidates, not approved substitutes for the
current panel profile.~\citep{exec25_panelcxx_gfe,exec25_panelcxx_nnr}

For fixest, the preferred initial route remains the original package behind
a typed worker, preserving its inspected R semantics. For xhdfe, retain
the C++ interface as a candidate for a narrowly qualified public profile,
while deferring adoption and porting pending the complete absorber,
covariance, backend and binding review. Its root MIT text does not cover
every dependency by implication: the acquired Eigen notice records MPL2
and additional third-party terms. Exact installed artifacts, compiler
settings, optional CUDA, licensing and maintenance all need separate
qualification. README runtime tables are upstream claims and are not
promoted to this proposal's measured speed or cost advantage.
\citep{exec20_fixest_source,exec25_panelcxx_xhdfe,exec25_panelcxx_reference}

The public test program should compare weighted residualization with
explicit-dummy solutions, verify rank and unsupported treatment targets,
permute rows while preserving panel IDs, vary gaps and missing groups,
exercise cluster and covariance-repair boundaries, and test every declared
failure and conversion path. It should report requested and resolved
backend and mode, all attempts and complete costs. An illustrative USD
3,600 experiment is USD 30 per submission for 120 cases, or USD 45 per
usable result for 80. These are budget identities, not measured prices
or a claim of economic benefit. Remaining source coverage and native
qualification stay explicit while independent proposal work continues.
\citep{exec25_panelcxx_reference}

The assembled comparisons return to the branch rollout with a more precise
choice. Cohort methods preserve treatment timing and supported populations;
synthetic comparisons require defended donor and calendar information;
feedback analysis requires a specified intervention and stable structural
meaning. A numerical engine can support any of these only through a matching
statistical design. The detailed failures therefore matter alongside the
successful algebra: lost rows, changed weights, substituted intervals and
unreported solver limits can alter the decision. The next chapter develops
trend-sensitivity inference, asking how much confidence the rollout conclusion
warrants when counterfactual trends remain uncertain.


